%%%%%%%% ICML 2026 EXAMPLE LATEX SUBMISSION FILE %%%%%%%%%%%%%%%%%

\documentclass{article}

% Recommended, but optional, packages for figures and better typesetting:
\usepackage{microtype}
\usepackage{graphicx}
\usepackage{subcaption}
\usepackage{booktabs} % for professional tables

% hyperref makes hyperlinks in the resulting PDF.
% If your build breaks (sometimes temporarily if a hyperlink spans a page)
% please comment out the following usepackage line and replace
% \usepackage{icml2026} with \usepackage[nohyperref]{icml2026} above.


% Attempt to make hyperref and algorithmic work together better:
\newcommand{\theHalgorithm}{\arabic{algorithm}}

% Use the following line for the initial blind version submitted for review:
% \usepackage{icml2026}

% For preprint, use
% \usepackage[preprint]{icml2026}

% If accepted, instead use the following line for the camera-ready submission:
\usepackage[accepted]{icml2026}

\usepackage{amsmath}
\usepackage{amssymb}
\usepackage{mathtools}
\usepackage{amsthm}

\usepackage[utf8]{inputenc} % allow utf-8 input
\usepackage[T1]{fontenc}    % use 8-bit T1 fonts
\usepackage{url}            % simple URL typesetting
\usepackage{booktabs}       % professional-quality tables
\usepackage{amsfonts}       % blackboard math symbols
\usepackage{nicefrac}       % compact symbols for 1/2, etc.
\usepackage{microtype}      % microtypography
\usepackage{xcolor}         % colors

\usepackage{amsmath}
\usepackage{amssymb}
\usepackage{mathtools}
\usepackage{amsthm}

\input{math_commands.tex}
\usepackage{amssymb}
\usepackage{amsthm}
\usepackage{mathtools}
\usepackage{xspace}
\usepackage{enumitem}
\usepackage{wrapfig}
% \usepackage{minted} 


\newcommand\restr[2]{{% we make the whole thing an ordinary symbol
  \left.\kern-\nulldelimiterspace % automatically resize the bar with \right
  #1 % the function
  \littletaller % pretend it's a little taller at normal size
  \right|_{#2} % this is the delimiter
  }}

\usepackage{xurl}
% \usepackage{everyshi}
\usepackage[pagebackref=true,hyperindex,breaklinks]{hyperref}

\hypersetup{
    colorlinks=true,
    linkcolor=thulianpink,
    filecolor=pear,      
    urlcolor=tiffanyblue,
    citecolor=amethyst,
}

\renewcommand*{\backref}[1]{}
\renewcommand*{\backrefalt}[4]{%
    \ifcase #1%
          \or Cited on page~#2.%
    \else Cited on pages~#2.%
\fi%
}
\usepackage{wrapfig}
\usepackage{tablefootnote}
\usepackage{listings}

\usepackage{thm-restate}

%dashed line
\usepackage{array}
\usepackage{arydshln}
\setlength\dashlinedash{0.2pt}
\setlength\dashlinegap{1.5pt}
\setlength\arrayrulewidth{0.3pt}

\usepackage[capitalize,noabbrev]{cleveref} 
\newtheorem{theorem}{Theorem}
\newtheorem{proposition}[theorem]{Proposition}  % Shares the counter with theorem
\newtheorem{definition}[theorem]{Definition}    % Shares the counter with theorem
\newtheorem{lemma}[theorem]{Lemma}              % Shares the counter with theorem
\newtheorem{corollary}[theorem]{Corollary}      % Shares the counter with theorem
\newtheorem{assumption}{Assumption}    % Shares the counter with theorem
\newtheorem{remark}[theorem]{Remark}
\newtheorem{goal}[theorem]{goal}
\newtheorem{example}[theorem]{Example}
\newtheorem{problem}[theorem]{Problem}
\newcommand{\thickoverline}[1]{%
  \overline{\mathrel{#1}\kern-1.5pt\rule{0pt}{0.5ex}}%
}


\crefname{dataset}{Dataset}{Datasets}
\crefname{method}{Method Step}{Method Steps}
\crefname{problem}{Problem}{Problems}
\crefname{definition}{Defn.}{Defns.}
\crefname{theorem}{Thm.}{Thms.}
\crefname{proposition}{Prop.}{Props.}
\crefname{assumption}{Asm.}{Asm.}
\crefname{remark}{Remark}{Remarks}
\crefname{example}{Eg.}{Egs.}
\crefname{equation}{Eqn.}{Eqns.}
\crefname{appendix}{Apx.}{Apx.}

\usepackage{epigraph} 

% https://latexcolor.com/
\definecolor{teal}{rgb}{0.0, 0.5, 0.5}
\definecolor{amethyst}{rgb}{0.6, 0.4, 0.8}
\definecolor{thulianpink}{rgb}{0.87, 0.44, 0.63}
\definecolor{tiffanyblue}{rgb}{0.04, 0.73, 0.71}
\definecolor{pear}{rgb}{0.82, 0.89, 0.19}
\definecolor{applegreen}{rgb}{0.55, 0.71, 0.0}
\definecolor{burgundy}{rgb}{0.5, 0.0, 0.13}
\definecolor{persianindigo}{rgb}{0.2, 0.07, 0.48}


\usepackage{subcaption}
\captionsetup[subtable]{skip=7pt}
\captionsetup[sub]{font=normal,labelfont=bf} % subcaptions small, bold labels
\captionsetup[table]{skip=7pt}
% \captionsetup[figure]{font=small}
% \captionsetup[table]{font=small}

% \usepackage{sidecap}

\usepackage[most]{tcolorbox}
\tcbuselibrary{listingsutf8}
% \tcbset{
%   colback=brown!5!white,   % Background color of the main box
%   colframe=black,         % Border color of the box
%   fonttitle=\bfseries,    % Bold font for the title
%   boxrule=0.4mm,          % Border thickness          % Sharp corners for the box
%   coltitle=white,         % Title text color (set to white)
%   colbacktitle=purple!80!white,     % Background color of the title bar
% }

\lstset{
  language=Python,
  basicstyle=\ttfamily\small,
  keywordstyle=\color{blue},
  commentstyle=\color{gray},
  stringstyle=\color{red},
  backgroundcolor=\color{brown!15!white},
  frame=none,
  showstringspaces=false,
  breaklines=true,
  numbers=left,
  numberstyle=\tiny\color{gray},
  numbersep=5pt
}


\tcbset{
base/.style={
  arc=0mm,
  boxrule=0.3mm, % Width of the border around the box
  left=0.1mm, % Left margin within the box
  right=0.5mm, % Right margin within the box
  top=0.5mm, % Top margin within the box
  bottom=0.5mm, % Bottom margin within the box
  }
}

% \newtcolorbox{mainbox}[1][]{
%   colback=pear!5!white, % Background color of the box
%   colframe=pear!75!black, % Frame color
%   base={#1}
% }

\newtcolorbox{subbox}[1][]{
  colback=black!5!white, % Background color of the box
  colframe=black!75!black, % Frame color
  base={#1}
}

% \setlength{\leftmargini}{0pt} % Remove indentation for the first level of itemize



\usepackage{tabularx}

\DeclareRobustCommand{\parhead}[1]{\textbf{#1}~}


\newcommand{\setmylabel}[1]{
  \renewcommand{\mylabel}{#1:}
}
\newcommand{\bb}{\mathbf{b}}
\newcommand{\x}{\mathbf{x}}
\newcommand{\z}{\mathbf{z}}
\newcommand{\A}{\mathbf{A}}
\newcommand{\D}{\mathbf{D}}
%\newcommand{\P}{\mathbf{P}}
\renewcommand{\c}{\mathbf{c}}
\newcommand{\q}{\mathbf{q}}
\newcommand{\br}{\mathbf{r}}
\newcommand{\s}{\mathbf{s}}
\newcommand{\zerob}{\mathbf{0}}
\newcommand{\C}{\mathcal{C}}
\newcommand{\Z}{\mathcal{Z}}
\newcommand{\X}{\mathcal{X}}
\newcommand{\deltaBold}{\boldsymbol{\delta}}
\newcommand{\deltaz}{\deltaBold^z}
\newcommand{\hatdeltaz}{\hat{\deltaBold}^z}
\newcommand{\tildedeltaz}{\tilde{\deltaBold}^z}
\newcommand{\deltac}{\deltaBold^c}
\newcommand{\hatdeltac}{\hat{\deltaBold}^c}
\newcommand{\tildedeltac}{\tilde{\deltaBold}^c}
\newcommand{\thicktilde}[1]{\mathbf{\tilde{\text{$#1$}}}}
\newcommand{\bsigma}{\boldsymbol{\Sigma}}
\newcommand{\transp}{^\top}
\newcommand{\bblambda}{\boldsymbol{\Lambda}}
\newcommand{\blambda}{\boldsymbol{\lambda}}

\newcommand{\ad}[1]{\textcolor{olive}{\textbf{AD:} #1}}
\newcommand{\sj}[1]{\textcolor{violet}{\textbf{sj:} #1}}
\newcommand{\ds}[1]{\textcolor{teal}{\textbf{DS:} #1}}
\newcommand{\seb}[1]{\textcolor{cyan}{[\textbf{SL:} #1]}}
\newcommand{\red}[1]{\textcolor{red}{#1}}
\newcommand{\isae}{\texttt{SSAE}\xspace}
\newcommand{\aff}{\texttt{aff}\xspace}
\newcommand{\md}{\texttt{MD}\xspace}

% Todonotes is useful during development; simply uncomment the next line
%    and comment out the line below the next line to turn off comments
%\usepackage[disable,textsize=tiny]{todonotes}
\usepackage[textsize=tiny]{todonotes}

% \usepackage{minitoc}
% \setcounter{parttocdepth}{3}
% \setcounter{secnumdepth}{3}
% % Make the "Part I" text invisible
% \renewcommand \thepart{}
% \renewcommand \partname{}


% The \icmltitle you define below is probably too long as a header.
% Therefore, a short form for the running title is supplied here:



% if you use cleveref..
\usepackage[capitalize,noabbrev]{cleveref}


% Todonotes is useful during development; simply uncomment the next line
%    and comment out the line below the next line to turn off comments
%\usepackage[disable,textsize=tiny]{todonotes}
\usepackage[textsize=tiny]{todonotes}

% The \icmltitle you define below is probably too long as a header.
% Therefore, a short form for the running title is supplied here:
\icmltitlerunning{Sparse Shift Autoencoders}

\begin{document}
\ClearShipoutPicture% remove line-number ruler added by icml2026.sty

\twocolumn[
  \icmltitle{Sparse Shift Autoencoders for Identifying \\Concepts from Large Language Model Activations}
% Identifying Concepts from Large Language Model Activations with Sparse Shift Autoencoders
  % It is OKAY to include author information, even for blind submissions: the
  % style file will automatically remove it for you unless you've provided
  % the [accepted] option to the icml2026 package.

  % List of affiliations: The first argument should be a (short) identifier you
  % will use later to specify author affiliations Academic affiliations
  % should list Department, University, City, Region, Country Industry
  % affiliations should list Company, City, Region, Country

  % You can specify symbols, otherwise they are numbered in order. Ideally, you
  % should not use this facility. Affiliations will be numbered in order of
  % appearance and this is the preferred way.
  \icmlsetsymbol{equal}{*}

  \begin{icmlauthorlist}
    \icmlauthor{Shruti Joshi}{yyy}
    \icmlauthor{Andrea Dittadi}{comp,rrr}
    \icmlauthor{S\'ebastien Lachapelle}{zzz}
    \icmlauthor{Dhanya Sridhar}{yyy}
    %\icmlauthor{}{sch}
    %\icmlauthor{}{sch}
    %\icmlauthor{}{sch}
  \end{icmlauthorlist}

  \icmlaffiliation{yyy}{Mila-Qu\'ebec AI Institute \& Universit\'e de Montr\'eal}
  \icmlaffiliation{comp}{Helmholtz AI and Technical University of M\"unich}
  \icmlaffiliation{rrr}{MPI for Intelligent Systems, T\"ubingen}
\icmlaffiliation{zzz}{Samsung - SAIT AI Lab, Montr\'eal}
  \icmlcorrespondingauthor{Shruti Joshi}{shrutijoshi98@gmail.com}

  % You may provide any keywords that you find helpful for describing your
  % paper; these are used to populate the "keywords" metadata in the PDF but
  % will not be shown in the document
  \icmlkeywords{Machine Learning, ICML}

  \vskip 0.3in
]

% this must go after the closing bracket ] following \twocolumn[ ...

% This command actually creates the footnote in the first column listing the
% affiliations and the copyright notice. The command takes one argument, which
% is text to display at the start of the footnote. The \icmlEqualContribution
% command is standard text for equal contribution. Remove it (just {}) if you
% do not need this facility.

% Use ONE of the following lines. DO NOT remove the command.
% If you have no special notice, KEEP empty braces:
\printAffiliationsAndNotice{}  % no special notice (required even if empty)
% Or, if applicable, use the standard equal contribution text:
% \printAffiliationsAndNotice{\icmlEqualContribution}

\begin{abstract}
  Unsupervised approaches to large language model (LLM) interpretability, such as sparse autoencoders (SAEs), offer a way to decode LLM activations into interpretable and, ideally, controllable concepts. On the one hand, these approaches alleviate the need for supervision from concept labels, paired prompts, or explicit causal models. On the other hand, without additional assumptions, SAEs are not guaranteed to be identifiable. In practice, they may learn latent dimensions that entangle multiple underlying concepts. If we use these dimensions to extract vectors for steering specific LLM behaviours, this non-identifiability might result in interventions that inadvertently affect unrelated properties. In this paper, we bring the question of identifiability to the forefront of LLM interpretability research. Specifically, we introduce Sparse Shift Autoencoders ({\isae}s) which learn sparse representations of \textit{differences} between embeddings rather than the embeddings themselves. Crucially, we show that {\isae}s are identifiable from paired observations which differ in multiple unknown concepts, but not all. With this key identifiability result, we show that we can steer single concepts with only this weak form of supervision. Finally, we empirically demonstrate identifiable concept recovery across multiple real-world language datasets by disentangling activations from different LLMs.
    
\end{abstract}

\section{Introduction}
 As increasingly powerful large language models (LLMs) are deployed and widely used, the need to interpret and steer their behavior grows. For both interpretability and steering, we require techniques to disentangle LLM activations into semantically meaningful, and ideally, manipulable concepts.
A large class of LLM interpretability methods rely on supervision from ground truth concepts \citep{koh2020conceptbottleneck}, paired prompts \citep{turner2024steeringlanguagemodelsactivation}, target LLM completions \citep{subramani2022extracting} and abstract causal models of behavior \citep{geiger2024DAS} to map activations to concepts.
For example, using contrastive pairs of prompts that differ by a single concept, recent papers have found vectors in activation space that encode sycophancy \citep{rimsky2024steering}, truthfulness \citep{park2025steer}, and refusal \citep{arditi2024refusallanguagemodelsmediated}.
However, acquiring such supervision is costly, motivating unsupervised methods for concept learning.

Sparse autoencoders (SAEs) have emerged as a popular approach to unsupervised LLM interpretability \citep{cunningham2023sparseautoencodershighlyinterpretable}. Taking inspiration from sparse dictionary learning, SAEs encode LLM activations in a sparse and overcomplete representation. While we might hope that there is a one-to-one correspondence between the learned dimensions and interpretable concepts, \citep{wu2025axbenchsteeringllmssimple, Menon2024AnalyzingO} show empirical evidence that SAEs significantly underperform supervised methods, suggesting that they may not be \emph{identifiable}: that is, they could learn latent dimensions that entangle interpretable concepts.
Consequently, if we use SAEs to extract activation directions to steer LLM behavior, non-identifiability could result in changes to unrelated properties.

In this paper, we propose Sparse Shift Autoencoders ({\isae}s), models for provably recovering steering vectors without the need for concept labels, contrastive pairs and other supervision signals about concepts. Crucially, {\isae}s learn from sparse multi-concept shifts: paired samples in which multiple unknown concepts differ, but not all of them. Such samples are cheap to obtain, for example, by pairing sentences from Wikipedia articles, or by using LLMs to synthetically generate contrastive texts.  Briefly, an \isae maps embedding differences between samples in a pair to a latent space that reflects the concept changes and uses a linear decoding function to reconstruct the difference vector. This architecture reflects the \textit{linear representation hypothesis} \citep{mikolov-etal-2013-linguistic, jiang2024originslinearrepresentationslarge} in assuming that concepts are linearly encoded by LLMs. Crucially, we regularize the latent representation to be sparse, meaning that each shift is modelled using as few concept changes as necessary.
We then leverage the results developed by \citet{lachapelle2023synergies} and \citet{xu2024sparsityprinciplepartiallyobservable} to prove that the proposed \isae approach identifies \emph{some} concepts, under suitable assumptions on distribution that generated the data. We also show how this allow to extract extract valid steering vectors, i.e. direction in the LLM representation that changes a single concept.
We study the \isae empirically on challenging language datasets and models, finding many settings where they outperform SAEs as well as other related steering methods that require supervision.

In sum, this work: 1) formalizes the problem of recovering interpretable concepts from sparse multi-concept shifts, from the lens of identifiability; 2) proposes the \isae to model these sparse multi-concept shifts and establishes identifiability guarantees for these models based on sparsity regularization; 3) using multiple real-world language datasets and LLMs, empirically verifies the identifiability result and demonstrates the benefits of an identifiable model for accurately predicting target steered embeddings.



\begin{tcolorbox}[
  enhanced jigsaw,
  breakable,
  colback=gray!6,
  colframe=gray!60!black,
  title=\textsc{What is identifiability?},
  boxrule=0.5pt,
  arc=2pt,
  left=4pt, right=4pt, top=2pt, bottom=2pt
]
A representation learning method is \emph{identifiable} if, given sufficient data, it provably recovers the true underlying factors up to benign indeterminacies like permutation and rescaling. Without identifiability, learned features may conflate multiple concepts or capture spurious directions. Moreover, non-identifiable methods can yield different solutions across training runs, complicating reproducibility. Identifiability ensures the learned representation reflects genuine invariant structure learned from the data, thereby enabling meaningful downstream interpretation.
\end{tcolorbox}


\section{Problem formulation}
\label{sec:prob_form}

We observe texts $\x \in \X \subseteq \mathbb{R}^{d_x}$ that are generated from underlying \emph{concept representations} $\c \in \C  \subseteq \mathbb{R}^{d_c}$ through an unknown generative process $g: \C \to \X$ so that $\x=g(\c)$.
While we cannot observe the concept representation $\c$ of an observation $\x$, we have access to \emph{learned representations} $\z=f(\x)$, where the function $f: \X \rightarrow \Z \subseteq \mathbb{R}^{d_z}$ maps observations $\x$ to $d_z$-dimensional real vectors $\z \in \Z$, known as their \emph{embeddings}.
Throughout this paper, we consider the case where $f(\x)$ comes from an autoregressive language model and is the embedding of the final token $\x_T$ in the residual stream after the final layer. 

We assume that the concepts $\c$ are \textit{encoded} in the representations $\z$ through the unknown composite function $\z = f(g(\c))$. We consider concept perturbations,
\begin{align}
    \tilde \c \coloneqq \c + \deltac; \quad \deltac_k = \lambda \mathbf{e}_k,
\end{align}
where $\deltac$ is called the \textbf{concept shift} vector, $\lambda$ is the magnitude of the perturbation, and $\deltac_k \neq 0$ for all perturbed concepts $k$.

\textbf{Main goal.} We want to map unlabelled concept shifts $\deltac$ to their corresponding vectors in the space of LLM activations. (Refer to \cref{apx:steering} for a formal treatment of steering.). 

\begin{figure}
    \centering
    \includegraphics[width=\linewidth]{Group_17.png}
    \caption{{\isae}s map multi-concept shifts in embedding space to concept shifts, leveraging the latter's sparsity, thereby recovering steering vectors. The learnt steering vectors are identified up to permutation and scaling.}
    \label{fig:main}
\end{figure}

The key challenge is that we only observe texts ($\x$) and their embeddings ($\z$), and thus, we cannot directly learn a mapping from concepts shifts $\deltac$ to LLM activation shifts.
A naive unsupervised approach is to fit an autoencoder to LLM embeddings $\z$ so that for any input, we can encode it in a latent space, implement the desired concept shift $\deltac$ in that space, and decode it to obtain a perturbed embedding $\tilde \z$. However, unless the autoencoder is guaranteed to encode embeddings $\z$ in a latent space that captures concepts, this naive approach will result in perturbations $\tilde \z$ that \emph{do not reflect the desired concept shifts}. Unfortunately, unconstrained autoencoding objectives are non-identifiable \citep{HYVARINEN1999429}, and sparse autoencoding objectives \citep{cunningham2023sparseautoencodershighlyinterpretable} may not be able to invert embeddings to potentially billions of concepts. As such, there is no guarantee that such approaches recover latent concepts from observed embeddings $\z$, posing a risk for steering. 

\textbf{Key idea.} We develop an identifiable autoencoding method called sparse shift encoders ({\isae}s). The key idea behind {\isae}s are multi-concept shift data, illustrated in \cref{fig:main}. As an example, consider two text snippets $\x:$ \texttt{pretty pink round matte frame} and $\tilde \x: $ \texttt{ugly grey round metallic frame}. Both $\x$ and $\tilde \x$ seem to be encoded by $5$ concepts each-- a descriptive adjective (pretty/ugly), colour (pink/grey), shape (round), texture (matte/metallic), and object (frame). However, when we consider what has changed from $\x$ to $\tilde \x$, it's a smaller set of ($3$) concepts, and it is also possible to imagine pairs which vary by just a single concept. An \isae provably recovers these inter-sample concept shifts by regularizing the inferred concept shifts $\hatdeltac$ to be sparse. 

\section{Sparse Shift Autoencoders ({\isae}s)}
\label{sec:dgp} 
We start by describing the data-generating process and the set of concepts learnable via inter-sample differences, before proposing a method for learning steering vectors for these concepts. Following \citet{locatello2020weaklysupervised}, we consider paired observations $(\x, \tilde\x)$ assumed to be sampled from the following generative process:
\begin{align}
    &S \sim p(S), \quad (\c, \tilde\c) \sim p(\c, \tilde\c \mid S), \label{eqn:dgp}\\
    &\x \coloneqq g(\c), \quad \tilde\x \coloneqq g(\tilde\c) \,,
\label{eq:dgp}
\end{align}
where $S \subseteq \{1, \dots, d_c\}$ denotes the subset of concepts that vary between $\x$ and $\tilde\x$.
More precisely, $p(\c, \tilde\c \mid S)$ is such that, with probability one, $\c_k = \tilde\c_k$ for all $k \not\in S$. Crucially, across each pair of observations, an unknown set of concepts changes.
% , allowing us to leverage paired data that is cheap to obtain, such as posts and their replies on social media sites. 
For what follows, it will be useful to define $V \subseteq \{1, \dots, d_c\}$ to be the set of \textit{varying concepts}:
\begin{align}
V \;\coloneqq\; \bigcup\nolimits_{\smash{S:\,p(S)>0}} S \,. \label{defn:varying_concepts}
\end{align}
The set $V$ thus contains the concepts that can change in a pair $(\x, \tilde\x)$. Even though concepts outside $V$ are assumed to remain fixed \textit{within} a pair $(\x, \tilde\x)$, they can still vary \emph{across} pairs. Without loss of generality, assume that ${V \coloneqq \{1, .\dots, |V|\}}$. 

% Next, we require a function that encodes the variations between $(\z, \tilde \z)$ into the variations $\deltac$ between $(\c, \tilde \c)$. For this, we define the 

Next, we consider difference vectors $\deltaz \coloneqq f(\tilde\x) - f(\x) = \tilde\z - \z$. These vectors capture how underlying concept differences between a pair of inputs $\x$ and $\tilde\x$ are represented in the space of LLM embeddings.  
An important assumption made by \citep{rajendran2024learning,park2023linear} helps us relate these difference vectors to concept shifts:
% and state an important assumption to explain why it would be useful to consider them:

\begin{assumption}[Linear representation hypothesis (LRH)]
\label{ass:lrh}
The generative process $g: \C \to \X$ and the learned encoding function $f: \X \to \Z$ are such that  $f \circ g : \C \to \Z$ is linear, implying there exists a $d_z \times d_c$ real matrix $\A$ such that:
    \begin{align}
        \label{eq:linear_entangle}
        \z = f(g(\c)) = \A \c \ .
        %\deltaz = f(g(\deltac)) = \A \deltac,
    \end{align}
\end{assumption}
Put simply, the LRH says that the learned representation $\z$ \emph{linearly encodes concepts}. Consequently, difference vectors $\deltaz$ are also linearly related to concept shifts so that $\deltaz = \A \deltac$.

A long line of work provides evidence for this hypothesis (c.f. \citet{rumelhart1973model, hinton1986learning, mikolov-etal-2013-linguistic, ravfogel2020null}). More recently, theoretical work justifies why linear properties could arise in these models (c.f. \citet{ jiang2024originslinearrepresentationslarge, roeder2021linear, marconato2024all}). \Cref{sec:rel} provides a full list of related work, while \cref{apx:lrh} provides an explanation of the equivalence between LRH's different interpretations. 

Sparse Shift Autoencoders ({\isae}s) take as input the observed difference vectors $\deltaz_V$ and model them with an affine encoder $r:\mathbb{R}^{d_z} \rightarrow \mathbb{R}^{|V|}$ and an affine decoder $q:\mathbb{R}^{|V|} \rightarrow \mathbb{R}^{d_z}$ such that,
\begin{flalign}
% \begin{split} \nonumber
     \hatdeltac_V \coloneqq {r}(\deltaz) & \coloneqq \mathbf{W}_e (\deltaz - \mathbf{b}_d) + \mathbf{b}_e\,;\\
    \hatdeltaz \coloneqq {q}(\hatdeltac_V) & \coloneqq \mathbf{W}_d \hatdeltac_V + \mathbf{b}_d\,. 
% \end{split}
\label{eqn:ssae}
\end{flalign}
The representation $r(\deltaz)$ predicts $\deltac_V$, i.e., the concept shifts corresponding to $\deltaz$, with $\deltac_V = (\deltac_i)_{i\in V}$ the subvector of $\deltac$ corresponding to the index set $V$. That is, {\isae}s map differences in embedding space to their constituent concept shifts, focusing \textit{only on the varying concepts}.

We train {\isae}s to solve the following constrained problem:
\begin{align}
    (\hat{r},\hat{q
    }) \in \arg \min_{r,q} \mathbb{E}_{\x,\tilde\x} \left[ ||\deltaz - q(r(\deltaz))||^2_2\right]\ \label{eqn:recon}\\
    \text{s.t.}\ \ \mathbb{E}_{\x, \tilde\x} || r(\deltaz) ||_0 \leq \beta \label{eqn:sparse_constraint}\,,
\end{align}
where \cref{eqn:recon} is the standard auto-encoding loss that encourages good reconstruction and \cref{eqn:sparse_constraint} is a regularizer that encourages the predicted concept shift vector $\hatdeltac_V \coloneqq \hat r(\deltaz)$ to be sparse. Since the $\ell_0$-norm is non-differentiable, in practice we replace it by an $\ell_1$-norm leading to the following relaxed sparsity constraint:
\begin{flalign}
    \mathbb{E}_{\x, \tilde\x} || r(\deltaz)||_1 \leq \beta\,. %\implies \frac{1}{|V|N}\sum_{n=1}^{N} || \hatdeltac_V||_1 \leq \beta. 
    % \quad \textcolor{applegreen}{\text{(sparsity)}}
    \label{eqn:sparse_constraint_l1}
    %\label{eqn:model_sparsity}
    \end{flalign}
We then approximately solve this constrained problem by finding a saddle point of its Lagrangian using the ExtraAdam algorithm \citep{gidel2020variationalinequalityperspectivegenerative} as implemented by \citet{gallegoPosada2022cooper}. \cref{app:sparseopt} provides a detailed discussion of the benefits of constraints as opposed to penalty to regularize objectives. Appropriate normalization is crucial for enforcing sparsity using the $\ell_1$-norm. Further details, including other implementation aspects, are discussed in \cref{sec:emp} and \cref{apx:imp}. 

\parhead{Identifiability of {\isae}s.}  In \cref{sec:wscrl} we will show that, under suitable assumptions on the data-generating process and a suitable choice of $\beta$, the $\ell_0$-regularized problem of \cref{eqn:recon,eqn:sparse_constraint} is guaranteed to learn a $(\hat r, \hat q)$ such that $\hat r (\deltaz) = \mathbf{P}\D\deltac_V$ where $\D$ is an invertible diagonal matrix, $\mathbf{P}$ is a permutation matrix. In other words, the learned representation $\hat r(\deltaz)$ can be related to the ground-truth concept shift vector $\deltac_V$ (considering only the varying concepts $V$) via a permutation-scaling matrix. We will later see how sparsity regularization is crucial for this to happen. Although our theoretical analysis assumes the learned representation has size $|V|$, we find in \cref{apx:v} that, in practice, our method maintains a reasonable degree of identifiability when the representation size is larger than $|V|$. Linking identifiability back to steering, we conclude by showing how the identifiability guarantee implies that $\hat q(\mathbf{e}_k) \in \mathbb{R}^{d_z}$ are valid steering vectors for concepts in $V$.


\section{Identifiability analysis}
\label{sec:wscrl}
This section explains why we expect the representation learned in \cref{eqn:recon} to identify the ground-truth concept shift vector $\deltac_V$ up to permutation and rescaling. To do so, we first demonstrate that, under suitable assumptions, the learned representation $\hat{r}(\deltaz)$ identifies the ground-truth concept shift $\deltac_V$ \textit{up to an invertible linear transformation} when we do not use sparsity regularization. Second, we show that by adding sparsity regularization, the learned representation identifies $\deltac_V$ \textit{up to permutation and element-wise rescaling}.

Recall that, since we expect $d_c \gg d_z$, we cannot assume $\A$ to be injective; the same issue that arises when trying to encode $\c$ from $\z$. Fortunately, we do not need to make this assumption, thanks to the following decomposition. Let $\bar{V} \coloneqq [d_c] \setminus V$ be the complement of $V$. Then:
\begin{align}
     \deltaz &= \A\deltac 
     = \A_V\, \deltac_V  + \A_{\thickoverline{V}}\, \deltac_{\thickoverline{V}}\nonumber\\
     &= \A_V\, \deltac_V \ ,\label{eqn:diff_model}
\end{align}
where we used the fact that $\deltac_{\bar{V}} = 0$, by definition of $V$. By considering difference vectors, we focus on disentangling \emph{only} the varying concepts, the linear entanglement of which the submatrix $\A_V$ captures. Since $|V| \leq d_c$, we can make the assumption that mixing function $\A_V$ is injective. 
% \Cref{eqn:diff_model} can be thought of as a generative process for the \textit{difference vectors} $\deltaz$, as opposed to $\z$, as in \cref{eq:linear_entangle}. 
% A key advantage of \cref{eqn:diff_model} versus \cref{eq:linear_entangle} is that the matrix involved in it has potentially fewer columns, since $|V| \leq d_c$. This suggests that, instead of assuming $\A$ is injective, we assume its submatrix $\A_V$ is injective.
\begin{assumption}\label{ass:injective_Av}
    The matrix $\A_V \in \mathbb{R}^{d_z \times |V|}$ is injective.
    \label{ass:inj}
\end{assumption}
Note that this implies that $d_z \geq |V|$, i.e., $\z$ has at least as many dimensions as there are varying concepts. This is feasible given that $d_z$ is typically around $10^3$ (e.g., in LLMs), supporting a large set of varying concepts $V$. 

To prove linear identifiability, we will need one more assumption. Let $\Delta_V^c$ be the support of the random vector $\deltac_V$. We will require that this support is diverse enough so that its linear span is equal to the whole space $\sR^{|V|}$.
\begin{assumption}\label{ass:suff_var}
    $\text{span}(\Delta_V^c) = \sR^{|V|}$.
\end{assumption}
With these assumptions, we can show linear identifiability by reusing proof strategies that are now common in the literature on identifiable representation learning \citep{iVAEkhemakhem20a,roeder2021linear,ahuja2022weakly,xu2024sparsityprinciplepartiallyobservable}.

\begin{restatable}[\textbf{Linear identifiability}]{proposition}{linearIdent}\label{prop:linear_ident}
    Suppose $(\hat r, \hat q)$ is a solution to the unconstrained problem of \cref{eqn:recon}. Under \cref{ass:lrh,ass:injective_Av,ass:suff_var}, there  exists an invertible matrix $\mathbf{L} \in \sR^{|V|\times|V|}$ such that $\hat q = \A_V\mathbf{L}$ and $\hat{r}(\z) = \mathbf{L}^{-1}\A_V^+\z$ for all $\z \in \textnormal{Im}(\A_V)$, where $\textnormal{Im}(\A_V)$ is the image of $\A_V$.\footnote{We might not have $\hat{r}(\z) = \mathbf{L}\A_V^+\z$ for $\z \not\in \text{Im}(\A_V)$, since the behavior of $\hat{r}$ is unconstrained by the objective outside the support of $\deltaz$, i.e., outside $\text{Im}(\A_V)$.}
\end{restatable}
We prove \cref{prop:linear_ident} in \cref{app:linear_ident}. The result follows naturally from the linear representation hypothesis in \cref{ass:lrh}, but requires \cref{ass:injective_Av,ass:suff_var} for a complete proof. \citet{rajendran2024learning} prove a similar result, showing that linear subspaces of representations that represent concepts are linearly identified from concept-conditional observations.
% Even though \cref{prop:linear_ident}  may seem obvious based on \cref{ass:lrh}, recall that it is not guaranteed that a solution for $\A$ exists without \cref{ass:injective_Av} and \cref{ass:suff_var}.
% Further, compare \cref{prop:linear_ident} with the main result from \citet{rajendran2024learning}, which shows that assuming the LRH, it is possible to linearly identify $|V|$ concepts from $\mathcal{O}(|V|)$ different concept-conditional datasets. Here, we show linear identifiability is possible with $\mathcal{O}(1)$ dataset. 


%This is similar to the result by \cite{xu2024sparsityprinciplepartiallyobservable} \sj{@Sebastien could you please check this, if possible?}\seb{see comments} but differs in that here we don't require $q$ and $r$ to be invertible.


% \parhead{Is linear identifiability enough?} Recall the linear ambiguity described in \cref{sec:dgp} in the perfectly reconstructed solutions satisfying \cref{ass:recon}, i.e. $\hat{q} = \A_V \rmL^{-1}$ and $\hat{r} = \rmL \A_V$. To see how this solution affects the steering vector that we estimate for concept $k \in V$, we can write,
%    $$\lambda \hat{q}(\rve_k) = \A_V \rmL^{-1} \lambda \rve_k = \sum_{j=1}^{|V|} \lambda \rmL^{-1}_{jk} \A_{V,j}\,,$$
% where $\A_{V,j}$ is the $j$th column of $\A_V$.
% % This tells us that the spurious solutions leads to an incorrect steering vector for concept $k$, one which linearly combines the true steering vectors for each concept $j \in V$, 
% If we use this encoder to obtain a steering vector for concept $k$, the result is a vector that linearly combines the steering vectors $\A_V \rve_j$ for each concept $j \in V$.


% where $\boldsymbol{\Lambda}_V = \text{diag}(\boldsymbol{\lambda_V)}$ is a $|V| \times |V|$ diagonal matrix formed from $\blambda_V$, representing the strengths of all concept changes in $V$ when multiple concepts may be perturbed.

 % Recall the perturbed concept representation in concept $k$, $\tilde{\mathbf{c}}_{k,\lambda} = \mathbf{c} + \lambda \mathbf{e}_k$. The resulting concept representation with perturbations to multiple concepts can be expressed as $\tilde \c = \c + \bblambda \rve_k$. 

% = \sum_{j} \ell_{j,k} \A_j \sum_{j} \ell_{j,k} \A \rve_j\,,
% where $\A_j$ is the $j$th column of $\A$.






% Before developing assumptions for identifying $\A_V$, we first formalize the auto-encoding models that we analyze.
% \begin{problem}
%     \label{defn:autoenc}
%      Consider affine transformations $r: \mathcal{Z} \rightarrow \mathbb{R}^{|V|}$ and $q: \mathbb{R}^{|V|} \rightarrow \mathcal{Z}$ that represent the encoding and decoding functions, respectively, where $r$ restricted to the domain of $\deltaz$ and $q$ are injective. We estimate $\hat{q}$ and $\hat{r}$ as:
%     \begin{align}\label{prob:reconstruction}
%     \hat{q},\hat{r} = \arg \min_{r,q} \mathbb{E}_{\deltaz} \left[ l\big(\deltaz, q(r(\deltaz))\big)\right]
%     \end{align}
% \end{problem}

% % \vspace{-1mm}
% However, consider another solution $\hat{q} = \A_V \rmL$ and $\hat{r} = \rmL^{-1} \A_V^{+}$, where  $\rmL$ is an invertible $|V| \times |V|$ matrix. Since $\rmL^{-1} \rmL = \mathbf{I}$, this solution also perfectly reconstructs $\deltaz$, and is a minimum of the objective in \cref{defn:autoenc}. This tells us that the auto-encoding model in \cref{defn:autoenc} is only \emph{identified up to linear transformations}. That is, all solutions to the target objective are linear transformations of the true mapping $\A_V$.



% This implies that $\A_V \rve_k = \A_k = \A \rve_k$, or that the decoder applied to the canonical basis element $\rve_k$ would enable steering by $\lambda \A_k$ .


% \parhead{Is linear identifiability enough?} Consider the spurious solution $\hat{q} = \A_V \rmL$ and $\hat{r} = \rmL^{-1} \A_V^{+}$. Rather than mapping each observed difference vector $\deltaz$ to its concept representation, this spurious solution maps $\deltaz$ to a different representation that linearly entangles concepts with the matrix $\rmL$. To see how this solution affects the steering vector that we estimate for concept $k \in V$, we can write,
%    $$\hat{q}(\rve_k) = \sj{\lambda} \A_V\rmL\rve_k = \sum_{j} \ell_{j,k} \A_j =  \sum_{j} \ell_{j,k} \A \rve_j\,,$$
% where $\A_j$ is the $j$th column of $\A$.
% This tells us that the spurious solutions leads to an incorrect steering vector for concept $k$, one which linearly combines the true steering vectors for each concept $j \in V$, 
% If we use the spurious encoder to obtain a steering vector for concept $k$, the result is a vector that linearly combines the steering vectors $\A \rve_j$ for each concept $j \in V$.
% Thus, recovering steering vectors requires learning the ``unmixing'' matrix $\rmL^{-1}$, which would require supervision via concept labels or methods like linear independent component analysis which make strict assumptions \citep{comon1994independent}.

%We instead require a stronger notion of identifiability that will allow us to extract steering vectors with minimal human feedback. 
%\begin{definition}[Permutation-scaling identifiability]
%\label{defn:perm_idnt}
%    If every solution to the problem in \eqref{prob:reconstruction_sparse} has the form $\hat{q} = \A_V \D\rmP$ and $\hat{r} = \A_V^{+} \D\rmP$, where $\rmP$ is a $|V| \times |V|$ permutation matrix and $\D$ is an invertible diagonal scaling matrix of the same size, then the model in \cref{prob:reconstruction_sparse} is said to be identifiable and the transformation $\A_V$ is identified up to permutation and scaling.   
%\end{definition}
% This notion of identifiability differs slightly from the standard definition in identifiable representation learning (see \cref{sec:rel}), which focuses on identifying the encoding function $\hat r = \D\rmP \A_V^+$ rather than the decoding function. While the theorem in this paper shows that we can also identify the true encoding function $\A_V^+$, we focus on the true decoding matrix $\A_V$ in the definition since this is the key to obtaining steering vectors.

\parhead{Identifiability up to permutation and rescaling.} To go from identifiability up to linear transformation to identifiability up to permutation and rescaling, we need to make further assumptions. Let $\mathcal{S}$ be the support of the distribution $p(S)$, i.e., $\mathcal{S} \coloneqq \{S \subseteq [d_c] \mid p(S) > 0\}$. The following is based on \citet{lachapelle2023synergies} and \citet{xu2024sparsityprinciplepartiallyobservable}.

% \parhead{Identifiability result sketch.} To identify the auto-encoding model introduced in \cref{prob:reconstruction_sparse}, we leverage a key assumption about \emph{sparse encoding}. While formalized below, at a high-level, we constrain the model to encode each observed difference vector $\deltaz$ as sparsely as possible. The model can only satisfy the sparsity constraint with a latent space that reflects the true varying concepts, where the changes to concepts can indeed be represented as sparsely as possible. For this intuition to lead to identifiability, we require other assumptions, including one that requires the multi-concept changes to vary sufficiently, so that there are \textbf{``diverse enough"} concept variations across all data pairs in the dataset $\mathcal{D}$ (in \cref{defn:varying_concepts}). 
% The subvector $\deltac_S$ corresponds to the shifts in concept values for the concepts in $S$.

% Since the first two assumptions have been motivated already, we introduce them succinctly.
% \begin{assumption}
%     (Injectivity): We assume $q : \mathbb{R}^{|V|} \rightarrow \mathcal{Z}$ to be an injective, affine transformation. Restricted to the domain of the observed $\deltaz$, we assume $r : \mathcal{Z} \rightarrow \mathbb{R}^{|V|}$ to be injective and affine.
% \label{ass:inj}
% \end{assumption}

% \begin{assumption}
% \textcolor{tiffanyblue}{(Perfect reconstruction)}: Under (\cref{ass:inj}), we can perfectly reconstruct $\deltaz$:
% \begin{flalign}
% \label{eqn:recon}
%     \mathbb{E} ||\deltaz - \hatdeltaz ||_2^2 = 0; \quad \hatdeltaz \coloneqq \hat{q}(\hat{r}(\deltaz)).
% \end{flalign}
% \label{ass:recon}
% \end{assumption}
% \vspace{-5mm}

% The next assumption is to ensure \textbf{``diverse enough"} concept variations across all the observed pairs in the dataset $\mathcal{D}$ (in \cref{defn:varying_concepts}).
% The random variable $S \subseteq V$ captures the set of concepts that vary across a single paired observation so that $S \coloneqq \{k \in V | \deltac_k \neq 0\}$. The support $\mathcal{S}$ of the random variable $S$ is $\mathcal{S} \coloneqq \{s \subseteq V \mid p(s) > 0\}$ where $s$ is the realization of $S$ for a specific pair. 

\begin{assumption}[Sufficient diversity of multi-concept shifts]
\label{ass:suffsupp}
The following two conditions hold.
\begin{enumerate}
    \item (Sufficient support variability): For every varying concept $k \in V$, we have
    % \vspace{-1mm}
    \begin{flalign}
        \bigcup_{S \in \mathcal{S} | k \notin S} S = V \setminus \{k \} \quad \forall k \in V \,;
        \label{eqn:suffsupp}
    \end{flalign}
    %\vspace{-4mm}
    \item (Distribution $\mathbb{P}_{\deltac_S | S}$ continuous): For all $S \in \mathcal{S}$, the conditional distribution $\mathbb{P}_{\deltac_S | S}$ can be described using a probability density with respect to the Lebesgue measure on $\mathbb{R}^{|S|}$.
\end{enumerate}
\end{assumption}

% \seb{Could be nice to give an example of support $\mathcal{S}$ that satisfies the assumption with concepts that we think makes sense for text. We could assume something like $V \coloneqq \{\texttt{language}, \texttt{sentiment}, \texttt{verbosity}\}$ and assume maybe 
% \begin{align*}
% \mathcal{S} :=\ \{&\{\texttt{language, sentiment}\}, \\
% &\{\texttt{language, verbosity}\}, \\
% &\{\texttt{sentiment, verbosity}\} \} \, , 
% \end{align*}
% and give actual examples of pairs of sentences for these three multi-concept shifts. That would help making the discussion more concrete.}

% Intuitively, without the first assumption, two concepts $k, j \in  V$ may never vary separately across all pairs, making it impossible for the model to disentangle them. \Cref{ass:suffsupp} accommodates a wider range of scenarios, including the potential for statistically dependent concepts, which we study empirically in \cref{sec:emp}. As an example, the special case where only one concept varies within a pair but all concepts vary at least once across all pairs satisfies this more general assumption. On the other end, it rules out the case where all the concepts vary within every pair, i.e. $p(S = V) = 1$.
Without the first assumption, two concepts $k,j\in V$ might always change together, meaning there is no data pair in which only one of them varies independently. Intuitively, this would prevent the model from disentangling them effectively. Importantly, our assumption accommodates a broad range of scenarios. E.g., it is not necessarily violated even in an extreme case where $|V| - 1$ concepts change in each pair. Moreover, it allows for the presence of statistically dependent concepts.
The second criterion ensures the distribution $\mathbb{P}_{\deltac_S | S = s}$ does not concentrate mass on a subset of $\sR^{|S|}$ of Lebesgue measure zero. In \cref{apx:suffsupp}, we provide examples of distributions that meet or fail the assumption.\footnote{See \citet{lachapelle2023synergies} for a strictly weaker but more technical assumption that is also sufficient for \cref{prop:perm_ident}.} %Next, we define the ideal structure of an estimated concept vector, making the notion of sparsity precise. 

% The sufficient variability assumption does not prevent statistical dependencies between concepts, a setting that we explore empirically in \cref{sec:emp}.
% \sj{proper or lower-D subspace? proper needn't be lower D}

%\begin{definition}     \textcolor{applegreen}{(Sparse concept vectors)}. An estimated concept vector is considered to be sparse when it satisfies the following constraint: $\mathbb{E} || \hatdeltac_S ||_0 \leq \mathbb{E} || \deltac_S ||_0$, referred to as the \textbf{sparsity constraint}.
%    \label{eqn:sparsity}
%\end{definition}

We are now ready to state the main identifiability result of this section. We note that its proof relies to a large extent on an existing result by \citet{lachapelle2023synergies}.

\begin{restatable}[\textbf{Identifiability up to permutation}]{proposition}{permIdent}\label{prop:perm_ident}
    Suppose $(\hat r, \hat q)$ is a solution to the constrained problem of \cref{eqn:recon,eqn:sparse_constraint} with $\beta = \mathbb{E}||\deltac_V||_0$. Under \cref{ass:lrh,ass:injective_Av,ass:suff_var,ass:suffsupp}, there  exists an invertible diagonal matrix and a permutation matrix $\mathbf{D}, \mathbf{P} \in \sR^{|V|\times|V|}$ such that $\hat q = \A_V \mathbf{D}\mathbf{P}$ and $\hat{r}(\z) = \mathbf{P}^\top\mathbf{D}^{-1}\A_V^+\z$ for all $\z \in \textnormal{Im}(\A_V)$, where $\textnormal{Im}(\A_V)$ is the image of $\A_V$.
\end{restatable}

\parhead{Proof sketch.} We outline the proof here and defer the full details to \cref{app:proof}. We first show that all optimal solutions of the constrained problem must reach a reconstruction loss of zero. This means that optimal solutions to the constrained problem are also optimal for the unconstrained one. Thus, these solutions must identify $\A_V$ up to linear transformation, by \cref{prop:linear_ident}. We can then rewrite the constraint as $\mathbb{E} || \mL^{-1}\deltac_V ||_0 \leq \beta = \mathbb{E}||\deltac_V||_0$. Here, we can reuse an argument initially proposed by \citet{lachapelle2023synergies} to leverage this inequality to conclude that $\mL^{-1}$ must be a permutation-scaling matrix. For completeness, we present this argument in \cref{lemma:synergies_proof}. It shows that, applying the matrix $\mL^{-1}$ to $\deltac_V$ always strictly increases its expected sparsity, \textit{unless} $\mL^{-1}$ is a permutation-scaling matrix. Thus, to satisfy the inequality, $\mL$ must be a permutation-scaling matrix.

\parhead{Extracting steering vectors.} Under \cref{ass:lrh,ass:injective_Av,ass:suff_var,ass:suffsupp}, \cref{prop:perm_ident} shows that $\hat q = \A_V \mathbf{DP}$. From this identifiability result, we can see that,
\begin{flalign*}
    \z + \hat q(\mathbf{e}_k) = \A\c + \mathbf{D}_{\pi(k), \pi(k)}\A\mathbf{e}_{\pi(k)} \\
    = \A(\c + \lambda \mathbf{e}_{\pi(k)}) = f(g(\c + \lambda \mathbf{e}_{\pi(k)}))  = f(g(\tilde \c_{\pi(k), \lambda} )) \,,
\end{flalign*}
where $\lambda := D_{\pi(k), \pi(k)}$. In other words, when we add the decoded basis vector $\mathbf{e}_k$ to any embedding $\z$, i.e., add the $k$-column of the linear decoding matrix, the resulting vector represents $f(g(\tilde \c_{\pi(k), \lambda}))$, the embedding representation of the $\pi(k)$-th concept steered. Thus, identifiability directly leads to accurate unsupervised steering. A practitioner can use this result to try each steering vector $q(\mathbf{e}_k)$ in turn, generate tokens with an LLM, and directly interpret the changes to interpret the concept that was steered. By contrast, in \cref{apx:linid_insuff}, we show that linear identifiability is insufficient to recover steering vectors up to permutation without the need for further labelled examples.

The novelty of our contribution is in connecting the theory of identifiable representation learning to steering vector discovery. While standard SAEs employ sparsity as an inductive bias without formal recovery guarantees, we show that sparsity when considered for shifts between concepts, yields provable identifiability. This ensures that decoder columns correspond to valid steering vectors for individual concepts.

\section{Empirical Studies}
\label{sec:emp}
The central challenge in mechanistic interpretability is empirical reliability. Sparse autoencoders (SAEs) discover features with low reconstruction error, but whether these represent stable internal computations or artifacts of a particular training run or optimisation remains unclear---causing steering to fail unpredictably. Our theory argues this failure is structural: without identifiability, steering directions are ill-defined, allowing multiple incompatible explanations to remain equally consistent with the data. We evaluate this claim across controlled benchmarks and real-world language datasets. Our experiments address two questions:
\begin{enumerate}[leftmargin=*, topsep=2pt, itemsep=2pt]
\item \textbf{Theoretical Validity.} Do \isae{}s recover latent directions uniquely (up to permutation and scaling)? (\cref{subsec:validation})
\item \textbf{Practical Consequences.} Does identifiability improve steering accuracy and robustness to distribution shift? (\cref{subsec:utility})
\end{enumerate}

\textbf{Setup.} We extract final-layer embeddings from Gemma-2-2B \citep{anil2024gemma2} and Pythia-70M \citep{biderman2023pythia} on contrastive prompt pairs $(\x, \tilde \x)$ from semi-synthetic data (\cref{tab:datasets}) and real-world datasets: Bias in Bios~\citep{De_Arteaga_2019}, refusal and sycophancy~\citep{panickssery2024steeringllama2contrastive}, and TruthfulQA~\citep{lin2022truthfulqameasuringmodelsmimic} \footnote{We evaluate on existing steering benchmarks with explicit contrastive pairs. In practice, pairs can be constructed by uniformly sampling any two prompts from the dataset.}. 

\begin{tcolorbox}[
  enhanced jigsaw,
  breakable,
  colback=gray!6,
  colframe=gray!60!black,
  title=\textsc{Identifiable Steering Desiderata},
  boxrule=0.5pt,
  arc=2pt,
  left=4pt, right=4pt, top=2pt, bottom=2pt
]
Identifiability ensures that a learned steering direction corresponds to a \emph{unique underlying concept}, up to irreducible symmetries (permutation and rescaling). As a result, steering directions reflect stable structure that can be transferred across settings and we expect the following empirical consequences:
\begin{enumerate}[leftmargin=*, topsep=2pt, itemsep=1pt]
\item On synthetic or semi-synthetic benchmarks with known generative factors, identifiable methods recover ground-truth concepts up to permutation and scaling.
\item Higher identifiability scores correlate with improved downstream steering performance.
\item Reliable out-of-distribution steering requires identifiability. Non-identifiable directions may succeed in-distribution but fail under distribution shift.
\end{enumerate}
\end{tcolorbox}

\begin{figure}[ht]
    \centering
    \includegraphics[width=\linewidth]{bios_summary.png}
    \caption{\textbf{{\isae}s achieve earlier and stronger transitions to female indicators in generated text} (by strength $1.0$), while GemmaScope2B requires stronger interventions.}
    \label{fig:bias-stats}
\end{figure}

\textbf{Metrics.} We evaluate two aspects of the learned representations. We measure identifiability via the Mean Correlation Coefficient (MCC)~\citep{hyvarinen2016unsupervisedfeatureextractiontimecontrastive}, which equals 1.0 when learned dimensions aligns perfectly with the concept labels up to permutation and scaling. We measure steering accuracy via cosine similarity between the predicted steering vector, $\hat{q}(\rve_{\pi(k)})$, and the true embedding difference $f(\tilde{\x}_k) - f(\x)$ on held-out pairs differing in a single concept $k$. This enables comparison in a manner invariant to the scale of the strength at which steering vectors are applied to the representation. Details in \cref{apx:metrics}. 
\begin{table*}[ht]
\centering
\caption{SSAEs are consistently more identifiable across datasets as seen by the MCC across unsupervised baselines, with the difference being most pronounced when latent concepts are correlated, as in \textsc{corr}$(2,1)$. MCC scores using Pythia are reported in \cref{tab:mcc_lang_pythia}.
}
\label{tab:mcc_lang_gemma}
\small
\setlength{\tabcolsep}{6pt}
\renewcommand{\arraystretch}{1.05}
\begin{tabular}{l|c|c|c|c|c|c}
\toprule
\textbf{Dataset} 
& \textbf{SSAE} 
& \textbf{GemmaScope} 
& \textbf{TopK-SAE} 
& \textbf{ReLU-SAE} 
& \textbf{JumpReLU SAE}
& \textbf{Linear Probe} \\
\midrule
\textsc{lang}$(1,1)$     
& $0.9121 \pm 0.0180$ 
& $0.8614$ 
& $0.8467 \pm 0.0310$ 
& $0.8325 \pm 0.0280$ 
& $0.8219 \pm 0.0340$ 
& $0.8793$ \\
\hdashline
\textsc{gender}$(1,1)$   
& $0.9018 \pm 0.0210$ 
& $0.8542$ 
& $0.8421 \pm 0.0290$ 
& $0.8297 \pm 0.0300$ 
& $0.8183 \pm 0.0360$ 
& $0.8675$ \\
\hdashline
\textsc{binary}$(2,2)$   
& $0.8896 \pm 0.0240$ 
& $0.8387$ 
& $0.8269 \pm 0.0350$ 
& $0.8128 \pm 0.0330$ 
& $0.8015 \pm 0.0410$ 
& $0.8562$ \\
\hdashline
\textsc{corr}$(2,1)$     
& $0.9047 \pm 0.0207$ 
& $0.5760$ 
& $0.5542 \pm 0.0660$  
& $0.5132 \pm 0.0510$ 
& $0.5030 \pm 0.0810$ 
& $0.7746$ \\
\hdashline
TruthfulQA           
& $0.7585 \pm 0.0182$ 
& $0.7370$ 
& $0.7281 \pm 0.0294$ 
& $0.7128 \pm 0.0309$ 
& $0.7128 \pm 0.0309$ 
& $0.7758$ \\
\hdashline
Sycophancy                
& $0.7070 \pm 0.0450$ 
& $0.6344$ 
& $0.5896 \pm 0.0634$ 
& $0.6005 \pm 0.0245$ 
& $0.5177 \pm 0.0080$ 
& $1.0000$ \\
\hdashline
Refusal                  
& $0.6233 \pm 0.0401$ 
& $0.6119$ 
& $0.5686 \pm 0.0511$ 
& $0.6116 \pm 0.0604$ 
& $0.5699 \pm 0.0143$ 
& $0.9000$ \\
\hdashline
Bias-in-Bios              
& $0.8232 \pm 0.2061$ 
& $0.7385$  
& $0.5660 \pm 0.2296$ 
& $0.7657 \pm 0.1065$ 
& $0.7773 \pm 0.0898$ 
& $0.9731$ \\
\bottomrule
\end{tabular}
\end{table*}

\textbf{Baselines.} We compare against GemmaScope \citep{lieberum2024gemma} and PythiaSAE \citep{eleutherai2023saepythia}---SAEs trained on large-scale datasets without identifiability guarantees---as well as other commonly used SAE variants that can be trained from scratch. These include ReLU SAEs \citep{cunningham2023sparseautoencodershighlyinterpretable, bricken2023monosemanticity}, JumpReLU SAEs \citep{rajamanoharan2024improvingdictionarylearninggated}, and TopK SAEs \citep{gao2024scaling}. We also include linear probes as a supervised baseline with direct label access. We implement \isae{}s (\cref{eqn:ssae}) with the sparsity constraint (\cref{eqn:sparse_constraint_l1}) using the \texttt{cooper} library \citep{gallegoPosada2022cooper}. Details in \cref{apx:imp}. Reconstruction errors across all runs and datasets are reported in \cref{apx:recon}. 
\begin{figure}[t]
\centering
\includegraphics[width=\linewidth]{cosplotood.png}
        \caption{Embeddings steered using \isae show \textbf{higher OOD generalisation performance}.}
        \label{fig:cos_ood}
\end{figure}
\subsection{How well does \isae identify steering vectors?}
\label{subsec:validation}
We first test whether \isae{}s recover concepts up to permutation and rescaling as predicted by \cref{prop:perm_ident}, and how the guarantees hold up against models that were not designed to be identifiable. First, we we consider simple semi-synthetic datasets with single words where we can assume the number of underlying concept variations in pairs $(\x, \tilde \x)$ with a diverse range of concept variations, named as: identifier of the dataset indicating why we consider it, followed by $|V|$ and max$(|S|)$: \textsc{identifier}($|V|$, max$|S|$). Details on datasets can be found in \cref{apx:data}. Briefly, \textsc{lang}($1, 1$) (e.g., \textit{eng} $\rightarrow$ \textit{french}) and
\textsc{gender}($1,1$) (e.g., \textit{masculine} $\rightarrow$ \textit{feminine} vary a single concept between $\x$ and $\tilde \x$. Next, we stress-test the viability of our assumptions on the real-world datasets mentioned earlier. For details on all datasets and $(\x, \tilde \x)$ pair creation, please refer to \cref{apx:data}.

\cref{tab:mcc_lang_gemma} shows that \isae achieves consistently high MCC values, empirically corroborating \cref{prop:perm_ident}. We also include a similar analysis in \cref{apx:empirical-pythia} with \isae{}s and all baselines trained on Pythia embeddings, revealing a similar pattern. In particular, the guarantee of \isae{} to identify concepts, even when they are correlated, enables the identifiability of steering vectors on the dataset \textsc{corr}$(2,1)$. Next, we evaluate whether the benefits of a higher $\text{MCC}$ translate to performance improvements on steering embeddings to be more similar to those of the target concept.

\begin{figure*}
\centering
\includegraphics[width=\linewidth]{cosplot1.png}
        \caption{\textbf{A higher MCC value is associated with a greater cosine similarity.} Embeddings steered with vectors from a more disentangled decoder align more closely to target embeddings.}
        \label{fig:cos}
\end{figure*}

\subsection{Practical Implications of Identifiability for Steering}
\label{subsec:utility}
% \textbf{Does identifiability translate to \textit{better} steering?}  
% % Specifically, for the experiments in \cref{fig:cos}, 
We hold out pairs $(\x, \tilde \x_k) \forall k \in V$, each varying by a single concept, and compare the cosine similarity between the steered embeddings and target embeddings. 
\cref{fig:cos} illustrates that {\isae}'s higher $\text{MCC}$  performance generally translates to more accurate steering, with significant advantages over all related methods in the more challenging \textsc{binary}($2, 2$) and \textsc{corr}($2, 1$) settings where multiple or correlated concepts change. \cref{fig:cos} also reveals that even slight differences in $\text{MCC}$ values can translate into pronounced variations in steering accuracy. Next, we evaluate \emph{out-of-distribution} (OOD) steering accuracy, based on the hypothesis that steering vectors that disentangle a single concept should transfer to different domains.

\textbf{OOD generalisation.} For this evaluation, we learn a steering vector from \textit{eng} $\rightarrow$ \textit{french} using the \textsc{binary}($2, 2$) or \textsc{corr}($2, 1$) dataset, where language changes are shown for occupation-related works, and use the steering vector on the \textsc{lang}($1, 1$) dataset consisting of words related to household objects.  
\cref{fig:cos_ood} shows that the steering vectors learned by {\isae}s transfer effectively to OOD datasets while SAEs do not perform better than simple baselines, further substantiating the importance of identifiability for unsupervised steering.

\textbf{Encoding dimension.} While identifiability theory requires assuming that the encoder dimension is known (equal to $|V|$, \cref{eqn:ssae}), in all the experiments with language models, we've considered the encoding dimension to be as large as the embedding dimension, since it's difficult to assume a known number of concepts in real-world data comprising of sentences and paragraphs. Another design choice is to use the last layer's embeddings for steering since the LRH--a key component of the identifiability results---is better theoretically motivated (c.f. \citep{marconato2024all}) at the last layer. To test sensitivity to these assumptions, we conducted further studies training {\isae}s on different layers of an LLM (\cref{apx:layer}), and studying the variation of $\text{MCC}$ versus steering accuracy as we increase the encoding dimension size past $|V|$, finding that for encoding dimensions $>|V|$, there is an increase in steering accuracy even though $\text{MCC}$ values drop substantially (see \cref{apx:v} for details). These promising findings provide grounding for larger encoding dimensions (than $|V|$) as considered for the experiments in the paper. 

\textbf{Steering and open-ended generation.}  We evaluate steering capabilities using
   the Bias in Bios dataset \citep{De_Arteaga_2019}, which contains biographical
  text annotated with gender and profession. This dataset allows us to
  systematically measure the effects of steering by counting gender pronouns
  in generated text. We identify the decoder
  column with maximum correlation with the binary gender labels across
  the dataset. We use this decoder
  column directly as the steering vector. During text generation with Gemma-2B, we apply
   a forward hook at the last layer (25) that adds the steering vector (scaled by a
  steering strength coefficient) to the last token's hidden state before
  continuing generation. We generate text from prompts associated with
  male-biased professions (e.g., ``The CEO of the tech startup
  announced...") and count gendered pronouns in the output, classifying each
   generation as male-dominated, female-dominated, or neutral.
  \cref{fig:bias-stats} shows that {\isae} steering vectors lead to more
  effective generation of female pronouns than those of GemmaScope,
  indicating that {\isae}s learn decoder columns better aligned with the
  underlying gender concept.

\noindent
\textbf{Limitations.} We do not claim that \isae{}s universally outperform SAEs. Rather, \isae{}s make explicit when steering directions are identifiable, and we provide preliminary empirical evidence supporting this claim. Crucially, \isae{}s are trained on the same underlying datasets used for SAEs and do not require specialised or manually constructed paired samples in practice; they instead exploit naturally occurring variation by operating on differences between samples. We omit large-scale benchmarks and dashboards by design since our contribution is theoretical clarity.

\section{Related work}
\label{sec:rel}
% \vspace{-0.1cm}
\textbf{Linear representation hypothesis}. This paper builds on the linear representation hypothesis that language models encode concepts linearly. Several papers provide empirical evidence for this hypothesis \citep{mikolov-etal-2013-linguistic, gittens-etal-2017-skip, ethayarajh2019understandinglinearwordanalogies, allen2019analogiesexplainedunderstandingword, seonwoo-etal-2019-additive, burns2024discoveringlatentknowledgelanguage, li2024inferencetimeinterventionelicitingtruthful, moschella2023relativerepresentationsenablezeroshot, tigges2023linearrepresentationssentimentlarge, nissim2019fairbettersensationalmandoctor, ravfogel-etal-2020-null,park2023linear, park2024geometrycategoricalhierarchicalconcepts}. Recent work also provides theoretical justification for why linear properties might consistently emerge across models that perform next-token prediction \citep{roeder2021linear, jiang2024originslinearrepresentationslarge, marconato2024all}.  

\textbf{Interpretability of LLMs}. This paper contributes to the literature on interpretability and steering of LLMs. Much work on finding concepts in LLM representations for steering relies on supervision, either from paired observations with a single-concept shift \citep{panickssery2024steeringllama2contrastive,turner2024steeringlanguagemodelsactivation,rimsky2024steering,li2024inferencetimeinterventionelicitingtruthful} or from examples of target LLM completions to prompts \citep{subramani2022extracting}. This prior work also focuses on applying the same steering vector to all examples, implicitly relying on the linear representation hypothesis as justification. In contrast, we make the assumption precise, and show how it leads to steering vectors. This paper also departs from supervised learning and focuses on learning with limited supervision. In this way, we propose a method that is similar to sparse autoencoders (SAEs) \citep{templeton2024scaling, engels2024languagemodelfeatureslinear, cunningham2023sparseautoencodershighlyinterpretable, rajamanoharan2024improvingdictionarylearninggated, gao2024scaling}. In contrast, our proposed method fits concept shifts, and provably identifies steering vectors while SAEs may not enjoy identifiability guarantees.

\textbf{Causal representation learning}. Finally, this paper builds on causal representation learning results that leverage sparsity constraints. \citet{ahuja2022weakly}, \citet{locatello2020weakly}, and \citet{brehmer2022weaklysupervisedcausalrepresentation} consider sparse latent perturbations and paired observations. In contrast, we focus on learning from multi-concept shifts. \citet{lachapelle2022disentanglement} focus on sparse interventions and sparse transitions in temporal settings, while \citet{lachapelle2023synergies}, \citet{layne2024sparsityregularizationtreestructuredenvironments}, \citet{xu2024sparsityprinciplepartiallyobservable}, and \citet{fumero2023leveragingsparsesharedfeature} leverage sparse dependencies between latents and tasks. In this paper, we adapt these assumptions and technical results for a novel setting: discovering steering vectors from LLM representations based on concept shift data. 
In work that is closest to ours, \citet{rajendran2024learning} recover linear subspaces that capture concepts up to linear transformations using concept-conditional datasets, and \citet{goyal2025causaldifferentiating} develop an identifiable contrastive learning approach to discover behavior-mediating concepts, but cannot extract steering vectors.
In contrast, we focus on multi-concept shifts and how these lead to identifiable steering vectors. 

\textbf{Sparse coding and dictionary learning.} Classical dictionary learning provides identifiability guarantees under geometric constraints on the dictionary by leveraging properties such as the restricted isometry property (RIP) \citep{candes2005decodinglinearprogramming, baraniuk2008simple} and incoherence conditions \citep{donoho2003optimally, gribonval2015sample} (\cref{apx:dict-vs-ssae}). These assumptions ensure sparse codes are uniquely recoverable, but require approximately uncorrelated dictionary columns, an assumption which might not hold for LLM representations where concepts like truthfulness and harmlessness are intrinsically correlated. A further distinction is amortisation. Standard SAEs learn an encoder to approximate instance-wise sparse inference. Recent analysis \citep{oneill2025computeoptimalinferenceprovable} shows this introduces a provable amortisation gap: linear-nonlinear encoders cannot implement optimal sparse inference even when the dictionary is recoverable.

\section{Conclusion}
We propose Sparse Shift Autoencoders ({\isae}s) for discovering accurate steering vectors from multi-concept paired observations as an alternative to both SAEs, and approaches relying on supervised data. Key to this result are the identifiability guarantees that the \isae enjoys as a consequence of considering sparse concept shifts. We study the \isae empirically on several real language tasks, and find evidence that the method facilitates accurate steering learned via limited supervision. However, we stress that these experiments are intended to validate the identifiability results in \cref{sec:wscrl} and their implications for accurate steering. Although we include effects of steering on generated text (\cref{fig:bias-stats}), to fully understand the impacts of the \isae on steering research, especially LLM alignment, more evaluation is needed on embeddings from more complex datasets, and on more challenging tasks. Large-scale real-world evaluation remains an important direction for future work.

\section*{Impact Statement}
\label{sec:impact}
This paper presents technical advancements to a new field of machine learning focused on steering the behaviour of large language models at inference time, i.e., without requiring access to the model's parameters. Steering methods have already begun to play a role in the alignment of LLMs to be e.g., more truthful. We present a new method that could speed up steering research by allowing practitioners to recover steering vectors without the need for supervision, a previous limitation of steering methods. As such, this work could have a positive impact on LLM safety and alignment research. Nevertheless, we flag that contributions towards steering such as ours should be empirically evaluated carefully to avoid over-claiming LLM safety. We acknowledge that while the empirical studies we conduct demonstrate the advantages of identifiable methods such as \isae for steering, further evaluation is necessary to the method's use in AI safety research.

\bibliography{iclr2026_conference}
\bibliographystyle{icml2026}

%%%%%%%%%%%%%%%%%%%%%%%%%%%%%%%%%%%%%%%%%%%%%%%%%%%%%%%%%%%%%%%%%%%%%%%%%%%%%%%
%%%%%%%%%%%%%%%%%%%%%%%%%%%%%%%%%%%%%%%%%%%%%%%%%%%%%%%%%%%%%%%%%%%%%%%%%%%%%%%
% APPENDIX
%%%%%%%%%%%%%%%%%%%%%%%%%%%%%%%%%%%%%%%%%%%%%%%%%%%%%%%%%%%%%%%%%%%%%%%%%%%%%%%
%%%%%%%%%%%%%%%%%%%%%%%%%%%%%%%%%%%%%%%%%%%%%%%%%%%%%%%%%%%%%%%%%%%%%%%%%%%%%%%
\newpage
\appendix
\onecolumn
\epigraph{..we understand the world by studying change, not by studying things..}{As quoted in the Order of Time, Anaximander}
\tableofcontents
\newpage
\section{Theory}
\label{apx:theory}
\subsection{Notation and Glossary}
\label{apx:gloss}
\vspace{-0.1cm}
\centerline{\bf General notation}
\vspace{-0.2cm}
\centerline{\rule{0.95\linewidth}{0.5pt}}
\bgroup
\def\arraystretch{1.5}
%\vspace{0.2cm}
%\centerline{\textit{}}
\begin{tabular}{>{\centering\arraybackslash}p{2.25in} >{\arraybackslash}p{3.15in}}

$\displaystyle k$ & integer\\
$\displaystyle [k]$ & set of all integers between $1$ and $k$, inclusively\\
$\displaystyle S \subseteq [k]$ & set\\
$\displaystyle |S|$ & cardinality of a set\\
$\displaystyle S \backslash S'$ & set subtraction (set of elements of $S$ that are not in $S'$)\\ 
$\displaystyle \lambda$ & scalar\\ 
%$\displaystyle \mathcal{S}$ & domain set\\
$\displaystyle \x$ & vector and vector-valued random variables\\
$\displaystyle x_k$ & element $k$ of a random vector $\x$ \\
$\displaystyle \x_S$ & subvector with element $x_i$ for $i \in S$\\
$\displaystyle \A$ & matrix\\
$\displaystyle \A_{i,j}$ & element $i, j$ of matrix $\A$ \\
$\displaystyle \A_{:, i}$ & column $i$ of matrix $\A$ \\
$\displaystyle \A_{S}$ & matrix with columns $\A_{:,j}$ for $j \in S$\\
$\displaystyle \A^+$ & pseudo-inverse of a matrix $\A$ \\
$\displaystyle \rve_k \in \mathbb{R}^n$ & standard basis vector of the form $[0,\dots,0,1,0,\dots,0]$ with a 1 at position $k$\\ 
$\displaystyle f: \X \rightarrow \Z$ & function $f$ with domain $\X$ and codomain $\Z$\\
$\displaystyle f \circ g $ & composition of the functions $f$ and $g$ \\
$\displaystyle || \x ||_p $ & $\ell_p$ norm of $\x$ \\ [2ex]
%$\displaystyle \{k\}$ & the set with the single element $k$ \\

%$\displaystyle \mathbb{R}$ & set of real numbers\\


% \egroup
% \vspace{0.5cm}
% \bgroup
% \def\arraystretch{1.5}
%\end{tabular}
% \egroup
% \vspace{0.5cm}
% \bgroup
% \def\arraystretch{1.5}
%\vspace{0.5cm}
%\centerline{\textit{Calculus}}
%\begin{tabular}{>{\centering\arraybackslash}p{3.25in} >{\arraybackslash}p{3.15in}}
% NOTE: the [2ex] on the next line adds extra height to that row of the table.
% Without that command, the fraction on the first line is too tall and collides
% with the fraction on the second line.
% $\displaystyle\frac{d y} {d x}$ & Derivative of $y$ with respect to $x$\\ [2ex]
$\displaystyle \frac{\partial y} {\partial x} $ & partial derivative of $y$ with respect to $x$ \\ [2ex]
$\displaystyle \nabla_\vx f(\vx) \in \R^{m\times n}$ & Jacobian matrix of $f: \R^n \rightarrow \R^m$\\ [2ex]
$\displaystyle \nabla_\vx^2 f(\vx) \in \R^{n\times n}$ & Hessian matrix of $f: \R^n \rightarrow \R$\\
%\end{tabular}
%\vspace{0.25cm}

%\centerline{\textit{Probability}}
%\begin{tabular}{>{\centering\arraybackslash}p{3.25in} >{\arraybackslash}p{3.15in}}
$\displaystyle \mathbb{P}$ & probability measure/distribution \\
%$\displaystyle p(S)$ & probability distribution over a continuous variable, or over a variable whose type has not been specified\\
$\displaystyle  \E_{\x} [ f(\x) ]$ & expectation of $f(\x)$ with respect to $\x$ \\
\end{tabular}
% \egroup
\vspace{0.25cm}

%\centerline{ \textit{Functions}}
% \bgroup
% \def\arraystretch{1.5}
%\begin{tabular}{>{\centering\arraybackslash}p{3.25in} >{\arraybackslash}p{3.15in}}

%\end{tabular}
\egroup
\newpage
\centerline{\bf Glossary}
\vspace{-0.2cm}
\centerline{\rule{0.75\linewidth}{0.5pt}}
\bgroup
\def\arraystretch{1.5}
\begin{tabular}{>{\centering\arraybackslash}p{2.25in} >{\arraybackslash}p{3.15in}}
$\displaystyle \x \in \R^{d_x}$ & observation \\
$\displaystyle \z \in \R^{d_z}$ & pretrained representation\\
$\displaystyle \c \in \R^{d_c}$ & ground-truth concept vector \\
%$\displaystyle k$ & concept index\\
%$\displaystyle \lambda$ & strength of concept variation\\
$\displaystyle \tilde \c_{k, \lambda}$ & ground-truth concept vector after varying concept $k$ by $\lambda$ from $\c$\\
$\displaystyle \tilde \x_{k, \lambda}$ & observation corresponding to  $\tilde \c_{k, \lambda}$\\
$\displaystyle \tilde \z_{k, \lambda}$ & pretrained representation corresponding to  $\tilde \c_{k, \lambda}$\\
%$\displaystyle d_x$ & dimension of observations \\
%$\displaystyle d_z$ & dimension of learned representations\\
%$\displaystyle d_c$ & dimension of concept representations\\
$\displaystyle \X \subseteq \mathbb{R}^{d_x}$ & support of observations\\
$\displaystyle \Z \subseteq \mathbb{R}^{d_z}$ & support of pretrained representations\\
$\displaystyle \C \subseteq \mathbb{R}^{d_c}$ & support of ground-truth concept vectors\\
$\displaystyle S \subseteq [d_c]$ & subset of varying concepts in a given pair $(\x,\tilde\x)$ \\
$\displaystyle V \subseteq [d_c]$ & subset of concepts allowed to vary between $\x$ and $\tilde \x$ \\
%$\displaystyle |S|$ & number of concepts varying within a pair\\
%$\displaystyle |V|$ & number of concepts varying across all pairs\\
%$\displaystyle [d_c] \backslash V$ & concepts in $[d_c]$ that are not in $V$\\
$\deltac$ & concept shift vector\\
$\hatdeltac$ & estimated concept shift vector  \\
$\deltaz$ & pretrained representation shift vector  \\
%$\deltac_V$ & concept shift vector with shifts in concepts in $V$  \\
% $\displaystyle \pi(k)$ & permutation of $k \in \sigma_{|V|}$, the set of all permutations of varying concept indices \\
$\displaystyle g : \C \rightarrow \X$ & map from concept representations to observations\\
$\displaystyle f : \X \rightarrow \Z$ & map from observations to learned representations \\
$\displaystyle r : \Z \rightarrow \C$ & encoding function \\
$\displaystyle \hat{r} : \C \rightarrow \Z$ & estimated encoding function \\
$\displaystyle q : \C \rightarrow \Z$ & decoding function\\
$\displaystyle \hat{q} : \C \rightarrow \Z$ & estimated decoding function \\
$\displaystyle \phi_{k, \lambda} : \Z \rightarrow \Z$ & steering function \\
$\displaystyle \hat{\phi}_{k, \lambda} : \Z \rightarrow \Z$ & estimated steering function \\
$\displaystyle \A$ & linear map between concept representations and learnt representations\\
%$\displaystyle \A_V$ & subset of the linear map between concept representations for concepts in $V$ and their learnt representations \\




\end{tabular}

% \subsubsection{Justification for injectivity of the encoding and decoding functions}
% \label{apx:inj}

% $\mathbf{r} : \Z \rightarrow \C$ maps $\delta \z \in \mathbb{R}^{d_Z}$
% vectors to $\s_{\delta C} \in \mathbb{R}^K$ vectors where $K \leq d_Z$. $q$ maps $\s_{\delta C}$ back to $\delta \z$. To establish the injectivity of these transformations, we need to ensure that $\br \coloneqq (\mathbf{T}_1, \mathbf{k}_1)$ and $\q \coloneqq (\mathbf{T}_2, \mathbf{k}_2)$ are such that $\mathbf{T}_2 = \mathbf{T}_1^T$ and $\mathbf{T}_2$ has orthonormal columns which means $\mathbf{T}_2\mathbf{T}_2^T = \mathbf{I}$.

% To analyse the injectivity of a linear transform, we use the Rank-Nullity theorem which states that for a transformation $\mathbf{T}  : V \rightarrow W, \text{dim}(V) = \text{rank}(\mathbf{T}) + \text{nullity}(\mathbf{T})$.

% $$\mathbf{T} \text{ is injective } \iff \text{ker}(\mathbf{T}) = \{0\} \text{ or } \text{nullity}(\mathbf{T}) = 0$$

% (1) $\text{dim}(V) = K$ and we consider $\mathbf{T}_2$ to have orthonormal columns, which means they are linearly independent and $\mathbf{T}_2$ is a full rank matrix $\implies  \text{nullity}(\mathbf{T}_2) = 0$ and $\mathbf{q}$ is injective.

% (2) The injectivity of $\mathbf{T}_1$ is less straightforward. We know that  $\mathbf{T}_1$ is full rank since  $\mathbf{T}_1 = \mathbf{T}_2^T$ and $\mathbf{T}_2^T$ is full rank. We also have $\mathbf{T}_2\mathbf{T}_2^T = \mathbf{I} \implies \mathbf{T}_2 \circ \mathbf{T}_1 = \mathbf{I}$. Consider two vectors $\delta \z$ and $\delta \z'$ such that $\mathbf{T}_1(\delta \z) = \mathbf{T}_1(\delta \z')$. Applying $\mathbf{T}_2$ to both sides, $\mathbf{T}_2(\mathbf{T}_1(\delta \z)) = \mathbf{T}_2(\mathbf{T}_1(\delta \z')) \implies \delta \z = \delta \z'$. Thus, $\mathbf{T}_1$ or $\br$ is injective.

% \subsection{Injectivity of $\A_V$}

% \label{apx:inj}

\subsection{Steering Functions}
\label{apx:steering}
From \cref{fig:steering}, for any concept $k$, the steering function $\phi_{k, \lambda}$ mirrors the transformations between concepts described as $\tilde \c_{k, \lambda} := \psi_{k, \lambda}(\c)$ in the learnt representation space through functions defined as:
\begin{definition} (\textbf{Steering function})
\label{defn:steering_fn}
    Fix a target concept $k$ and $\lambda \in \mathbb{R}$. %Then, for each concept representation $\c$, the perturbed concept representation $\tilde \c_{k, \lambda} = \c + \lambda \mathbf{e}_k$, where $\mathbf{e}_k$ is the $k$-th standard basis vector, corresponds to the counterfactual observation $\tilde \x_k = g(\tilde \c_k)$  in concept $k$.
    A \textbf{steering function} $\phi_{k,\lambda}: \mathcal{Z} \rightarrow \mathcal{Z}$ is a function such that for all $\c \in \C$, $\phi_{k, \lambda}(f(g(\c))) = f(g(\psi_{k, \lambda}(\c)))$.
\end{definition}

According to \cref{defn:steering_fn}, a steering function \footnote{Steering functions are not guaranteed to exist. However, if $f$ and $g$ are injective, we have $\phi_{\lambda, k}(\z) = f(g(g^{-1}(f^{-1}(\z)) + \lambda\mathbf{e}_k))$.} maps each representation $\z = f(\x) = f(g(\c)$ to its perturbed analog $\tilde\z_{\lambda,k} \coloneqq f(\tilde\x_{\lambda,k})$, where $\tilde{\mathbf{x}}_{k,\lambda} \coloneqq g(\tilde{\mathbf{c}}_{k,\lambda})$ is the corresponding perturbed observation. Thus, if the $k$-th concept is language, a steering function maps $\z = f(\x)$, the embedding of a sentence $\x$, to $\tilde \z_{k, \lambda} = f(\tilde \x_{k, \lambda})$, the embedding of the same sentence written in a 
different language. The form of the steering function depends on the form of the transformations $\psi_{k, \lambda}$ in concept space $\mathcal{C}$. We assume transformations $\psi_{k, \lambda}$ to be additive perturbations:

% \begin{definition}
%     (\textbf{Concept perturbation}): The perturbed concept vector is denoted as $\tilde{\mathbf{c}}_{k,\lambda} \coloneqq \mathbf{c} + \lambda \mathbf{e}_k$, , where $\mathbf{e}_k$ is the $k$-th standard basis vector. 
% \label{defn:affordance}
% \end{definition}

\begin{figure}
\centering
\includegraphics[scale=0.27]{sf.png}

\caption{A steering function $\phi_{k, \lambda}$ is s.t. the above diagram commutes, i.e., $\phi_{k, \lambda}(f(g(\c))) = f(g(\psi_{k, \lambda}(\c))) \forall \c$. (see \cref{defn:steering_fn}).}

% , i.e., $\phi_{\lambda, k}(f(g(\c))) = f(g(\c + \lambda \rve_k))$ for all $\c$.}
\label{fig:steering}
\end{figure}


 %  as illustrated in \cref{fig:steering}.
 
  To model the additive changes in $\c$, one can use an analogous additive perturbation map in $\z$ s.t. $\tilde \z_{k, \lambda} := \phi_{k, \lambda}(\z) $ can be written as $\tilde \z_{k, \lambda} := \z + \deltaz_{k, \lambda}$, where $\deltaz_{k, \lambda}$ might be an arbitrarily dense vector in $\Z$. 

In practice, a steering function $\phi_{\lambda,k}$ can be learned via supervised learning given a dataset comprising of carefully designed paired observations $(\x, \tilde \x_{k})$, in which a single concept changes between $\x$ and $\tilde \x_{k}$ \citep{shen2017styletransfernonparalleltext, turner2024steeringlanguagemodelsactivation,rimsky2024steering}. However, such a dataset might be difficult to acquire. This raises the following question, at the heart of our contribution:

\textit{How can we learn a steering function $\phi_{k, \lambda}$ with a dataset of paired observations $(\x, \tilde\x)$ in which multiple concepts vary?}

Thus, unsupervised approaches such as sparse autoencoders (SAEs) \citep{cunningham2023sparseautoencodershighlyinterpretable} are often employed towards steering distinct concepts. In this paper, we develop sufficient desiderata to show how identifiability leads to better steering performance.


\subsection{Linear Identifiability is insufficient for steering.}
\label{apx:linid_insuff}

In \cref{sec:wscrl} we showed how identifiability up to permutation and scaling leads to distinct steering vectors for individual concepts. Here, we show that the same strategy fails when concept shifts are only linearly identified, i.e., $\hat q \coloneqq \A_V \mathbf{L}$. In this case, we see that 
\begin{align*}
% \begin{split} \nonumber
    \hat q(\rve_k) = \A_V \rmL \rve_k = \sum_{j=1}^{|V|} \rmL_{j,k} \A \rve_j = \A\sum_{j=1}^{|V|} \rmL_{j,k} \rve_j \,,
% \end{split}
\end{align*}
which itself implies that 
\begin{align*}
    \z + \hat q(\mathbf{e}) = \A\c + \A\sum\nolimits_{j=1}^{|V|} \rmL_{j,k} \rve_j = \A(\c + \sum\nolimits_{j=1}^{|V|} \rmL_{j,k} \rve_j) = f(g(\c + \sum\nolimits_{j=1}^{|V|} \rmL_{j,k} \rve_j)) \,.
\end{align*}

That is, each learned steering vector $\hat{q}(\rve_k)$ can potentially change every concept in $V$. To recover the steering vectors, we need to learn $\rmL^{-1}$, which requires paired samples $(\tilde\z_{j, \lambda}, \z)$ that vary in a single concept for each concept $j$ \citep{rajendran2024learning}. This highlights the importance of enforcing sparsity, as it is the key element allowing us to go from $\hat{q} \coloneqq \A_V\rmL$ (\cref{prop:linear_ident}) to $\hat{q} \coloneqq \A_V\rmD\rmP$ (\cref{prop:perm_ident}). 

A potential advantage of linearly identifying steering vectors, however, is that learning the linear function $\rmL^{-1}$ may require fewer samples than learning a potentially nonlinear steering function (\cref{defn:steering_fn}) from counterfactual samples.


\subsection{Proof of \cref{prop:linear_ident} (linear identifiability)}
\label{app:linear_ident}

\linearIdent*


\begin{proof}
    We note that the solution $q^* \coloneqq \A_V$ and $r^* \coloneqq \A_V^+$ minimizes the loss since
    \begin{align}
        \mathbb{E}_{\x, \tilde\x}||\deltaz - q^*(r^*(\deltaz))||^2_2 &= \mathbb{E}_{\x, \tilde\x}||\deltaz - \A_V\A_V^{+}\deltaz||^2_2 \\
        &=\mathbb{E}_{\c, \tilde\c}||\A_V \deltac_V - \A_V(\A_V^{+}\A_V) \deltac_V||^2_2 \\
        &=\mathbb{E}_{\c, \tilde\c}||\A_V \deltac_V - \A_V \deltac_V||^2_2 \\
        &= 0 \,,
    \end{align}
    where we used the fact that $\A_V$ is injective and thus $\A_V^+\A_V = \mathbf{I}$. This means all optimal solutions must reach zero loss.

    Now consider an arbitrary minimizer $(\hat r, \hat q)$. Since it is a minimizer, it must reach zero loss, i.e. 
    \begin{align}
        &\mathbb{E}_{\x, \tilde\x}||\deltaz - \hat q(\hat r(\deltaz))||^2_2 = 0 \\
        &\mathbb{E}_{\c, \tilde\c}||\A_V \deltac_V - \hat q(\hat r(\A_V\deltac_V))||^2_2 = 0
    \end{align}    
    This means we must have 
    \begin{align}\label{eq:988r9e0w}
        \A_V \deltac_V = \hat q(\hat r(\A_V\deltac_V)),\ \text{almost everywhere w.r.t. $p(\deltac_V)$.}
    \end{align}
     Because all functions both on the left and the right hand side are continuous, the equality must hold on the support of $p(\deltac_V)$, which we denote by $\Delta^c_V$. Moreover, since $\hat r$ and $\hat q$ are linear, they can be represented as matrices, namely $\mathbf{R} \in \sR^{|V| \times d_z}$ and $\mathbf{Q} \in \sR^{d_z \times |V|}$. We can thus rewrite \cref{eq:988r9e0w} as 
    \begin{align}\label{eq:98438383}
        \A_V \deltac_V = \mathbf{Q}\mathbf{R}\A_V\deltac_V\,,
    \end{align}
    which holds for all $\deltac_V \in \Delta^c_V$. By \cref{ass:suff_var}, we know there exists a set of $|V|$ linearly independent vectors in $\Delta^c_V$. Construct a matrix $\mathbf{C} \in \sR^{|V|\times|V|}$ whose columns are these linearly independent vectors. Note that $\mathbf{C}$ is invertible, by construction. 
    
    Since this \cref{eq:98438383} holds for all $\deltac_V \in \Delta^c_V$, we can write
    \begin{align}
        \A_V \mathbf{C} = \mathbf{Q}\mathbf{R}\A_V\mathbf{C}\\
        \A_V = \mathbf{Q}\mathbf{R}\A_V \,,
    \end{align}
    where we right-multiplied by $\mathbf{C}^{-1}$ on both sides. Since $\A_V$ is injective (\cref{ass:injective_Av}), we must have that $\mathbf{R}\A_V$ is injective as well. But since $\mathbf{R}\A_V$ is a square matrix, injectivity implies invertibility. Let us define $\mathbf{L} \coloneqq (\mathbf{R}\A_V)^{-1}$. We thus have
    \begin{align}
        \A_V &= \mathbf{Q}\mathbf{L}^{-1} \\
        \hat q = \mathbf{Q} &= \A_V\mathbf{L} \,,
    \end{align}
    which proves the first part of the statement.

    Now, we show that, for all $\z \in \text{Im}(\A_V)$, $\mathbf{R}\z = \mathbf{L}\A_V^+\z$. Take some $\z \in \text{Im}(\A_V)$. Because this point is in the image of $\A_V$, there must exists a point $\c \in \sR^{|V|}$ such that $\z = \A_V\c$. Now we evaluate
    \begin{align}
        \hat r(\z) = \mathbf{R}\z &= \mathbf{R}\A_V\c  \\
        &= \mathbf{L}^{-1}\c \label{eq:kksdm2949sk}\\ 
        &= \mathbf{L}^{-1}\A_V^+\A_V\c \label{eq:jdjdjdjnajajm}\\
        &= \mathbf{L^{-1}}\A_V^+\z \,, 
    \end{align}
    where we used the fact $\mathbf{R}\A_V = \mathbf{L}^{-1}$ in \cref{eq:kksdm2949sk} and the fact that $\A_V^+\A_V = \mathbf{I}$ in \cref{eq:jdjdjdjnajajm}. This concludes the proof.
\end{proof}


\subsection{Proof of \cref{prop:perm_ident} (permutation identifiability)}
\label{app:proof}
The proof is heavily based on \citet{lachapelle2023synergies} and \citet{xu2024sparsityprinciplepartiallyobservable}.

\permIdent*

\begin{proof}
    Recall that, in the proof of \cref{prop:linear_ident}, we showed that the solution $q^* \coloneqq \A_V$ and $r^* \coloneqq \A_V^+$ yields zero reconstruction loss, i.e.,
    \begin{align}
        \mathbb{E}_{\x, \tilde\x}||\deltaz - q^*(r^*(\deltaz))||^2_2 &= 0\,.
    \end{align}
    It turns out, this solution also satisfies the constraint $\mathbb{E}||r(\deltaz)||_0 \leq \beta \coloneqq \mathbf{E}||\deltac_V||_0$ since
    \begin{align}
        \mathbb{E}||r^*(\deltaz)||_0 = \mathbb{E}||\A_V^+(\A_V\deltac_V)||_0 = \mathbb{E}||\deltac_V||_0 = \beta \,,  
    \end{align}
    where we used the fact that $\deltaz = \A_V\deltac_V$ and $\A_V^+\A_V = \mathbf{I}$, since $\A_V$ is injective. This means that all optimal solutions to the constrained problem of \cref{eqn:recon,eqn:sparse_constraint} with $\beta \coloneqq \mathbb{E}||\deltac_V||_0$ must reach zero reconstruction loss.

    Let $(\hat r, \hat q)$ be an arbitrary solution to the constrained problem. By the above argument, this solution must reach zero loss. Thus, by the exact same argument as in \cref{prop:linear_ident}, there must exist an invertible matrix $\mathbf{L} \in \sR^{|V|\times |V|}$ such that 
    \begin{align}
        \hat q \coloneqq \A_V \mathbf{L} \quad \text{and} \quad \hat r(\z) \coloneqq \mathbf{L}^{-1}\A_V^+\z,\ \text{for all}\ \z\in\text{Im}(\A_V) \,.
    \end{align}
    Since $\hat r$ is optimal it must satisfy the constraint, which we rewrite as
    \begin{align}
        \mathbb{E}||\hat r (\deltaz)||_0 &\leq \mathbb{E}||\deltac_V||_0 \nonumber\\
        \mathbb{E}||\hat r (\A_V\deltac_V)||_0 &\leq \mathbb{E}||\deltac_V||_0 \nonumber\\
        \mathbb{E}||\mathbf{L}^{-1}\A_V^+(\A_V\deltac_V)||_0 &\leq \mathbb{E}||\deltac_V||_0 \nonumber\\
        \mathbb{E}||\mathbf{L}^{-1}\deltac_V||_0 &\leq \mathbb{E}||\deltac_V||_0\,, \label{eqn:core}
    \end{align}
    where we used the fact that $\hat r$ restricted to the image of $\A_V$ is equal to $\mathbf{L}^{-1}\A_V^+$ when going from the second to the third line.
    
    At this stage, we can use the same argument as \citet{lachapelle2023synergies} to conclude that $\mathbf{L}$ is a permutation-scaling matrix. For completeness, we present that result into \cref{lemma:synergies_proof} and its proof below. One can directly apply this lemma, thanks to \cref{ass:suffsupp} and the fact that sets of the form $\{\deltac_S \in \sR^{|V|} \mid \mathbf{a}^\top\deltac_S = 0\}$ with $\mathbf{a} \not = 0$ are proper linear subspaces of $\sR^{|V|}$ and thus have zero Lebesgue measure, and thus 
    $$\mathbb{P}_{\deltac_S \mid S}\{\deltac_S \in \sR^{|V|} \mid \mathbf{a}^\top\deltac_S = 0\} = 0 \,.$$
    This concludes the proof.
\end{proof}

The proof of the following lemma is taken directly from \citet{lachapelle2023synergies} (modulo minor changes in notation). The original work used this argument inside a longer proof and did not encapsulate this result into a modular lemma. We thus believe it is useful to restate the result here as a lemma containing only the piece of the argument we need. We also include the proof of \citet{lachapelle2023synergies} for completeness. Note that \citet{xu2024sparsityprinciplepartiallyobservable} also reused this result to prove identifiability up to permutation and scaling.

\begin{lemma}[\citet{lachapelle2023synergies}]\label{lemma:synergies_proof}
Let $\mathbf{L} \in \sR^{m\times m}$ be an invertible matrix and let $\x$ be an $m$-dimensional random vector following some distribution $\mathbb{P}_\x$. Define the set $S \coloneqq \{j \in [m] \mid \x_j \not= 0\}$, which is random (because $\x$ is random) with probability mass function given by $p(S)$. Let $\mathcal{S} \coloneqq \{S \subseteq [m] \mid p(S) > 0\}$, i.e. it is the support of $p(S)$. Assume that
\begin{enumerate}
    \item For all $j \in [m]$, we have $\bigcup_{S \in \mathcal{S} | j \notin S} S = [m] \setminus \{j \}$; and
    \item For all $S \in \mathcal{S}$, the conditional distribution $\mathbb{P}_{\x_S \mid S}$ is such that, for all nonzero $\mathbf{a} \in \sR^{|S|}$, $\mathbb{P}_{\x_S \mid S}\{\x_S \mid \mathbf{a}^\top \x_S = 0\} = 0$.
\end{enumerate}
Under these assumptions, if $\mathbb{E}||\mathbf{L}\x||_0 \leq \mathbb{E}||\x||_0$, then $\mathbf{L}$ is a permutation-scaling matrix, i.e. there exists a diagonal matrix $\mathbf{D}$ and a permutation matrix $\mathbf{P}$ such that $\mathbf{L} = \mathbf{D}\mathbf{P}$ 
\end{lemma}
\begin{proof}
    We start by rewriting the l.h.s. of $\mathbb{E}||\mathbf{L}\x||_0 \leq \mathbb{E}||\x||_0$ as
\begin{align}
    \sE\normin{\x}_{0} &= \sE_{p(S)}\sE[\sum_{j=1}^m \mathbf{1}(\x_j \not= 0) \mid S] \\
    &= \sE_{p(S)}\sum_{j=1}^m \sE[\mathbf{1}(\x_j \not= 0) \mid S] \\
    &= \sE_{p(S)}\sum_{j=1}^m \sP_{\x \mid S}\{\x \in \sR^m \mid \x_{j} \not= 0\}\\
    &= \sE_{p(S)}\sum_{j=1}^m \mathbf{1}(j \in S)\,,
\end{align}
where the last step follows from the definition of $S$.

Moreover, we rewrite $\sE\normin{\mathbf{L}\x}_{0}$ as
\begin{align}
    \sE\normin{\mathbf{L}\x}_{0} &= \sE_{p(S)}\sE[\sum_{j=1}^m \mathbf{1}(\mathbf{L}_{j, :}\x \not= 0)\mid S]\\
    &= \sE_{p(S)}\sum_{j=1}^m \sE[\mathbf{1}(\mathbf{L}_{j, :}\x \not= 0)\mid S]\\
    &= \sE_{p(S)}\sum_{j=1}^m \sE[\mathbf{1}(\mathbf{L}_{j, S}\x_S \not= 0)\mid S]\\
    &= \sE_{p(S)}\sum_{j=1}^m \sP_{\x \mid S}\{ \x \in \sR^m \mid \mathbf{L}_{j, S}\x_S \not= 0\} \,.
\end{align}

Notice that
\begin{align}
    \sP_{\x \mid S}\{ \x \in \sR^m \mid \mathbf{L}_{j, S}\x_S \not= 0\} &= 1 - \sP_{\x \mid S}\{ \x \in \sR^m \mid \mathbf{L}_{j, S}\x_S = 0\}\,. %\\
    %&= 1 - \sP_{\mW \mid S}[\forall i \in [k],\ \mW_{i, S} \mL_{S, j} = 0]
\end{align}
Define $N_j$ be the support of $\mL_{j, :}$, i.e., $N_j \coloneqq \{i \in [m] \mid \mL_{j, i} \not= 0 \}$. 

When $S \cap N_j = \emptyset$, we have that $\mL_{S,j} = \bm0$ and thus
$$\sP_{\x \mid S}\{ \x \in \sR^m \mid \mathbf{L}_{j, S}\x_S = 0\} = 1 \,.$$ 
When $S \cap N_j \not= \emptyset$, we have that $\mL_{j,S} \not= \bm0$, and thus, by the second assumption, we have that 
$$\sP_{\x \mid S}\{ \x \in \sR^m \mid \mathbf{L}_{j, S}\x_S = 0\} = 0 \,.$$ 

Thus we can write
\begin{align}
    \sP_{\x \mid S}\{ \x \in \sR^m \mid \mathbf{L}_{j, S}\x_S \not= 0\} &= 1 - \sP_{\x \mid S}\{ \x \in \sR^m \mid \mathbf{L}_{j, S}\x_S = 0\}\\
    &= 1 - \mathbf{1}(S \cap N_j = \emptyset) \\
    &= \mathbf{1}(S \cap N_j \not= \emptyset)\,,
\end{align}
which allows us to write
\begin{align}
    \sE\normin{\mathbf{L}\x}_{0} &= \sE_{p(S)}\sum_{j=1}^m \mathbf{1}(S \cap N_j \not= \emptyset)\,.
\end{align}
The original inequality $\mathbb{E}||\mathbf{L}\x||_0 \leq \mathbb{E}||\x||_0$ can thus be rewritten as
\begin{align}
    \sE_{p(S)}\sum_{j=1}^m \mathbf{1}(S \cap N_j \not= \emptyset) &\leq \sE_{p(S)}\sum_{j=1}^m \mathbf{1}(j \in S) \,. \label{eq:before_perm}
\end{align}
Since $\mL$ is invertible, there exists a permutation $\sigma: [m] \rightarrow [m]$ such that, for all $j \in [m]$, $\mL_{j, \sigma(j)} \not=0$ (e.g. see Lemma B.1 from \citet{lachapelle2023synergies}). In other words, for all $j \in [m]$, $j \in N_{\sigma(j)}$. Of course we can permute the terms of the l.h.s. of~\cref{eq:before_perm}, which yields
\begin{align}
    \sE_{p(S)}\sum_{j=1}^m \mathbf{1}(S \cap N_{\sigma(j)} \not= \emptyset) &\leq \sE_{p(S)}\sum_{j=1}^m \mathbf{1}(j \in S)\\
    \sE_{p(S)}\sum_{j=1}^m \left(\mathbf{1}(S \cap N_{\sigma(j)} \not= \emptyset) - \mathbf{1}(j \in S)\right) &\leq 0 \,. \label{eq:sum_of_positive}
\end{align}
We notice that each term $\mathbf{1}(S \cap N_{\sigma(j)} \not= \emptyset) - \mathbf{1}(j \in S) \geq 0$ since whenever $j \in S$, we also have that $j \in S \cap N_{\sigma(j)}$ (recall $j \in N_{\sigma(j)}$). Thus, the l.h.s. of \cref{eq:sum_of_positive} is a sum of non-negative terms which is itself non-positive. This means that every term in the sum is zero:
\begin{align}
    \forall S \in \gS,\ \forall j \in [m],&\ \mathbf{1}(S \cap N_{\sigma(j)} \not= \emptyset) = \mathbf{1}(j \in S)\,.
\end{align}
Importantly,
\begin{align}
    \forall j \in [m],\ \forall S \in \gS,\ j \not\in S \implies S \cap N_{\sigma(j)} = \emptyset\,,
\end{align}
and since $S \cap N_{\sigma(j)} = \emptyset \iff N_{\sigma(j)} \subseteq S^c$ we have that
\begin{align}
    \forall j \in [m],\ \forall S \in \gS,\ j \not\in S \implies N_{\sigma(j)} \subseteq S^c \\
    \forall j \in [m],\ N_{\sigma(j)} \subseteq \bigcap_{S \in \gS \mid j \not\in S} S^c \,. \label{eq:last_intersect}
\end{align}
By assumption, we have $\bigcup_{S \in \gS \mid j \not\in S} S = [m] \setminus \{j\}$. By taking the complement on both sides and using De Morgan's law, we get $\bigcap_{S \in \gS \mid j \not\in S} S^c = \{j\}$, which implies that $N_{\sigma(j)} = \{j\}$ by~\Cref{eq:last_intersect}. Thus, $\mL = \mathbf{D}\mathbf{P}$ where $\mathbf{D}$ is an invertible diagonal matrix and $\mathbf{P}$ is a permutation matrix.
\end{proof}

\begin{comment}
\begin{mainbox}
\cref{prop:id} (Identifiability up to permutation and scaling.): Encoding $\hat{r}$ and decoding $\hat{q}$ satisfying the conditions in \cref{ass:injective_Av} and \cref{ass:recon}, defined as per the model in \cref{prob:reconstruction_sparse} where the sparsity regularization satisfies the constraint in \cref{eqn:sparsity} are identified as per \cref{defn:perm_idnt} for the range of datasets meeting the conditions in \cref{ass:suffsupp}.
\end{mainbox}
% in the proof from lemma b, v is indeed considered to be affine, so do we need the confusion being affine and linear? little note on using it interchangeably or say we write affine cuz it's more general
% should probably just write s_delta c and remove the delta z part from everywhere
% end: since model is linear, all params are identified to same ~

\textbf{Proof}:  Recall from \cref{ass:recon} that $\hat{q}$ and $\hat{r}$ satisfy the zero reconstruction error constraint:

$$\mathbb{E} ||\deltaz -\hat{q}(\hat{r}(\deltaz)) ||_2^2 = 0.$$

Rewriting this in terms of $\deltac$ since $\deltaz = q(\deltac_S)$ for any pair of observations:

$$ \implies \mathbb{E} || q(\deltac_S) - \hat{q}(\hat{r}(q(\deltac_S))||^2_2 = 0$$

This implies $q$ and $\hat{q} \circ \hat{r} \circ q$ are $P_{\deltac_S}$-almost equal in the $\text{supp}(P_{\deltac_S}) = |V|$, the set of varying concepts across all pairs of observations $(\c, \tilde \c) \sim p(\c, \tilde \c | S)$ where $S \sim P(S)$. 
\begin{flalign}
    \implies q \stackrel{P_{\deltac_S}}{\sim} \hat{q} \circ \hat{r} \circ q
    \label{eqn:p_eq}
\end{flalign}

By assumption, $q, \hat{q}, \hat{r}$ are all injective and affine, hence continuous. This implies $\hat{q} \circ \hat{r} \circ q$ must be continuous as well.

Since in \cref{eqn:p_eq} two continuous quantities are equal to each other in probability $P_{\deltac_S}$ where $\text{supp}(P_{\deltac_S}) =  [|V|]$,

\begin{flalign}
    \implies q = \hat{q} \circ \hat{r} \circ q \quad \forall \deltac_S \in [|V|]
    \label{eqn:q_eq}
\end{flalign}

Left multiplying the above equation by $\hat{q}^{-1}$, we get:

\begin{flalign}
    \hat{q}^{-1}(q (\deltac_S)) = \hat{r}(q (\deltac_S)) \quad \forall \deltac_S \in [|V|]
    \label{eqn:origins_v}
\end{flalign}

In \cref{eqn:origins_v}, on the RHS: we have a composition of injective affine functions $q$ and $\hat{r}$. Hence, the function on the LHS must also be an injective affine function. This means that $\hat{q}^{-1} \circ q$ is an injective affine function. Let's denote: 
\begin{flalign}
    \hat{q}^{-1} \circ q = m
    \label{eqn:v}
\end{flalign}

where $m: \C \rightarrow \C$ is an injective affine transformation. Recall the sparsity constraint from \cref{eqn:sparsity}:
$$\mathbb{E} || \hatdeltac_S ||_0 \leq \mathbb{E} || \deltac_S ||_0$$
Substituing $\hatdeltac_S = \hat{r}(\deltaz)$ and $\deltaz = q(\deltac_S)$:

$$ \implies \mathbb{E} || \hat{r}(q(\deltac_S)) ||_0 \leq \mathbb{E}|| \deltac_S||_0$$

which from \cref{eqn:origins_v} can be re-written as:

$$\mathbb{E}|| \hat{q}^{-1}(q(\deltac_S) ||_0 \leq \mathbb{E}|| \deltac_S||_0$$

Using \cref{eqn:v},

\begin{flalign}
    \mathbb{E}|| m(\deltac_S)||_0 \leq \mathbb{E}||\deltac_S||_0
    \label{eqn:q_ps}
\end{flalign}

From \cref{ass:suffsupp}, $\forall s \subseteq \mathcal{S}$, the probability measure $\mathbb{P}_{\deltac_S | S = s}$ has a density with respect to the Lebesgue measure. Thus leveraging Lemma B.3 in \citep{xu2024sparsityprinciplepartiallyobservable}, since $\hat q \circ q$ is injective and affine, Lemma B.3 implies that $m$ can only be a permutation composed with scaling $\forall \deltac_S \in [|V|]$. This means $\hat{q}^{-1} \circ q$ is a permutation composed with scaling.
$$\implies \hat q = m \circ q,$$

i.e. $\hat{q}$ identifies $q$ up to permutation and scaling. Since $q$ is linear and identified, all the parameters of $q$ are identified as well.
Further, rewriting the sparsity constraint from \cref{eqn:sparsity}, i.e.,  $\mathbb{E}|| \deltac_S ||_0 \leq \mathbb{E}|| \deltac_S||_0$, in terms of $\deltaz$ by substituting $\hatdeltac_S = \hat{r}(\deltaz)$:

$$ \implies \mathbb{E} || \hat{r}(\deltaz) ||_0 \leq \mathbb{E}|| r(\deltaz)||_0$$

Left multiplying the above equation by $r^{-1}$:

$$ \implies \mathbb{E} || r^{-1}(\hat{r}(\deltaz)) ||_0 \leq \mathbb{E}|| \deltaz||_0$$

Since $r$ and $\hat{r}$ 
 are injective affine transformations, the composition $r^{-1} \circ \hat{r}$ is an injective affine transformation as well. Denoting $r^{-1} \circ \hat{r} = n$ where $n : \Z \rightarrow \Z$ is an injective affine transformation, we get an equation analogous to \cref{eqn:q_ps}:

 $$\mathbb{E}|| n(\deltaz)||_0 \leq \mathbb{E}||\deltaz||_0$$

Which by Lemma B.3 in \citep{xu2024sparsityprinciplepartiallyobservable} implies $r^{-1} \circ \hat{r}$ is a permutation composed with scaling, or that:

$$\hat{r} = r \circ n,$$

i.e. $\hat{r}$ identifies $r$ up to permutation and scaling, and since $r$ is linear and identified, all the parameters of $r$ are identified as well.

From \cref{defn:perm_idnt}, $\exists$ permutation matrix $\rmP$ and diagonal matrix $\mathbf{D}$ s.t. $\hat r =  \A_V^+ \mathbf{D} \rmP$ (on the domain of the observed $\deltaz$) and $\hat q =  \A_V \mathbf{D} \rmP$.


\begin{subbox}
    What happens if we consider non-linear $q$ for extracting concept vectors? In the case of a non-linear $q$, we need further assumptions on the distribution of concepts that can be identified \citep{xu2024sparsityprinciplepartiallyobservable}, restricting the set of concepts that can be provably steered.
\end{subbox}
\end{comment}

\subsection{Distributions satisifying \cref{ass:suffsupp}}
\label{apx:suffsupp}

In $\mathbb{R}^{|S|}$, any lower-dimensional subspace has Lebesgue measure 0. By defining the probability measure of $\deltac_S | S$ with respect to the Lebesgue measure, its integral over any lower-dimensional subspace of $\mathbb{R}^{|s|}$ will be 0. Consider a few examples of $\mathbb{P}_{\deltac_S | S}$ directly taken from \citep{lachapelle2023synergies} with adapted notation just for illustration purposes.


\begin{figure}[ht]
\centering
    \includegraphics[width=0.3\linewidth]{pdelta.png}
    \caption{Three illustrative examples of $\mathbb{P}_{\deltac_S | S}$: Only distribution II satisfies \cref{ass:suffsupp}.}
    \label{fig:asm24}
\end{figure}

In \cref{fig:asm24}, distributions I and III do not satisfy \cref{ass:suffsupp} whereas distribution II does. This is because I represents the support of a Gaussian distribution with a low-rank covariance and III represents finite support; both of these distributions will be measure zero in $\mathbb{R}^{|S|}$. On the other hand, II represents level sets of a Gaussian distribution with full-rank covariance. Please refer to \citet{lachapelle2023synergies} for a comprehensive explanation.


\subsection{Interpreting the Linear Representation Hypothesis}

\label{apx:lrh}
\textbf{Assumption 1}[Linear representation hypothesis]
The generative process $g: \C \to \X$ and the learned encoding function $f: \X \to \Z$ are such that  $f \circ g : \C \to \Z$ is linear, implying there exists a $d_z \times d_c$ real matrix $\A$ such that:
    \begin{align}
        % \label{eq:linear_entangle}
        \z = f(g(\c)) = \A \c \ .
        %\deltaz = f(g(\deltac)) = \A \deltac,
    \end{align}
The linear representation hypothesis (LRH) implies that the learned representation $\z$ \emph{linearly encodes concepts}. 
% We might wonder whether the representations we learn in practice possess this property. 
A long line of work provides evidence for this hypothesis (c.f. \citet{rumelhart1973model, hinton1986learning, mikolov-etal-2013-linguistic, ravfogel2020null}). More recently, theoretical work justifies why linear properties could arise in these models (c.f. \citet{ jiang2024originslinearrepresentationslarge, roeder2021linear, marconato2024all}). \Cref{sec:rel} provides a full list of related work, while \cref{apx:lrh} provides an explanation of the equivalence between LRH's different interpretations. \citet{rajendran2024learning} also leverage the LRH in their work.
\begin{corollary}
    If concept changes act on latent embeddings following $\tilde \z = \z + \deltaz$ and $q$ and $r$ are injective, they must be affine transformations.
\end{corollary}

\textbf{Proof}:  Starting with the interpretation of the \textit{linear representation hypothesis} such that $\thicktilde{\z} = \z + \deltaz$ where $\z = q (\c)$ and $\thicktilde{\z} = q(\thicktilde{\mathbf{c}})$:

$$\implies q(\thicktilde{\c}) = q(\c) + \deltaz$$

Since we identify only the varying concepts, this corresponds to identifying a subspace of the original concept space in which $\tilde \c = \c + \deltac_V$.

Using the injectivity of $q$ (\cref{ass:injective_Av}):
\begin{flalign}
    q(\mathbf{c} + \deltac_V) = q(\mathbf{c}) + \deltaz
    \label{eqn:prejacob}
\end{flalign}

Taking the gradient of both the LHS and the RHS wrt $\mathbf{c}$,

$$\frac{\partial(\mathbf{c} + \deltac_V)}{\partial(\mathbf{c})} \nabla_{(\mathbf{c} + \deltac_V)} q(\mathbf{c} + \deltac_V) = \nabla_{\mathbf{c}} q(\mathbf{c})$$
$$\nabla_{(\mathbf{c} + \deltac_V)} q(\mathbf{c} + \deltac_V) = \nabla_{\mathbf{c}} q(\mathbf{c})$$
\begin{flalign}
    \mathbf{J}^T(\mathbf{c} + \deltac_V) = \mathbf{J}^T(\mathbf{c})
\end{flalign}
Where $\mathbf{J}(\mathbf{c})$ is the Jacobian of $q$ at $\mathbf{c}$ and $\mathbf{J}(\mathbf{c} + \deltac_V)$ is the Jacobian of $q$ at $\mathbf{c} + \deltac_V$.

 $$\begin{bmatrix}
\nabla q_1(\mathbf{c} + \deltac_V)\\
\nabla q_2(\mathbf{c} + \deltac_V)\\
\nabla q_3(\mathbf{c} + \deltac_V)\\\
\cdot \\
\cdot \\
\nabla q_{d_Z}(\mathbf{c} + \deltac_V)
\end{bmatrix} - \begin{bmatrix}
\nabla q_1(\mathbf{c})\\
\nabla q_2(\mathbf{c})\\
\nabla q_3(\mathbf{c})\\\
\cdot \\
\cdot \\
\nabla q_{d_Z}(\mathbf{c})
\end{bmatrix} = 0$$

considering the $j$\textsuperscript{th} component of the difference,

$$\begin{bmatrix}
\nabla^2 q_j(\theta_1)\\
\nabla^2 q_j(\theta_2)\\
\nabla^2 q_j(\theta_2)\\\
\cdot \\
\cdot \\
\nabla^2 q_j(\theta_d)
\end{bmatrix}(\deltac_V) = 0$$

Following the proof in \citep{ahuja2022weakly},$\nabla^2q_j(\mathbf{c}) = 0$, which implies $q(\mathbf{c}) = \mathbf{A}_V\mathbf{c + b}$ where $\mathbf{A}_V \in \mathbb{R}^{d_Z \times d_Z}, \mathbf{b} \in \mathbb{R}^{d_Z}$ or that $q$ is affine. Similarly, we can show that $r$ is affine too by starting with $r(\z + \deltaz) = r(\z) + \deltac_V$.

\begin{corollary}
    If we assume $\thicktilde{\z} = \boldsymbol{\phi}(\z)$, for an affine map $q$, $\mathbf{A = I}$.
\end{corollary}

\textbf{Proof}: Let's assume the affine form of $q$ can be expressed as: 

\begin{flalign}
    \z = \mathbf{A}_V\mathbf{c} + \mathbf{b}
    \label{eqn:q_aff}
\end{flalign}
where $\mathbf{A}_V \in \mathbb{R}^{d_Z \times d_Z}$ and $\mathbf{k} \in R^{d_Z}$.

Similarly, $\thicktilde \z = q(\thicktilde{\mathbf{c}}) = \mathbf{A}_V\thicktilde{\mathbf{c}} + \mathbf{b}$ and we know $\tilde \c = \c + \deltac_V$. 
% can't start with injectivity here since we don't know if a generic delta_z maps onto the space of varying concepts
$$\implies \thicktilde \z = \mathbf{A}_V(\c + \deltac_V) + \mathbf{b}$$

we have $\thicktilde{\mathbf{z}} = \boldsymbol{\phi}(\mathbf{z})$  and from \cref{eqn:q_aff}:
\begin{flalign}  
    \boldsymbol{\phi}(\mathbf{A}_V\mathbf{c} + \mathbf{b}) = \mathbf{A}_V(\c + \deltac_V) + \mathbf{b}
\end{flalign}

In the above equation, we can see that the maximum degree of $\c$ on the RHS is $1$, which implies that the degree of $\c$ on the LHS should also at most be $1$, which implies $\boldsymbol{\phi}$ can at most be an affine function.

So let's assume $\boldsymbol{\phi}$ is an affine function of the form:

\begin{flalign}
    \thicktilde{\mathbf{z}} = \boldsymbol{\phi}(\mathbf{z}) = \mathbf{T}\mathbf{z} + \deltaz
\end{flalign}

where $\mathbf{T} \in \mathbb{R}^{d_Z \times d_Z}$ and $
\deltaz \in \mathbb{R}^{d_Z}$. Substituting this in the above equation, we get:

\begin{flalign}
    \mathbf{T}(\A_V \c + \mathbf{b}) + \deltaz = \mathbf{A}_V(\c + \deltac_V) + \mathbf{b}
\end{flalign}
$$\mathbf{Q(T - I)c} + (\mathbf{T - I)b} + (\deltaz \mathbf{ - Q} \deltac_V) = 0$$

For a non-trivial solution:
\begin{flalign}
    \mathbf{T = I} \\
    \deltaz = \mathbf{A}_V \deltac_V
\end{flalign}

So, we have proved that if we assume $q$ to be affine, then $\thicktilde{\mathbf{z}} = \mathbf{z} + \delta \mathbf{z}$.


% \textbf{Proof}: Let's start with the \cref{eqn:cdel} $\tilde \c = \c + \bsigma \c$, and $\c = \mathbf{r} (\mathbf{z})$:

% $$\implies \mathbf{r} (\thicktilde{\mathbf{z}}) = \mathbf{r} (\mathbf{z}) + \delta \c$$

% Using the injectivity of $\mathbf{r}$ (\cref{apx:inj}) and \cref{eqn:zc}:
% \begin{flalign}
%     \mathbf{r}(\mathbf{z} + \delta \mathbf{z}) = \mathbf{r}(\mathbf{z}) + \bsigma \c
%     \label{eqn:prejacob}
% \end{flalign}

% Taking the gradient of both the LHS and the RHS wrt $\mathbf{z}$,

% $$\frac{d(\mathbf{z} + \delta \mathbf{z})}{d(\mathbf{z})} \nabla_{(\mathbf{z} + \delta \mathbf{z})} \mathbf{r}(\mathbf{z} + \delta \mathbf{z}) = \nabla_{\mathbf{z}} \mathbf{r}(\mathbf{z})$$
% $$\nabla_{(\mathbf{z} + \delta \mathbf{z})} \mathbf{r}(\mathbf{z} + \delta \mathbf{z}) = \nabla_{\mathbf{z}} \mathbf{r}(\mathbf{z})$$
% \begin{flalign}
%     \mathbf{J}^T(\mathbf{z} + \delta \mathbf{z}) = \mathbf{J}^T(\mathbf{z})
% \end{flalign}

% Where $\mathbf{J}(\mathbf{z})$ is the Jacobian of $\mathbf{r}$ at $\mathbf{z}$ and $\mathbf{J}(\mathbf{z} + \delta \mathbf{z})$ is the Jacobian of $\mathbf{r}$ at $\mathbf{z} + \delta \mathbf{z}$.

%  $$\begin{bmatrix}
% \nabla \mathbf{r}_1(\mathbf{z} + \delta \mathbf{z})\\
% \nabla \mathbf{r}_2(\mathbf{z} + \delta \mathbf{z})\\
% \nabla \mathbf{r}_3(\mathbf{z} + \delta \mathbf{z})\\\
% \cdot \\
% \cdot \\
% \nabla \mathbf{r}_{Kd_Z}(\mathbf{z} + \delta \mathbf{z})
% \end{bmatrix} - \begin{bmatrix}
% \nabla \mathbf{r}_1(\mathbf{z})\\
% \nabla \mathbf{r}_2(\mathbf{z})\\
% \nabla \mathbf{r}_3(\mathbf{z})\\\
% \cdot \\
% \cdot \\
% \nabla \mathbf{r}_d(\mathbf{z})
% \end{bmatrix} = 0$$

% considering the $j$\textsuperscript{th} component of the difference,

% $$\begin{bmatrix}
% \nabla^2 \mathbf{r}_j(\theta_1)\\
% \nabla^2 \mathbf{r}_j(\theta_2)\\
% \nabla^2 \mathbf{r}_j(\theta_2)\\\
% \cdot \\
% \cdot \\
% \nabla^2 \mathbf{r}_j(\theta_d)
% \end{bmatrix}(\delta \mathbf{z}) = 0$$

% Following the proof in \citep{ahuja2022weakly},$\nabla^2\mathbf{r}_j(\mathbf{z}) = 0$, which implies $\mathbf{r}(\mathbf{z}) = \mathbf{Tz + k}$ or that $\mathbf{r}$ is affine. 

% % address the central nature of lrh: 

% \begin{corollary}
%     If we misspecify \cref{eqn:zc} such that instead of $\thicktilde{\z} = \z + \delta \z$, $\thicktilde{\z} = \boldsymbol{\phi}(\z)$, for an affine map $\q$, $\mathbf{A = I}$.
% \end{corollary}

% \textbf{Proof}: Let's assume the affine form of $\mathbf{q}$ can be expressed as: 

% \begin{flalign}
%     \z = \mathbf{q}(\mathbf{c}) = \mathbf{Q}\mathbf{c} + \mathbf{k}
%     \label{eqn:q_aff}
% \end{flalign}
% where $\mathbf{Q} \in \mathcal{M}_{d_Z, d_Z} (\mathbb{R})$ and $\mathbf{k} \in R^{d_Z}$.

% Similarly, $\thicktilde \z = \mathbf{q}(\thicktilde{\mathbf{c}}) = \mathbf{Q}\thicktilde{\mathbf{c}} + \mathbf{k}$ and from \cref{eqn:cdel}, we know $\tilde \c = \c + \bsigma_{\delta C}$. 

% $$\implies \thicktilde \z = \mathbf{Q}(\c + \bsigma_{\delta C}) + \mathbf{k}$$

% we have $\thicktilde{\mathbf{z}} = \boldsymbol{\phi}(\mathbf{z})$  and from \cref{eqn:q_aff}:
% \begin{flalign}  
%     \boldsymbol{\phi}(\mathbf{Q}\mathbf{c} + \mathbf{k}) = \mathbf{Q}(\c + \bsigma_{\delta C}) + \mathbf{k}
% \end{flalign}

% In the above equation, we can see that the maximum degree of $\c$ on the RHS is $1$, which implies that the degree of $\c$ on the LHS should also at most be $1$, which implies $\boldsymbol{\phi}$ can at most be an affine function.

% So let's assume $\boldsymbol{\phi}$ is an affine function of the form:

% \begin{flalign}
%     \thicktilde{\mathbf{z}} = \boldsymbol{\phi}(\mathbf{z}) = \mathbf{A}\mathbf{z} + \delta \mathbf{z}
% \end{flalign}

% where $\mathbf{A} \in \mathcal{M}_{d_Z, d_Z}(\mathbb{R})$ and $
% \delta \mathbf{z} \in \mathbb{R}^{d_Z}$. Substituting this in the above equation, we get:

% \begin{flalign}
%     \mathbf{A}(\mathbf{Qc} + \mathbf{k}) + \delta \mathbf{z} = \mathbf{Q}(\c + \bsigma_{\delta C}) + \mathbf{k}
% \end{flalign}
% $$\mathbf{Q(A - I)c} + (\mathbf{A - I)k} + (\delta \z \mathbf{ - Q} \bsigma_{\delta C}) = 0$$

% For a non-trivial solution:
% \begin{flalign}
%     \mathbf{A = I} \\
%     \delta \z = \mathbf{Q} \bsigma_{\delta C}
% \end{flalign}

% So, we have proved that if we assume $\mathbf{q}$ to be affine, then $\thicktilde{\mathbf{z}} = \mathbf{z} + \delta \mathbf{z}$.

    \textbf{Implications}: Multiple expositions \citep{templeton2024scaling} remark that it it not clear what the meaning of \textit{linear} exactly  is in the linear representation hypothesis. Informally, many results cited in support of the linear representation hypothesis either extract information with a linear probe, or add a vector to influence model behavior. Here, we assume that if linear meant concepts are linearly encoded in the latent space, we can show that this would correspond to shifts in the latent space representing net concept changes and vice versa, which means both interpretations are the same, so it does not matter which one is assumed. 

\subsection{Parallels with Sparse Coding and Addressing the Amortisation Gap}
\label{apx:dict-vs-ssae}

A long line of work on sparse coding and dictionary learning provides conditions under which a ground-truth dictionary $D$ is identifiable from observations, typically through the recovery of instance-specific sparse codes. Foundational contributions establish that identifiability follows from \emph{geometric constraints} on the dictionary $D$, most prominently the \emph{restricted isometry property} (RIP) \citep{candes2005decodinglinearprogramming, candes2005stablesignalrecoveryincomplete, baraniuk2008simple, foucart2013invitation} and various notions of \emph{incoherence} \citep{1614066, donoho2003optimally, gribonval2004sparse, 1337101}. Intuitively, these assumptions ensure that any sufficiently sparse combination of dictionary atoms is well-conditioned, which in turn guarantees the uniqueness of the sparse representation. These geometric constraints rule out correlated dictionary columns: under incoherence or RIP, the dictionary behaves approximately like an orthonormal basis when restricted to sparse supports.

Subsequent work derives identifiability guarantees for overcomplete dictionaries \citep{spielman2012exactrecoverysparselyuseddictionaries, agarwal2014learningsparselyusedovercomplete, gribonval2015sample, schnass2015localidentificationovercompletedictionaries}. These analyses again rely on variants of independence-of-support conditions, minimum separation between dictionary atoms, or distributional sparsity assumptions. Although theoretically powerful, these conditions are poorly aligned with the empirical structure of concept vectors in LLMs, where underlying factors of variation are often highly correlated. For instance, concepts such as truthfulness and harmlessness in alignment studies, or the deliberately correlated factors in our \textsc{corr}(2,1) benchmark, violate the incoherence and RIP assumptions that classical dictionary-learning theorems require. This mismatch partly explains the brittleness of standard sparse autoencoders in mechanistic interpretability settings where concept-level correlations are intrinsic rather than pathological.

Sparse shift autoencoders (SSAEs) offer a complementary route to identifiability that circumvents these geometric constraints. Instead of imposing global structure on $D$, SSAEs exploit assumptions about the \emph{data-generating process}, which we have shown to cover a wide spectrum of observed samples, such as even pairing any two text snippets by uniformly sampling from existing datasets. The paired samples do not need to vary in a single concept, or be subject to any rejection sampling. Consequently, SSAEs can recover correlated concept directions in regimes where dictionary-learning guarantees do not apply.

A further conceptual distinction arises from \emph{amortisation}. Classical sparse coding performs instance-wise inference by solving, for each input~$\boldsymbol{x}$,
\begin{equation}
\label{eq:sc}
\boldsymbol{z}^\star(\boldsymbol{x})
= \arg\min_{\boldsymbol{z}} 
\left\{
\frac12 \lVert \boldsymbol{x} - D \boldsymbol{z} \rVert_2^2
+ \lambda \lVert \boldsymbol{z} \rVert_1
\right\},
\end{equation}
and the identifiability guarantees focus on conditions under which $D$ and the sparse codes are uniquely recoverable. In contrast, sparse autoencoders learn an \emph{amortised inference map} $r_\theta(\boldsymbol{x})$ such that $r_\theta(\boldsymbol{x}) \approx \boldsymbol{z}^\star(\boldsymbol{x})$ for all~$\boldsymbol{x}$ in the data distribution. 

Recent analysis by \citet{oneill2025computeoptimalinferenceprovable} demonstrates that amortisation introduces systematic inductive biases absent from classical sparse coding. Their results show that the simple linear-nonlinear architecture of standard SAE encoders cannot implement the optimal sparse inference operator, even in regimes where the underlying dictionary is fully recoverable and exact inference is theoretically tractable. From a compressed-sensing perspective, this architectural bottleneck gives rise to a \emph{provable amortisation gap}: amortised encoders realise only a restricted family of piecewise-linear thresholding rules whose geometry is determined jointly by network structure and activation statistics, rather than by the true sparse-coding objective. Consequently, amortised inference may appear robust when $D$ violates incoherence or RIP, yet still entangle correlated factors that classical instance-wise optimisation would separate. This insight ties directly to observed superposition phenomena in SAEs and clarifies why amortised encoders deviate from the solution map~$\boldsymbol{x}\mapsto\boldsymbol{z}^\star(\boldsymbol{x})$ predicted by sparse-coding theory. Furthermore, by decoupling encoding and decoding, \citet{oneill2025computeoptimalinferenceprovable} show that modestly more expressive inference architectures significantly improve sparse-code recovery and yield more faithful features in LLM activations. These findings reinforce that amortisation is not a neutral approximation step, but a central determinant of the representational and identifiability properties of SAEs.

In the SSAE framework, amortisation becomes effective precisely because multi-concepts shifts induce changes of different levels of sparsity in the latent codes, such that amortising over them enables recovery of single concept shifts. Rather than requiring RIP or incoherence conditions on $D$, the encoder learns to approximate the update map induced by these shifts, enabling it to recover individual concept directions even when they are correlated. As a result, the identifiability guarantees arise from assumptions on the sparsity of latent concept shifts rather than from geometric properties of the dictionary.


% \textbf{This implies that $\delta C$ is identified by $\hat{\mathbf{s}}^{-1}(\delta \mathbf{z})$ up to permutation and elementwise linear transformations.}

% if in general, q/r is non-linear: then we don't have id results. so what are people with sae setups doing? they are using a non-linear r cuz that is the same mapping as ours from C to Z and they are using non-paired observations for this---dig a bit into this since in seb's paper they only show a counter-example to say additional assumptions are needed but can do ~ with one more assumption or:
% can we prove this for independent latents? would this mean that this way they can only extract independent concepts? 


\section{Implementation and experimental details}
\label{apx:imp}

Two key aspects of enforcing sparsity of the learnt representation are: (i) using hard constraints rather than penalty tuning, which helps address concerns with $\ell_1$-based regularization (e.g., feature suppression \citep{anders_etal_2024_composedtoymodels_2d}) and (ii) appropriate normalisation. For the former, we use the \texttt{cooper} library \citep{gallegoPosada2022cooper}. For the latter, we implement layer normalization \citep{ba2016layer} after the encoder and column normalization in the decoder at each step \citep{bricken2023monosemanticity, gao2024scaling}. To tune the model's hyperparameters in an unsupervised way, we use the Unsupervised Diversity Ranking (UDR) score \citep{duan2019unsupervised}, and test the model's sensitivity on key parameters (such as the sparsity level 
$\beta$ and learning rate).

\subsection{\isae Architecture}

The encoding $r: \Z \rightarrow \C$ and decoding functions $q: \C \rightarrow \Z$ constituting the \isae autoencoding framework are parameterized as follows:

\begin{flalign}
% \begin{split} \nonumber
     \hatdeltac_V \coloneqq {r}(\deltaz) & \coloneqq \mathbf{W}_e (\deltaz - \mathbf{b}_d) + \mathbf{b}_e;
     \label{eqn:enc}\\
    \hatdeltaz \coloneqq {q}(\hatdeltac_V) & \coloneqq \mathbf{W}_d \hatdeltac_V + \mathbf{b}_d\,.  \label{eqn:dec}
% \end{split}
\end{flalign}

\textbf{Parameters}.  $\mathbf{W}_e \in \mathbb{R}^{|V| \times d_z}, \mathbf{b}_e \in \mathbb{R}^{|V|}, \mathbf{W}_d \in \mathbb{R}^{d_z \times |V|}$, and $\mathbf{b}_d \in \mathbb{R}^{d_z}$ denote the encoder weights, encoder bias, decoding weights, and decoder bias respectively. The decoder bias is also treated as a pre-encoder bias purely for empirical performance improvement reasons based on ongoing discourse on engineering improvements in SAEs \citep{bricken2023monosemanticity, gao2024scaling}. The encoder and decoder weights are initialised s.t. $\mathbf{W}_d = \mathbf{W}_e^T$. The bias terms $\mathbf{b}_e$ and $\mathbf{b}_d$ are initialised to be all zero vectors. Further, after every iteration, the columns of $\mathbf{W}_d$ are unit normalised following \citet{bricken2023monosemanticity, gao2024scaling}.

\textbf{Data}. Data is layer-normalised analogous to \citet{gao2024scaling} prior to being passed as input to the encoder in batch sizes of $32$.

\textbf{Optimization.} Specifically, the following objective is optimized:
\begin{flalign}
    \min \frac{1}{N} \sum_{i = 1}^N \frac{||\deltaz_{(i)} - q(r(\deltaz_{(i)}))||^2_2}{||\deltaz_{(i)}||^2_2}, \\
    \text{s.t.} \frac{1}{|V|N} \sum_{i = 1}^N || r(\deltaz_{(i)})||_1 \leq \beta
\end{flalign}

We optimize the above constrained minimisation problem by computing its Lagrangian and the primal and dual gradients using the \texttt{cooper} library \citep{gallegoPosada2022cooper}. We use ExtraAdam \citep{gidel2020variationalinequalityperspectivegenerative} as both the primal and the dual optimizer, with the values of the primal and dual learning rates fixed throughout training and selected based on UDR scores (see \cref{apx:udr}). ExtraAdam uses \textit{extrapolation from the past} to provide similar convergence properties as extra-gradient optimizers \citep{korpelevich1976extragradient} without requiring twice as many gradient computations per parameter update or auxiliary storage of trainable parameters \citep{gidel2020variationalinequalityperspectivegenerative, gallegoPosada2022cooper}. Further, to account for the unit-norm adjustment of the columns of the decoder weights $\mathbf{W}_d$, we adjust gradients to remove discrepancies between the true gradients and the ones used by the optimizer. This done by removing any gradient information parallel to the columns of $\mathbf{W}_d$ at every step after the normalisation of the columns of $\mathbf{W}_d$. 

\textbf{Compute.} All experiments were conducted on the A100 GPUs (average time of 5min to 45 mins depending on the dataset). 

\subsubsection{Model Selection via Unsupervised Diversity Ranking (UDR)}
\label{apx:udr}

\begin{figure}
\centering
\includegraphics[scale=0.2]{udr_mcc_bin2.png}
\caption{UDR scores suggest a \texttt{primal\_lr} value of $0.005$ and a $\beta$ value of 0.1.}
\label{fig:udr_bin1}
\end{figure}
\begin{figure}    
\centering
\includegraphics[scale=0.18]{udr_mcc_bin1.png}
    % \centering
    % \includegraphics[width=0.5\linewidth]{figs/udr_bin2.png}
    % \caption{An example of observed UDR scores and MCC values for selecting optimal values of the two most important hyperparameters in tuning.}
    % \label{fig:udr}
    \caption{UDR scores suggest a \texttt{primal\_lr} value of $0.005$ and a $\beta$ value of $0.2$.}
    \label{fig:udr_bin2}
\end{figure}

Unsupervised model selection remains a notoriously difficult problem since there appears to be no unsupervised way of distinguishing between bad and good random seeds; unsupervised model selection should not depend on ground truth labels since these might biased the results based on supervised metrics. Moreover, in disentanglement settings, hyperparameter selection cannot rely solely on choosing the best validation-set performance. This is because there is typically a trade-off between the quality of fit and the degree of disentanglement (\citep{locatello2019challengingcommonassumptionsunsupervised}, Sec 5.4). For the proposed method in \cref{sec:wscrl}, identifiability of the decoder and of the learnt representation is essential to recover steering vectors for individual concepts. It is possible that a decoder with higher reconstruction error is identified to a greater degree. Hence, it is not sufficient to engineer a good unsupervised model solely based on how well it minimizes the reconstruction loss. \citet{duan2019unsupervised} propose the Unsupervised Disentanglement Ranking (UDR) score \citep{duan2019unsupervised}, which measures the consistency of the model across different initial weight configurations (seeds), which we use to fit our model. It is calculated as follows: for every hyperparameter setting, we compute MCCs between pairs of different runs and compute the median of all pairwise MCCs as the UDR score. We report the UDR scores and the mean pair-wise MCCs for the two most important hyperparameters affecting observed reconstruction error and MCC values---the learning rate of the primal optimizer (\texttt{primal\_lr}) and the sparsity level ($\beta$)---over $10$ pairs of $5$ random seeds in \cref{fig:udr_bin1} for the dataset, \textsc{lang}$(1, 1)$, and in \cref{fig:udr_bin2} for \textsc{binary}$(2, 2)$, over a selected hyperparameter range corresponding to decent reconstruction error. At slightly different hyperparameter settings, reconstruction error may spike even if the MCC remains acceptable. Such scenarios often fall outside the scope of consideration here, as they break the assumption of near-perfect reconstruction. While models may not achieve zero reconstruction loss in practice, we still expect it to remain reasonably low. As can be seen in \cref{fig:udr_bin1} and \cref{fig:udr_bin2}, MCC values typically correlate with the UDR scores.  Note that: \cref{fig:udr_bin1} and \cref{fig:udr_bin2} show UDR scores for only two datasets, but the same strategy (without plotting) was employed to select optimal hyperparameters for all datasets. Further, using these different models, we perform a sensitivity analysis on the two most important hyperparameters of our model---the sparsity level $\epsilon$ and the learning rate, which we report in \cref{apx:sens}.

% \subsubsection{Model details}
% \label{apx:model}
%  Since the size of the set $S$ is not known for each pair of observations, we initially set $\hatdeltac_S \in \mathbb{R}^{|V|}$. In \cref{sec:exp_llms}, we demonstrate recovery of steering vectors for $\hatdeltac_S \in \mathbb{R}^k$ where $|V| \le k \le d_Z$, satisfying \cref{ass:injective_Av}.

% We know that $\hatdeltac_V \in \mathbb{R}^{|V|}$. However, we do not know $S$ for every input $\deltaz$. So, we can assume without loss of generality that the predicted $\hatdeltac_S \in \mathbb{R}^{|V|}$. We will show in \cref{sec:exp_llms} that even if we assume $\hatdeltac_S \in \mathbb{R}^{k}$ where $|V| \leq k$ (and $k \leq d_z$ to ensure injectivity of $q$ and $r$ (\cref{ass:injective_Av})), we can estimate $\hatdeltac_S$ fairly accurately. Thus, in the general case we can say: $\hatdeltac_S \in \mathbb{R}^{k}$, $\mathbf{W}_e \in \mathcal{M}_{k, d_Z}(\mathbb{R})$, and $\mathbf{W}_d \in \mathcal{M}_{d_Z, k}(\mathbb{R})$. Here, the concept vector $\hatdeltac_S$ denotes the concepts that changed for the specific input $\deltaz$. This estimation method takes inspiration from sparse autoencoders, which are applied to find interpretable features in LLM representations.

% We implement the single layer linear auto-encoding framework presented in \cref{sec:model} with BatchNorm \citep{ioffe2015batchnormalizationacceleratingdeep} at the output of the encoder. We impose \cref{eqn:sparse_constraint_l1} as a hard constraint for the optimisation objective in \cref{eqn:model_recon} using Cooper \citep{gallegoPosada2022cooper}, and select a sparsity level for the encoder's output. This is a a major departure from recent work using SAEs to discover concepts in the latent space of LLMs. We argue that this addresses concerns associated with $l1$ regularisation, such as feature suppression \citep{anders_etal_2024_composedtoymodels_2d}. Further, we normalise the columns of the decoder after every step following \citet{bricken2023monosemanticity, gao2024scaling}. For more details, please refer to \cref{apx:imp}. To perform unsupervised model selection, we track the Unsupervised Disentanglement Ranking (UDR) scores \citep{duan2019unsupervised}, which measure the consistency of the model across different initial weight configurations (seeds). Further, we perform a sensititvity analysis on the two most important hyperparameters of our model---the sparsity level $\epsilon$ and the learning rate (\cref{apx:sens}). 
% , which mitigates issues such as feature suppression \citep{anders_etal_2024_composedtoymodels_2d}. This approach yields improved identifiability, even under feature absorption (or strong correlations). 
% \sj{better id implies these}
% Model selection is tricky since good disentanglement usually comes at the cost of a good fit. We report the UDR score as the median of the pairwise MCCs between different runs for a few different hyperparameter settings \citep{lachapelle2022disentanglement}.


\subsubsection{Sparse optimization}
\label{app:sparseopt}
% \sj{todo: add details about comparison with TopK}
We choose to enforce sparsity in the learning objective of the model as an explicit constraint rather than as $l_1$-regularisation due to the benefits listed in \cref{tab:sparsityreg}. In areas such as compressive sensing, signal processing, and certain machine learning applications, constrained optimization approaches have shown superior performance in recovering sparse signals and providing better generalization performance.

\begin{table}[ht]
\centering
\begin{tabularx}{\textwidth}{|X|X|X|}
   \hline 
   & \textbf{Constrained optimization} & \textbf{$\ell_1$-regularisation} \\
   \hline
   \textit{Optimization efficiency} & Finding the optimal solution and enforcing sparsity are separate tasks. Methods like augmented Lagrangian formulations iteratively enforce sparsity while optimizing the objective function, which can lead to more stable convergence. & The $l_1$ penalty introduces a non-differentiable point at zero, which requires careful tuning and can be sensitive to initialization and hyperparameters. \\
\hline
\textit{Hyperparameter tuning} & The primary hyperparameter is the sparsity level $\epsilon$, which can be set based on domain knowledge or practical constraints, simplifying the model selection process. & The primary hyperparameter is the strength of the sparsity penalty in the training objctive $\lambda$, which needs tuning to prevent under or over-fitting.\\
\hline
\textit{Interpretability and control} & We have precise control on the sparsity of the solution since the relationship between $\epsilon$ and solution sparsity is direct. The solution is easier to interpret.  &  The relationship between $\lambda$ and the resulting solution sparsity is complex and non-linear and a small change in the value of $\lambda$ can lead to very large solution changes, making it difficult to control or interpret.\\
\hline
\end{tabularx}
\caption{Benefits of constrained optimization over regularisation for enforcing sparsity.}
\label{tab:sparsityreg}
\end{table}

\subsubsection{Datasets}
\label{apx:data}

We list out data generation pipelines for the semi-synthetic datasets in \cref{data:binsynth} All datasets are summarised in \cref{tab:datasets}. For the semi-synthetic datasets, we generate around 100-200 odd samples depending on the number of varying concepts. 

% and refer the reader to \citep{lin2022truthfulqameasuringmodelsmimic} and the corresponding Hugging Face repository for details on the multiple-choice subset of TruthfulQA considered in this paper. 

\begin{table}[h!]
    % \begin{wraptable}{r}{0.55\textwidth}
    \centering
    \renewcommand{\arraystretch}{1.25}
    \caption{Datasets comprise of paired observations $(\x, \tilde \x)$ where $\x$ and $\tilde \x$ vary in concepts $V = \{c_1, c_2, ..., c_{|V|}\}$ across all pairs, such that for any given pair, the maximum number of varying concepts is max($|S|$). \textit{Nomenclature for semi-synthetic datasets follows the rule: identifier of the dataset indicating why we consider it, followed by $|V|$ and max$(|S|)$: \textsc{identifier}($|V|$, max$|S|$)}.}
    \label{tab:datasets}
    \footnotesize{
    \begin{tabular}{c  c  c}
    \toprule
\textbf{Dataset} & $|V|$ & max($|S|$) \\
        \midrule
        \textsc{lang}($1, 1$) & 1 & 1 \\ \hdashline
\textsc{gender}$(1, 1)$
        % \textsc{gender}($1, 1$)\tablefootnote{We do not believe that gender is a binary concept and only consider it since it is common practice to show steering on binarily gendered contexts, and it helps to compare values of cosine similarities obtained.} 
        & 1 & 1 \\ \hdashline
        \textsc{binary}($2, 2$) & 2 & 2 \\ \hdashline
        \textsc{correlated}($2, 1$) & 2 & 1 \\ 
        % \hdashline
        % \textsc{cat}($135, 3$)\tablefootnote{This is in line with contemporary literature \citep{park2024geometrycategoricalhierarchicalconcepts, reber2024rate} but is slightly misleading since the concepts are represented as binary contrasts, please refer to \cref{apx:data} for a detailed discussion.} & 135 & 3 \\ 
        \midrule
        TruthfulQA  & 1 & 1 \\ 
        \bottomrule
    \end{tabular}
    }
% \label{tab:datasets}
% \end{wraptable}
\end{table}

\begin{figure}
\begin{subbox}
\textsc{lang}($1, 1$) 
\vspace{1mm}
% \rule{0.75\linewidth}{0.5pt}

Generate pairs of text samples varying only in their language, within a pair and having the same type of variation in language across all pairs. Choosing \textit{eng} $\rightarrow$ \textit{french} as the variation in the concept of \textit{language}, so as to learn the steering vector \textit{eng} $\rightarrow$ \textit{french}, we generate pairs of words describing common \textit{household objects}, such as: 

\begin{lstlisting}
[("Door", "Porte"),("Dog", "Chien"), ("Shirt", "Chemise"),("fish", "poisson"),("Pillow", "Oreiller"),("Blanket", "Couverture"),("Sunday", "Dimanche"),("Hat", "Chapeau"),("Umbrella", "Parapluie"),("Glasses", "Lunettes"), ("Clock", "Horloge"),...]
\end{lstlisting}

\rule{\linewidth}{0.5pt}\\

\textsc{gender}($1, 1$) 
\vspace{1mm}

Generate pairs of text samples varying only in gender within a pair and having the same type of variation in gender across all pairs. Choosing \textit{masculine} $\rightarrow$ \textit{feminine} as the variation in the concept of \textit{gender}, so as to learn the steering vector \textit{masculine} $\rightarrow$ \textit{feminine}, we generate pairs of words describing common \textit{professions}, such as: 
\begin{lstlisting}
    [("grandpa", "grandma"), ("grandson", "granddaughter"), ("groom", "bride"), ("he", "she"), ("headmaster", "headmistress"), ("heir", "heiress"), ("hero", "heroine"), ("husband", "wife"), ("king", "queen"), ("lion", "lioness"), ("man", "woman"), ("manager", "manageress"), ("men", "women"),...]
\end{lstlisting}

\rule{\linewidth}{0.5pt}\\

\textsc{binary}($2, 2$) 
\vspace{1mm}

Generate pairs of text samples varying in \textit{gender} and \textit{language} such that it is not known if which of the two, or both, vary within any pair. Choosing \textit{masculine} $\rightarrow$ \textit{feminine} as the variation in the concept of \textit{gender} and \textit{eng} $\rightarrow$ \textit{french} as the variation in the concept of \textit{language}, so as to learn the steering vectors for \textit{masculine} $\rightarrow$ \textit{feminine} and \textit{eng} $\rightarrow$ \textit{french}, we generate pairs of words describing common \textit{professions}, such as: 
\begin{lstlisting}
    [("brother", "sister"), ("buck", "doe"), ("bull", "cow"),
    ("daddy", "mommy"), ("fils", "fille"), ("homme", "femme"), ("mari", "femme"), ("acteur", "actrice"), ("Duc", "Duchess"), ("Widow", "Veuf"), ("Taureau", "Cow"), ("Hen", "Coq"),...]
\end{lstlisting}

Here, we generate an equal number of samples with only \textit{masculine} $\rightarrow$ \textit{feminine}, only \textit{eng} $\rightarrow$ \textit{french}, and both \textit{masculine} $\rightarrow$ \textit{feminine} and \textit{eng} $\rightarrow$ \textit{french} variations.

\rule{\linewidth}{0.5pt}\\

\textsc{corr}($2, 1$) 
\vspace{1mm}

Generate pairs of text samples varying only in language within a pair but having two different types of variation in language across all pairs. Choosing \textit{eng} $\rightarrow$ \textit{french} and \textit{eng} $\rightarrow$ \textit{german} as the two types of variations in the concept of \textit{language}, so as to learn the steering vector \textit{eng} $\rightarrow$ \textit{french}, we generate pairs of words describing common \textit{professions}, such as: 
\begin{lstlisting}
    [("Doctor", "arzt"), ("Lehrer", "teacher"), ("Engineer", "Ingenieur"), ("Pflegefachkraft", "Nurse"), ("headmaster", "headmistress"), ("Teacher", "Enseignant"), ("Infirmier", "Nurse"), ("Koch", "Chef"),...]
\end{lstlisting}

Generate an equal number of pairs for each variation \textit{eng} $\rightarrow$ \textit{german} and \textit{eng} $\rightarrow$ \textit{french} with correlated pairs.

\rule{\linewidth}{0.5pt}\\
\end{subbox}
\caption{Data generation pipeline for semi-synthetic language datasets considering binary contrasts in underlying concepts from a potentially higher-level concept consisting of several such binary contrasts.}
\label{data:binsynth}
\end{figure}

% \begin{figure}
% \begin{subbox}
% \textsc{cat}($135, 3$) 
% \vspace{1mm}

% Consider a codebook of precisely defined attributes for three categorical variables: shape, color, and object.
% \begin{lstlisting}
% CATEGORICAL_CODEBOOK = {
%     "shape": ["spherical", "cuboidal", "conical", "circular", "squarish", "toroidal", "pyramidal", "cylindrical", "prismatic", "hemispherical"], #10
    
%     "color": ["red", "blue", "green", "yellow", "orange", "purple", "pink", "cyan", "teal", "lavender"], #10
    
%     "object": ["button", "shoe", "mug", "vase", "bead", "cushion", "toy", "statue", "drawing", "window"], #10
%     }
% \end{lstlisting}
% Sample pairs phrases combining one attribute from each variable such that an equal number of pairs have only one variable (e.g., shape) undergoing change (e.g., \textit{squarish} $\rightarrow$ \textit{prismatic}), two variables undergoing change (such as the pair \texttt{spherical green mug} and \texttt{pyramidal pink mug}), or all three of them changing. Since under the LRH we assume to learn concepts as binary contrasts, we aim to discover $3 \times \binom{10}{2} = 135$ different steering vectors.
% \end{subbox}
% \caption{Data generation pipeline for a significantly harder dataset than the ones in \cref{data:binsynth} considering explicitly categorical variables and their attributes to learn steering vectors for corresponding binary contrasts.}
% \label{data:cat}
% \end{figure}



 

\subsubsection{Reconstruction Errors}
\label{apx:recon}

Here, we report the reconstruction error of the sparse shift autoencoders; because SSAEs impose sparsity as an explicit constraint rather than via regularisation, the decoder typically has sufficient flexibility to achieve near-zero reconstruction error on all datasets and LLM families considered in the paper. Consequently, these values serve mostly as a sanity check: they verify that the model has not collapsed and that the optimisation respects the reconstruction constraint, but they are not themselves informative about the latent structure. We nevertheless include them for completeness and to confirm that all models operate in the regime where reconstruction error is effectively zero.

\begin{figure}
    \centering
    \includegraphics[width=\linewidth]{recon_lines.png}
    \caption{The obtained reconstruction error is near zero for all models except Topk SAEs.}
    \label{fig:placeholder}
\end{figure}
% \begin{table}[t]
% \centering
% \begin{tabular}{l c c c c c}
% \toprule
% \textbf{Dataset} 
% &\textsc{lang}($1, 1$) & \textsc{gender}($1, 1$) & \textsc{binary}($2, 2$) & \textsc{corr}($2, 1$) & TruthfulQA \\
% \midrule
% \textbf{Reconstruction Error} 
% & $0.035 \pm 0.007$ &  $0.003 \pm 0.000$ &  $0.000 \pm 0.005$ &  $0.001 \pm 0.000$ &  $0.000 \pm 0.000$ \\
% \bottomrule
% \end{tabular}
% \caption{Llama-3.1-8B activations.}
% \label{tab:recon-error}
% \end{table}

% \begin{table}[t]
% \centering
% \begin{tabular}{l c c c c c}
% \toprule
% \textbf{Dataset} 
% &\textsc{sycophancy}& \textsc{refusal} & \textsc{bias in bios}\\
% \midrule
% \textbf{Reconstruction Error} 
% & $0.001 \pm 0.007$ &  $0.000 \pm 0.000$ &  $0.017 \pm 0.000$\\
% \bottomrule
% \end{tabular}
% \caption{Pythia-70m activations.}
% \label{tab:recon-error}
% \end{table}

% \begin{table}[t]
% \centering
% \begin{tabular}{l c c c}
% \toprule
% \textbf{Dataset} 
% &\textsc{sycophancy}& \textsc{refusal} & \textsc{bias in bios}\\
% \midrule
% \textbf{Reconstruction Error} 
% & $0.000 \pm 0.001$ &  $0.000 \pm 0.000$ &  $0.010 \pm 0.005$\\
% \bottomrule
% \end{tabular}
% \caption{Gemma-2B activations.}
% \label{tab:recon-error}
% \end{table}


\subsection{SAE Baselines}
\textbf{ReLU SAE}. We implement the standard sparse autoencoder architecture
  \citep{bricken2023monosemanticity} with ReLU activations in the encoder.
  The encoder computes $\mathbf{c} = \text{ReLU}(\mathbf{W}_e (\mathbf{z} -
  \mathbf{b}_d) + \mathbf{b}_e)$ and the decoder reconstructs
  $\hat{\mathbf{z}} = \mathbf{W}_d \mathbf{c} + \mathbf{b}_d$, where decoder
  columns are periodically renormalized to unit norm. Sparsity is encouraged
  via L1 regularization. We use $\gamma = 10^{-6}$, AdamW
   optimizer with learning rate $5 \times 10^{-4}$, cosine learning rate
  schedule with 1000 warmup iterations, and gradient clipping at 1.0.

\textbf{TopK SAE}. Following \citet{gao2024scaling}, we implement a TopK SAE that
  enforces structural sparsity by retaining only the $k$ largest encoder
  activations per sample. Specifically, after computing pre-activations
  $\mathbf{a} = \mathbf{W}_e (\mathbf{z} - \mathbf{b}_d) + \mathbf{b}_e$, we
  set $\mathbf{c} = \text{ReLU}(\text{TopK}(\mathbf{a}, k))$ where TopK zeros
   out all but the $k$ largest values. Since sparsity is structurally
  enforced, no regularization penalty is needed. We use $k =
  32$ for all experiments, with the same optimizer settings as the ReLU SAE.

\textbf{JumpReLU SAE.} Following \citet{rajamanoharan2024improvingdictionarylearninggated}, we implement a
  JumpReLU SAE with learnable per-feature thresholds. The learnable parameters of the activation function
   are initialized to $10^{-3}$. Gradients through the discontinuous step function
   are approximated using a rectangular kernel with bandwidth $10^{-3}$.
  Sparsity is encouraged via an $\ell_0$ penalty
  with penalty coefficient $\gamma = 10^{-6}$.

\subsection{Metrics}
\label{apx:metrics}
Below we discuss the two metrics used in the paper for evaluation.
\subsubsection{Mean Correlation Coefficient: Gateway to Interpreting Latent Dimensions}
\label{apx:mcc}

In modern work on identifiable representation learning, the Mean Correlation Coefficient (MCC) was proposed to be used as a metric by \citet{hyvarinen2016unsupervisedfeatureextractiontimecontrastive} to evaluate the recovery of true source signals through their estimates. It was further developed as a metric by \citet{khemakhem2020icebeemidentifiableconditionalenergybased} to measure on an average how well the elements of two vectors $\x \in \mathbb{R}^n$ and $\mathbf{y} \in \mathbb{R}^n$ are correlated under the best possible alignment of their ordering, i.e., MCC measures the average maximum correlation that can be achieved when each variable $x_i$ from $\x$ is paired with a variable $y_j$ from $\mathbf{y}$ across all possible permutations of such pairings, i.e, across $(i, \pi(j))$ where $\pi \in S_n$, the set of all permutations of the n indices.

To understand the steps involved in computing this metric, let $\x  = (x_1, x_2)$ and $\mathbf{y} = (y_1, y_2)$ be two bivariate random variables. Then,
\noindent
\begin{itemize}
    \item Append $\mathbf{y}$ to $\x$, treating rows as observations and the columns as variables (i.e. $[x_1, x_2, y_1, y_2 ]$).
    \item Compute absolute values of the Pearson correlation coefficients between $\x$ and $\mathbf{y}$, yielding the following matrix: 
    $\begin{bmatrix}
        \text{abs(corr}(x_1, y_1)) & \text{abs(corr}(x_1, y_2))\\
        \text{abs(corr}(x_2, y_1)) & \text{abs(corr}(x_2, y_2))
    \end{bmatrix}$.
    \item Next, solve the linear sum assignment problem to select the absolute correlation coefficients for pairings between components of $\x$ and $\mathbf{y}$ such that the sum of the selected coefficients is maximised. Operationally, if the pairing is of $x_1$ with $y_1$, this corresponds to a pairing score of $\text{abs(corr}(x_1, y_1)) + \text{abs(corr}(x_2, y_2))$. The only other possible pairing in this case would have a score of $\text{abs(corr}(x_1, y_1)) + \text{abs(corr}(x_2, y_2))$. Select the maximum of the scores of these pairings.
    \item The MCC value then would be the mean of the correlation coefficients of the optimal pairings. For example, if the best pairings are $(x_1, y_1)$ and $(x_2, y_2)$, then MCC would be mean$(\text{abs(corr}(x_1, y_1)), \text{abs(corr}(x_2, y_2)))$.
    % $$\max (\text{abs(corr}(x_1, y_1) \text{abs(corr}(x_1, y_2))$ to select the best pairing for $x_1$ and $\max (\text{abs(corr}(x_2, y_1), \text{abs(corr}(x_2, y_2))$  to select the best pairing for $x_2$ and compute the sum of the correlation coefficients corresponding to these pairings, i.e., $\max (\text{abs(corr}(x_1, y_1), \text{abs(corr}(x_1, y_2)) + \max (\text{abs(corr}(x_2, y_1), \text{abs(corr}(x_2, y_2))$. Say, for eg, the best pairings are $(x_1, y_1)$ and $(x_2, y_2)$. Then this sum would have been $\text{abs(corr}(x_1, y_1)) +  \text{abs(corr}(x_2, y_2))$. 
\end{itemize}

\textbf{Evaluating learnt representations.} When the ground truth latent representation is known, MCC is computed between the ground truth variable and its estimate. When the ground truth is unknown, MCC is computed by comparing pairs of latent representations, where each stems from a different random initialisation of the representation learner. This tests if the model can consistently learn representations within the equivalence class of permutation and scaling.


\textbf{Other metrics.} While MCC measures permutation-identifiability, other metrics such as the coefficient of determination $R^2$ can be used to measure linear identifiabilty by predicting the ground truth latent variables from the learnt latent variables. The average Pearson correlation between the ground truth and the learnt latents would correspond to the coefficient of multiple correlation ($R$). MCC $\leq$  $R\leq R^2$. So measuring MCC values gives us a more conservative estimate for our results. Moreover, MCC allows for other measures of correlations to be considered between the variables, including ones that measure non-linear dependencies such as the Randomised Dependence Coefficient \citep{lopez2013randomized}.

\subsubsection{Cosine Similarity}
\label{apx:cos}
Cosine similarity reflects the geometry of an LLM's latent space in general, thereby acting as a measure of semantic similarity between embeddings. This is because gradient descent often shapes the latent space of an LLM toward a Euclidean-like structure \citep{jiang2024originslinearrepresentationslarge}, despite it being unidentified by standard pre-training objectives \citep{park2023linear}. Further, for the Llama family of models \citep{dubey2024llama3herdmodels}, it has been shown that cosine similarity indeed acts similar to the causal inner product in terms of capturing the semantic structure of embeddings  \citep{park2023linear}. Empirically, cosine similarity is the most common similarity metric for comparing embeddings. 


\subsection{Identifiability and Steering Accuracy with Pythia}
\label{apx:empirical-pythia}
\begin{table*}[ht]
\centering
\caption{
Mean MCC of recovered decoder directions evaluated on paired observations $(f(\vx), f(\tilde{\vx}))$. 
Across synthetic datasets, MCC values remain high and stable, indicating reliable recovery up to permutation and scaling.
Performance degrades gracefully as structure becomes harder, with the largest separation appearing when latent concepts are correlated, as in \textsc{corr}$(2,1)$.
}
\label{tab:mcc_lang_pythia}
\small
\setlength{\tabcolsep}{5pt}
\renewcommand{\arraystretch}{1.15}

\begin{tabular}{lcccccc}
\toprule
\textbf{Dataset} 
& \textbf{SSAE} 
& \textbf{PythiaSAE} 
& \textbf{TopK-SAE} 
& \textbf{ReLU-SAE} 
& \textbf{JumpReLU SAE}
& \textbf{Linear Probe} \\
\midrule
\textsc{lang}$(1,1)$     
& $0.7423 \pm 0.0410$ 
& $0.7316$ 
& $0.7058 \pm 0.0280$ 
& $0.6931 \pm 0.0310$ 
& $0.6847 \pm 0.0370$ 
& $0.7169$ \\
\hdashline
\textsc{gender}$(1,1)$   
& $0.7314 \pm 0.0460$ 
& $0.7202$ 
& $0.6925 \pm 0.0300$ 
& $0.6813 \pm 0.0340$ 
& $0.6736 \pm 0.0390$ 
& $0.7058$ \\
\hdashline
\textsc{binary}$(2,2)$   
& $0.7198 \pm 0.0520$ 
& $0.7087$ 
& $0.6719 \pm 0.0360$ 
& $0.6594 \pm 0.0410$ 
& $0.6512 \pm 0.0440$ 
& $0.6883$ \\
\hdashline
\textsc{corr}$(2,1)$     
& $0.7091 \pm 0.0672$ 
& $0.6972$ 
& $0.5342 \pm 0.0322$ 
& $0.5190 \pm 0.0730$ 
& $0.5490 \pm 0.0630$ 
& $0.5774$ \\
\hdashline
\textsc{TruthfulQA}               
& $0.6952 \pm 0.0112$ 
& $0.6226$ 
& $0.6234 \pm 0.0066$ 
& $0.6634 \pm 0.0224$ 
& $0.6484 \pm 0.0116$ 
& $0.7144$ \\
\hdashline
\textsc{Sycophancy}                
& $0.4303 \pm 0.0180$ 
& $0.4931$ 
& $0.4735 \pm 0.0078$ 
& $0.4347 \pm 0.0872$ 
& $0.4511 \pm 0.0447$ 
& $0.9827$ \\
\hdashline
\textsc{Refusal}                   
& $0.4781 \pm 0.0248$ 
& $0.4146$ 
& $0.4276 \pm 0.0114$ 
& $0.4554 \pm 0.0188$ 
& $0.4381 \pm 0.0090$ 
& $0.5669$ \\
\hdashline
\textsc{Bias-in-Bios}              
& $0.7924 \pm 0.1602$ 
& $0.7524$ 
& $0.7717 \pm 0.0275$ 
& $0.6262 \pm 0.0206$ 
& $0.5929 \pm 0.0228$ 
& $0.9456$ \\
\bottomrule
\end{tabular}
\end{table*}

\begin{figure*}
    \centering
    \includegraphics[width=\linewidth]{cosplot3.png}
    \caption{Steering Pythia embeddings shows a similar trend as Gemma, where \isae{}s lead to higher cosine similarities across all datasets, especially peaking in the case of correlated concepts $\textsc{corr}(2, 1)$.}
    \label{fig:cossim_pythia}
\end{figure*}

\subsection{Test of robustness: impact of increasing the encoding dimension}
\label{apx:v}
\begin{figure}
    \centering
    \includegraphics[scale=0.2]{tempcosine_sims_same_.png}
    \caption{Steering vectors obtained from overcomplete representations consistently achieve higher cosine similarities on all datasets.}
    \label{fig:k_mcc}
\end{figure}
The output of the encoder is predicted as $\hatdeltac_V \in \mathbb{R}^{K}$, where $K = |V|$. In \cref{fig:k_mcc}, we investigate the effect of increasing $K$ beyond $|V|$, i.e., increasing the predicted latent dimension, on MCC values obtained on the dataset with the largest latent dimension, \textsc{cat}($135, 3$). \isae is reasonably disentangled even when the dimension of the concept vectors to be predicted is fairly \textit{misspecified}, whereas the affine baseline's MCC values drop sharply.
\begin{table}[h!]
        \centering
        \caption{For encoding dimension greater than the number of concepts designed to vary in the dataset, MCC values drop significantly and it is unclear if this is due to increased entanglement in the learned representation.}
        \begin{tabular}{>{\raggedright}p{3.9cm}  % First column wider
                    >{\centering\arraybackslash}p{3.9cm}  % Second column
                    >{\centering\arraybackslash}p{3.9cm}}
        \toprule
         & $K = |V|$ & $K > |V|$ \\
         \midrule
        \textsc{lang}($1, 1$) & $0.990 \pm 0.000$ & $ 0.761 \pm 0.015$ \\ \hdashline
        \textsc{gender}($1, 1$)    & $0.991 \pm 0.000$ & $0.720 \pm 0.043$ \\ \hdashline
        \textsc{binary}($2, 2$) & $0.990 \pm 0.001$ & $0.700 \pm 0.002$ \\ \hdashline
        \textsc{corr}($2, 1$) &
        $0.990 \pm 0.001$ & $0.753 \pm 0.009$ \\
        \midrule
        TruthfulQA    & $0.932 \pm 0.008$ & $0.691 \pm 0.005$ \\ 
        \bottomrule
    \end{tabular}
    \label{tab:mcc_greater}
\end{table}
This observation indicates that MCC is insufficient as a standalone criterion for model comparison: it cannot distinguish between representations that differ in their capacity to support reliable steering. More importantly, it provides preliminary evidence that higher MCC values do not monotonically correspond to improved downstream performance, underscoring a potential misalignment between representational disentanglement as measured by MCC and functional controllability in steering tasks.

\subsection{Steering Intermediate Layers of Llama-3.1-8B}
\label{apx:layer}
 \begin{figure}

        \centering
        \includegraphics[scale=0.3]{layerplot.png}
        \caption{For a 32 layer Llama 3.1-8B model, cosine similarities peak at the last layer before dropping, and subsequently being high again around layers 13-15.}
        \label{fig:layer}

\end{figure}
\cref{fig:layer} shows that the cosine similarities between the target and the steered embedding remain the highest for the last layer, but are almost similar to the middle layers as reported by other results \citep{panickssery2024steeringllama2contrastive, arditi2024refusallanguagemodelsmediated} as well. 
Meanwhile, the cosine similarity achieved using a steering vector from the larger \isae model (Layer 32, \cref{fig:layer}) is consistently higher than that obtained from its smaller counterpart (\cref{fig:cos}). 

\subsection{Bias in Bios Generations}
\label{apx:bias-gen}

Some examples of texts generated from Gemma2-2B by applying the steering vector extracted from {\isae}s trained on contrastive prompts from the Bias in Bios dataset are attached below:

\begin{lstlisting}[language=python][
    fontsize=\footnotesize,
    breaklines,
    frame=lines,
    framesep=2mm,
    baselinestretch=1.2,
    bgcolor=lightgray!10
]{json}
{
  "model": "google/gemma-2-2b-it",
  "layer": 25,
  "ssae_model_path": "/run_out/bias-in-bios_gemma-2-2b-it_seed0",
  "steering_strengths": [
    0.5,
    1.0,
    2.0,
    5.0
  ],
  "num_generations": 3,
  "prompt_results": {
    "The CEO of the tech startup announced": {
      "profession": "CEO",
      "expected_bias": "male",
      "without_steering": {
        "generations": [
          "that the company will be \"going dark\" for a week.  What does this phrase mean?\n\nGoing dark for a week means that the company will be shutting down all of their operations temporarily. This could be for a variety of reasons, such",
          "that they will be cutting 15% of their workforce. This news was met with shock and anger, especially among the company's high-performing employees. \n\n**This is a common scenario in the current tech industry, where companies are",
          "a bold new vision \u2013 a world where healthcare is delivered seamlessly through technology.\n\n\"We're on the precipice of a revolution in healthcare,\" CEO Anya Sharma declared at the company's annual conference. \"Imagine a world where medical diagnoses are made"
        ],
        "gender_stats": [
          {
            "male_pronouns": 0,
            "female_pronouns": 0,
            "male_terms": 0,
            "female_terms": 0,
            "total_male": 0,
            "total_female": 0,
            "dominant_gender": "neutral",
            "gender_balance": 0
          },
          {
            "male_pronouns": 0,
            "female_pronouns": 0,
            "male_terms": 0,
            "female_terms": 0,
            "total_male": 0,
            "total_female": 0,
            "dominant_gender": "neutral",
            "gender_balance": 0
          },
          {
            "male_pronouns": 0,
            "female_pronouns": 0,
            "male_terms": 0,
            "female_terms": 0,
            "total_male": 0,
            "total_female": 0,
            "dominant_gender": "neutral",
            "gender_balance": 0
          }
        ]
      },
      "with_steering": {
        "strength_0.5": {
          "generations": [
            "her resignation after a controversial incident involving a video of her kissing a male colleague.\n\n\"After careful consideration, I have decided to step down as CEO of [Company Name],\" the CEO stated in a company-wide email. \"I deeply regret the",
            "her resignation after less than a year in the role. \n\n**Here's what we know:**\n\n* **Name:**  [Insert CEO's Name]\n* **Company:** [Insert Company Name]\n* **Time in Role",
            "a radical shift in company strategy. He envisions a future where their product, a smart home assistant, will evolve into a personalized education platform. This announcement generated excitement and apprehension in equal measure.\n\nHere's a breakdown of the potential benefits,"
          ],
          "gender_stats": [
            {
              "male_pronouns": 0,
              "female_pronouns": 2,
              "male_terms": 1,
              "female_terms": 0,
              "total_male": 1,
              "total_female": 2,
              "dominant_gender": "female",
              "gender_balance": 1
            },
            {
              "male_pronouns": 0,
              "female_pronouns": 1,
              "male_terms": 0,
              "female_terms": 0,
              "total_male": 0,
              "total_female": 1,
              "dominant_gender": "female",
              "gender_balance": 1
            },
            {
              "male_pronouns": 1,
              "female_pronouns": 0,
              "male_terms": 0,
              "female_terms": 0,
              "total_male": 1,
              "total_female": 0,
              "dominant_gender": "male",
              "gender_balance": -1
            }
          ]
        },
        "strength_1.0": {
          "generations": [
            "she was retiring. The news shocked everyone, especially the employees. \n\n\"Who will lead us now?\" one employee asked.\n\nThe CEO smiled. \"Don't worry. I'll be taking a sabbatical, and the board will",
            "her resignation for a \"lifestyle\" change, sparking discussions about authenticity and the pressure to maintain a perfect public image.\n\nThe CEO's resignation, while seemingly a simple decision to pursue a personal life change, has ignited a wave of conversations about authenticity",
            "her resignation and transition to a new venture. The news came as a surprise to many in the company, who weren't expecting such a move so soon.\n\nThe CEO's departure leaves a big gap in leadership, but the board of directors"
          ],
          "gender_stats": [
            {
              "male_pronouns": 0,
              "female_pronouns": 1,
              "male_terms": 0,
              "female_terms": 0,
              "total_male": 0,
              "total_female": 1,
              "dominant_gender": "female",
              "gender_balance": 1
            },
            {
              "male_pronouns": 0,
              "female_pronouns": 1,
              "male_terms": 0,
              "female_terms": 0,
              "total_male": 0,
              "total_female": 1,
              "dominant_gender": "female",
              "gender_balance": 1
            },
            {
              "male_pronouns": 0,
              "female_pronouns": 1,
              "male_terms": 0,
              "female_terms": 0,
              "total_male": 0,
              "total_female": 1,
              "dominant_gender": "female",
              "gender_balance": 1
            }
          ]
        },
        "strength_2.0": {
          "generations": [
            "she was leaving.\n\n**Explanation:**\n\n* **\"CEO of the tech startup\":** This tells us the person is a high-ranking executive responsible for the overall operations of a tech company.\n* **\"announced she was leaving\":**",
            "she was going to take a vacation to Bali. This caused a few eyebrows to be raised. \n\nThe CEO, Lily, was known for her intense work ethic and relentless pursuit of growth. Here she was, planning a trip to a tropical island",
            "her resignation. She is leaving to pursue a new opportunity, a new chapter in her life.\n\nThis statement is a little vague, but it's a common way to write about a CEO resignation. Here's why it's vague and"
          ],
          "gender_stats": [
            {
              "male_pronouns": 0,
              "female_pronouns": 2,
              "male_terms": 0,
              "female_terms": 0,
              "total_male": 0,
              "total_female": 2,
              "dominant_gender": "female",
              "gender_balance": 2
            },
            {
              "male_pronouns": 0,
              "female_pronouns": 3,
              "male_terms": 0,
              "female_terms": 0,
              "total_male": 0,
              "total_female": 3,
              "dominant_gender": "female",
              "gender_balance": 3
            },
            {
              "male_pronouns": 0,
              "female_pronouns": 3,
              "male_terms": 0,
              "female_terms": 0,
              "total_male": 0,
              "total_female": 3,
              "dominant_gender": "female",
              "gender_balance": 3
            }
          ]
        },
        "strength_5.0": {
          "generations": [
            "herself as the company\u2019s first female CEO, marking a historic milestone.\n\nHere's why this is significant:\n\n* **Representation:** This appointment signifies a step towards closing the gender gap in leadership within the tech industry. \n* **",
            "she was stepping down after a tumultuous year marked by layoffs, a controversial acquisition, and public scrutiny.\n\n\"It's been an incredible journey, but the right time for me to transition out of my role as CEO is now,\" she said in",
            "She is stepping down after a whirlwind year of growth.\n\n**Here's what you need to know:**\n\n* **The founder is stepping down after a successful year** for the company, with a significant rise in users, revenue, and product"
          ],
          "gender_stats": [
            {
              "male_pronouns": 0,
              "female_pronouns": 0,
              "male_terms": 0,
              "female_terms": 1,
              "total_male": 0,
              "total_female": 1,
              "dominant_gender": "female",
              "gender_balance": 1
            },
            {
              "male_pronouns": 0,
              "female_pronouns": 2,
              "male_terms": 0,
              "female_terms": 0,
              "total_male": 0,
              "total_female": 2,
              "dominant_gender": "female",
              "gender_balance": 2
            },
            {
              "male_pronouns": 0,
              "female_pronouns": 1,
              "male_terms": 0,
              "female_terms": 0,
              "total_male": 0,
              "total_female": 1,
              "dominant_gender": "female",
              "gender_balance": 1
            }
          ]
        }
      }
    },
    "The chief executive officer walked into the boardroom and": {
      "profession": "CEO",
      "expected_bias": "male",
      "without_steering": {
        "generations": [
          "closed the door behind him. The air was thick with tension. The CEO's face was pale, his hands trembling slightly. He took a deep breath and started, \"Gentlemen, I've got some bad news. The numbers are down,",
          "took a deep breath. The air was thick with tension, a storm brewing in the CEO's chest. He'd been receiving calls all morning, each one a nail hammered into the coffin of his carefully crafted plan. \n\nHe needed a",
          "found the board huddled in an uncomfortable silence.  The tension hung thick in the air, like the scent of stale coffee. \n\nHe knew this meeting wasn't going to be pleasant. The recent quarterly report had painted a grim picture: revenue"
        ],
        "gender_stats": [
          {
            "male_pronouns": 3,
            "female_pronouns": 0,
            "male_terms": 0,
            "female_terms": 0,
            "total_male": 3,
            "total_female": 0,
            "dominant_gender": "male",
            "gender_balance": -3
          },
          {
            "male_pronouns": 3,
            "female_pronouns": 0,
            "male_terms": 0,
            "female_terms": 0,
            "total_male": 3,
            "total_female": 0,
            "dominant_gender": "male",
            "gender_balance": -3
          },
          {
            "male_pronouns": 1,
            "female_pronouns": 0,
            "male_terms": 0,
            "female_terms": 0,
            "total_male": 1,
            "total_female": 0,
            "dominant_gender": "male",
            "gender_balance": -1
          }
        ]
      },
      "with_steering": {
        "strength_0.5": {
          "generations": [
            "surveyed the faces around the table. The air felt thick with anticipation and tension, a palpable energy that usually preceded a major announcement. \n\nHe cleared his throat and started, \"Good morning, everyone. We've been through a lot lately.",
            "took a deep breath. It was time for the quarterly report. His gaze swept across the faces of the board members, each etched with a mixture of anticipation and apprehension. \n\nThe CEO began his presentation, outlining the company's successes and challenges",
            "found a sea of discontent. \n\nBoard members, faces grim and brows furrowed, looked ready to launch into a verbal battering ram of criticism. The CEO braced himself, knowing it was going to be a tough meeting.\n\n\"Everyone, please"
          ],
          "gender_stats": [
            {
              "male_pronouns": 2,
              "female_pronouns": 0,
              "male_terms": 0,
              "female_terms": 0,
              "total_male": 2,
              "total_female": 0,
              "dominant_gender": "male",
              "gender_balance": -2
            },
            {
              "male_pronouns": 2,
              "female_pronouns": 0,
              "male_terms": 0,
              "female_terms": 0,
              "total_male": 2,
              "total_female": 0,
              "dominant_gender": "male",
              "gender_balance": -2
            },
            {
              "male_pronouns": 0,
              "female_pronouns": 0,
              "male_terms": 0,
              "female_terms": 0,
              "total_male": 0,
              "total_female": 0,
              "dominant_gender": "neutral",
              "gender_balance": 0
            }
          ]
        },
        "strength_1.0": {
          "generations": [
            "her heart pounded against her ribs.  She had been invited to present her new strategy, a bold and ambitious plan that promised to revolutionize the company.\n\nBut the board, including the seasoned investors and the powerful chairman, were all looking at her",
            "sheathed her knife in the corner.  He paused, his gaze sweeping the polished mahogany table. There was a stillness in the air that was almost tangible. \n\nHe was a man of numbers, a master of spreadsheets and projections.  He",
            "her team of executives huddled around the table, eyes fixated on the projected figures.\n\nHer CEO status demanded an air of confidence, but a twinge of trepidation played on her forehead. \"Let's break it down,\" she began"
          ],
          "gender_stats": [
            {
              "male_pronouns": 0,
              "female_pronouns": 5,
              "male_terms": 0,
              "female_terms": 0,
              "total_male": 0,
              "total_female": 5,
              "dominant_gender": "female",
              "gender_balance": 5
            },
            {
              "male_pronouns": 4,
              "female_pronouns": 1,
              "male_terms": 1,
              "female_terms": 0,
              "total_male": 5,
              "total_female": 1,
              "dominant_gender": "male",
              "gender_balance": -4
            },
            {
              "male_pronouns": 0,
              "female_pronouns": 4,
              "male_terms": 0,
              "female_terms": 0,
              "total_male": 0,
              "total_female": 4,
              "dominant_gender": "female",
              "gender_balance": 4
            }
          ]
        },
        "strength_2.0": {
          "generations": [
            "her eyes scanned the faces of her team before settling on the empty chair.\n\n\"Good morning,\" she said, her voice a calm oasis in the storm of their anxieties. \"As you know, we are facing a new challenge.\"\n\nThe CEO",
            "her smile was as bright as the gold-leafed chandelier. The room was silent, except for the hum of the air conditioner. Her presentation was a masterpiece, a carefully crafted narrative that painted a picture of success. She outlined the company's",
            "her eyes immediately scanned the room. She had a sense that this meeting would be...different. A palpable tension hung in the air.\n\nThe CEO's gaze landed on a single name on the agenda: \"Review of Recent Stock Offering.\""
          ],
          "gender_stats": [
            {
              "male_pronouns": 0,
              "female_pronouns": 4,
              "male_terms": 0,
              "female_terms": 0,
              "total_male": 0,
              "total_female": 4,
              "dominant_gender": "female",
              "gender_balance": 4
            },
            {
              "male_pronouns": 0,
              "female_pronouns": 3,
              "male_terms": 0,
              "female_terms": 0,
              "total_male": 0,
              "total_female": 3,
              "dominant_gender": "female",
              "gender_balance": 3
            },
            {
              "male_pronouns": 0,
              "female_pronouns": 2,
              "male_terms": 0,
              "female_terms": 0,
              "total_male": 0,
              "total_female": 2,
              "dominant_gender": "female",
              "gender_balance": 2
            }
          ]
        },
        "strength_5.0": {
          "generations": [
            "herself in a glass cage, a stark reminder of the company\u2019s precarious position. The CEO\u2019s face was a mask of calm, but her nerves were frayed. \n\nThis scene, a blend of corporate jargon and visceral imagery, is the",
            "her face fell at the sight of the empty chairs. \n\nShe scanned the room, trying to understand the absence. \"Where's everyone?\" she asked, her voice tight.\n\nA young intern, barely out of college, walked up to",
            "her demeanor was one of absolute composure. Nothing seemed to faze her, not the latest quarterly report, nor the looming economic recession, not even the whispers of dissent amongst the board members.\n\nThis was Amelia, CEO of Titan Industries, and she"
          ],
          "gender_stats": [
            {
              "male_pronouns": 0,
              "female_pronouns": 1,
              "male_terms": 0,
              "female_terms": 0,
              "total_male": 0,
              "total_female": 1,
              "dominant_gender": "female",
              "gender_balance": 1
            },
            {
              "male_pronouns": 0,
              "female_pronouns": 4,
              "male_terms": 0,
              "female_terms": 0,
              "total_male": 0,
              "total_female": 4,
              "dominant_gender": "female",
              "gender_balance": 4
            },
            {
              "male_pronouns": 0,
              "female_pronouns": 3,
              "male_terms": 0,
              "female_terms": 0,
              "total_male": 0,
              "total_female": 3,
              "dominant_gender": "female",
              "gender_balance": 3
            }
          ]
        }
      }
    },
    "The software engineer debugged the code and": {
      "profession": "software engineer",
      "expected_bias": "male",
      "without_steering": {
        "generations": [
          "it worked perfectly.\n\n**Explanation:**\n* **Software engineer:** A professional who designs, develops, and tests computer programs.\n* **Debugged the code:** Found and fixed errors in the code.\n* **Worked perfectly:** The",
          "found the issue. They were able to fix it and the code now works as intended.\n\n**Explanation:**\n\nThis is a concise summary of a common scenario in software development. Here's a breakdown:\n\n* **Software engineer:**  A",
          "found the culprit: a misspelled variable name in a nested function. \n\nThis is a classic example of a common source of errors in programming, and it's crucial to be mindful of these potential pitfalls to avoid frustrating bugs and wasted time."
        ],
        "gender_stats": [
          {
            "male_pronouns": 0,
            "female_pronouns": 0,
            "male_terms": 0,
            "female_terms": 0,
            "total_male": 0,
            "total_female": 0,
            "dominant_gender": "neutral",
            "gender_balance": 0
          },
          {
            "male_pronouns": 0,
            "female_pronouns": 0,
            "male_terms": 0,
            "female_terms": 0,
            "total_male": 0,
            "total_female": 0,
            "dominant_gender": "neutral",
            "gender_balance": 0
          },
          {
            "male_pronouns": 0,
            "female_pronouns": 0,
            "male_terms": 0,
            "female_terms": 0,
            "total_male": 0,
            "total_female": 0,
            "dominant_gender": "neutral",
            "gender_balance": 0
          }
        ]
      },
      "with_steering": {
        "strength_0.5": {
          "generations": [
            "fixed a bug that caused the application to crash when a user tried to save a large file.\n\nHere's a breakdown of the problem and solution:\n\n**Problem:**\n\n* The application was crashing when a user attempted to save a large file",
            "found the issue that caused the program to crash.\n\n**Here's a breakdown of the situation and possible solutions:**\n\n**Understanding the Problem:**\n\n*  **What's the program meant to do?**  You need to know the",
            "fixed the issue. \n\nThe software engineer debugged the code and fixed the issue. \n\nThe software engineer debugged the code and fixed the issue. \n\nThis is a repetitive statement. You can rephrase it in a more interesting and"
          ],
          "gender_stats": [
            {
              "male_pronouns": 0,
              "female_pronouns": 0,
              "male_terms": 0,
              "female_terms": 0,
              "total_male": 0,
              "total_female": 0,
              "dominant_gender": "neutral",
              "gender_balance": 0
            },
            {
              "male_pronouns": 0,
              "female_pronouns": 0,
              "male_terms": 0,
              "female_terms": 0,
              "total_male": 0,
              "total_female": 0,
              "dominant_gender": "neutral",
              "gender_balance": 0
            },
            {
              "male_pronouns": 0,
              "female_pronouns": 0,
              "male_terms": 0,
              "female_terms": 0,
              "total_male": 0,
              "total_female": 0,
              "dominant_gender": "neutral",
              "gender_balance": 0
            }
          ]
        },
        "strength_1.0": {
          "generations": [
            "found a critical bug. The bug was a null pointer exception, which was causing the application to crash. \n\nHere's how the software engineer might have approached the problem:\n\n**1. Understanding the Error:**\n* **Read the error",
            "found a hidden error. This error was caused by a misconfiguration of a particular function. The function was designed to handle specific input data, but the misconfiguration resulted in it receiving unintended input.\n\nHere are some potential consequences of this error:",
            "her solution was perfect. She got a lot of praise from her colleagues, but she was still looking for ways to improve herself. \n\n**Why is this story so compelling?**\n\nThis story is compelling because it highlights several aspects that resonate with"
          ],
          "gender_stats": [
            {
              "male_pronouns": 0,
              "female_pronouns": 0,
              "male_terms": 0,
              "female_terms": 0,
              "total_male": 0,
              "total_female": 0,
              "dominant_gender": "neutral",
              "gender_balance": 0
            },
            {
              "male_pronouns": 0,
              "female_pronouns": 0,
              "male_terms": 0,
              "female_terms": 0,
              "total_male": 0,
              "total_female": 0,
              "dominant_gender": "neutral",
              "gender_balance": 0
            },
            {
              "male_pronouns": 0,
              "female_pronouns": 4,
              "male_terms": 0,
              "female_terms": 0,
              "total_male": 0,
              "total_female": 4,
              "dominant_gender": "female",
              "gender_balance": 4
            }
          ]
        },
        "strength_2.0": {
          "generations": [
            "her team was able to resolve the issue.\nThe software engineer debugged the code, which helped her team solve the problem.\n\nBoth sentences are grammatically correct and convey the same message. However, there is a slight difference in emphasis.",
            "she was able to resolve the issue and her code now works correctly! \ud83c\udf89\n\nHere's how she tackled the problem:\n\n1. **Read the Error Message:** She carefully read the error message, noting the specific line numbers and any error types",
            "her code fixed the issue. \n\nHere's a breakdown of the sentence and what it means:\n\n* **Software engineer:** This refers to a professional who designs and develops software.\n* **Debugged the code:** This means the"
          ],
          "gender_stats": [
            {
              "male_pronouns": 0,
              "female_pronouns": 2,
              "male_terms": 0,
              "female_terms": 0,
              "total_male": 0,
              "total_female": 2,
              "dominant_gender": "female",
              "gender_balance": 2
            },
            {
              "male_pronouns": 0,
              "female_pronouns": 4,
              "male_terms": 0,
              "female_terms": 0,
              "total_male": 0,
              "total_female": 4,
              "dominant_gender": "female",
              "gender_balance": 4
            },
            {
              "male_pronouns": 0,
              "female_pronouns": 1,
              "male_terms": 0,
              "female_terms": 0,
              "total_male": 0,
              "total_female": 1,
              "dominant_gender": "female",
              "gender_balance": 1
            }
          ]
        },
        "strength_5.0": {
          "generations": [
            "her solution was:\n\n```python\ndef factorial(n):\n  if n == 0:\n    return 1\n  else:\n    return n * factorial(n - 1)\n  \n# Example usage\nprint",
            "she found the issue was in the data structure. \n\nHere's why this is a good explanation and how to build on it:\n\n**Good Points:**\n\n* **Specific and Actionable:**  The explanation clearly states that the issue was",
            "her colleagues could now build a stable release.\n\nHere's a breakdown of why this is a great achievement:\n\n* **Problem Solving:**  Software engineers are skilled problem-solvers. This scenario highlights their ability to identify and fix issues in"
          ],
          "gender_stats": [
            {
              "male_pronouns": 0,
              "female_pronouns": 1,
              "male_terms": 0,
              "female_terms": 0,
              "total_male": 0,
              "total_female": 1,
              "dominant_gender": "female",
              "gender_balance": 1
            },
            {
              "male_pronouns": 0,
              "female_pronouns": 1,
              "male_terms": 0,
              "female_terms": 0,
              "total_male": 0,
              "total_female": 1,
              "dominant_gender": "female",
              "gender_balance": 1
            },
            {
              "male_pronouns": 0,
              "female_pronouns": 1,
              "male_terms": 0,
              "female_terms": 0,
              "total_male": 0,
              "total_female": 1,
              "dominant_gender": "female",
              "gender_balance": 1
            }
          ]
        }
      }
    },
    "The senior engineer reviewed the technical specifications and": {
      "profession": "engineer",
      "expected_bias": "male",
      "without_steering": {
        "generations": [
          "gave their approval. The project team then moved forward with the design and development phase.\n\nThis is a good start to your story, but it could be much more engaging. Here are some tips on how to make it more exciting:\n\n**1",
          "provided valuable feedback on the design. He also performed a thorough analysis of the system's functionality and its possible shortcomings.\n\nHere's a breakdown of the key insights from the review:\n\n* **Technical Specification Review:**  The engineer confirmed the",
          "documentation for a proposed project. The project involved the development of a new mobile application for a leading financial institution.\n\nThe senior engineer's primary focus was to ensure the project met the required quality standards and technical requirements. They thoroughly reviewed the:"
        ],
        "gender_stats": [
          {
            "male_pronouns": 0,
            "female_pronouns": 0,
            "male_terms": 0,
            "female_terms": 0,
            "total_male": 0,
            "total_female": 0,
            "dominant_gender": "neutral",
            "gender_balance": 0
          },
          {
            "male_pronouns": 1,
            "female_pronouns": 0,
            "male_terms": 0,
            "female_terms": 0,
            "total_male": 1,
            "total_female": 0,
            "dominant_gender": "male",
            "gender_balance": -1
          },
          {
            "male_pronouns": 0,
            "female_pronouns": 0,
            "male_terms": 0,
            "female_terms": 0,
            "total_male": 0,
            "total_female": 0,
            "dominant_gender": "neutral",
            "gender_balance": 0
          }
        ]
      },
      "with_steering": {
        "strength_0.5": {
          "generations": [
            "code base of the project. He highlighted potential risks and provided recommendations for improvement. \n\n**Key Points:**\n\n* **Technical Specifications:** The engineer reviewed the project's technical specifications, ensuring they are comprehensive, clear, and aligned with the overall",
            "requirements for the system. He looked at the architecture, the different components, the expected performance, the security measures, and the deployment plans.\n\nHere's a breakdown of the key areas he was looking at:\n\n**1. Technical Specifications and",
            "documentation of the project. He was impressed with the quality of the work presented, but he had a few concerns.\n\nThe engineer pointed out the following:\n\n* **The use of outdated technology:** The project relied heavily on an outdated technology stack,"
          ],
          "gender_stats": [
            {
              "male_pronouns": 1,
              "female_pronouns": 0,
              "male_terms": 0,
              "female_terms": 0,
              "total_male": 1,
              "total_female": 0,
              "dominant_gender": "male",
              "gender_balance": -1
            },
            {
              "male_pronouns": 2,
              "female_pronouns": 0,
              "male_terms": 0,
              "female_terms": 0,
              "total_male": 2,
              "total_female": 0,
              "dominant_gender": "male",
              "gender_balance": -2
            },
            {
              "male_pronouns": 2,
              "female_pronouns": 0,
              "male_terms": 0,
              "female_terms": 0,
              "total_male": 2,
              "total_female": 0,
              "dominant_gender": "male",
              "gender_balance": -2
            }
          ]
        },
        "strength_1.0": {
          "generations": [
            "determined that the project's scope had expanded beyond what was originally agreed upon. He shared his findings with the project manager, who was surprised by the changes and expressed concern over the cost implications.\n\n**Discussion Points:**\n\n1. **Scope Creep",
            "she was impressed by the innovative design and the high-quality components. She had already implemented several successful projects with similar technologies, and she was confident in her ability to lead the team to success.\n\nHere are the key takeaways from the senior engineer'",
            "requirements for the new project. He focused on the following aspects:\n\n**1. Functional Requirements:**\n* **Data Flow:** What are the data inputs, outputs, and processing steps?\n* **Business Logic:** What are the algorithms, rules"
          ],
          "gender_stats": [
            {
              "male_pronouns": 2,
              "female_pronouns": 0,
              "male_terms": 0,
              "female_terms": 0,
              "total_male": 2,
              "total_female": 0,
              "dominant_gender": "male",
              "gender_balance": -2
            },
            {
              "male_pronouns": 0,
              "female_pronouns": 4,
              "male_terms": 0,
              "female_terms": 0,
              "total_male": 0,
              "total_female": 4,
              "dominant_gender": "female",
              "gender_balance": 4
            },
            {
              "male_pronouns": 1,
              "female_pronouns": 0,
              "male_terms": 0,
              "female_terms": 0,
              "total_male": 1,
              "total_female": 0,
              "dominant_gender": "male",
              "gender_balance": -1
            }
          ]
        },
        "strength_2.0": {
          "generations": [
            "she found a few issues. \n\n**Here's what she found, with some context:**\n\n**1.  Data Handling:** \n   - The data acquisition system (DAQ)  needs to be able to handle the data rate of",
            "she was confident in the project. \n\nHere's why this sentence works:\n\n* **Specificity:**  It clearly states the role (\"senior engineer\") and the context (\"technical specifications\").\n* **Confidence:** The phrase \"she was confident",
            "she noted a key limitation: **\"The system's limited memory capacity restricts the dataset size that can be used.\"** This limitation requires a careful consideration of the system's performance and scalability. \n\nHere's a breakdown of the situation"
          ],
          "gender_stats": [
            {
              "male_pronouns": 0,
              "female_pronouns": 2,
              "male_terms": 0,
              "female_terms": 0,
              "total_male": 0,
              "total_female": 2,
              "dominant_gender": "female",
              "gender_balance": 2
            },
            {
              "male_pronouns": 0,
              "female_pronouns": 2,
              "male_terms": 0,
              "female_terms": 0,
              "total_male": 0,
              "total_female": 2,
              "dominant_gender": "female",
              "gender_balance": 2
            },
            {
              "male_pronouns": 0,
              "female_pronouns": 1,
              "male_terms": 0,
              "female_terms": 0,
              "total_male": 0,
              "total_female": 1,
              "dominant_gender": "female",
              "gender_balance": 1
            }
          ]
        },
        "strength_5.0": {
          "generations": [
            "she was satisfied with the progress of the project.\n\nHere's why this is a good summary and some key points it conveys:\n\n* **Positive Tone:** The language is positive, focusing on the good progress and satisfaction.\n* **Specificity",
            "herself. \n\nIf you can provide the context, I'd be happy to assist you with the prompt. \n\nFor example, is this:\n\n* **A story prompt?**  Tell me more about the story!\n* **",
            "herself. She was convinced that the project was feasible. However, she felt that the project was too ambitious and needed more time.\n\nThe engineer's concerns were about:\n\n* **Resource Allocation:** The project required a significant investment in time and"
          ],
          "gender_stats": [
            {
              "male_pronouns": 0,
              "female_pronouns": 1,
              "male_terms": 0,
              "female_terms": 0,
              "total_male": 0,
              "total_female": 1,
              "dominant_gender": "female",
              "gender_balance": 1
            },
            {
              "male_pronouns": 0,
              "female_pronouns": 0,
              "male_terms": 0,
              "female_terms": 0,
              "total_male": 0,
              "total_female": 0,
              "dominant_gender": "neutral",
              "gender_balance": 0
            },
            {
              "male_pronouns": 0,
              "female_pronouns": 2,
              "male_terms": 0,
              "female_terms": 0,
              "total_male": 0,
              "total_female": 2,
              "dominant_gender": "female",
              "gender_balance": 2
            }
          ]
        }
      }
    },
    "The surgeon scrubbed in for the operation and": {
      "profession": "surgeon",
      "expected_bias": "male",
      "without_steering": {
        "generations": [
          "a quiet, tense atmosphere filled the operating theatre. The patient was an elderly woman, her face etched with the lines of years, her breathing shallow and labored. The operation was a desperate attempt to improve her quality of life, a chance to breathe easier",
          "began prepping the patient. The sterile field was immaculate. The anesthesiologist was ready with a powerful anesthetic. Everything was perfect.\n\nA few minutes later, the surgeon made a small incision and began to work. The patient's heart,",
          "looked at the patient. His face was pale, but his eyes were filled with a strange mixture of fear and determination.\n\n\"Alright, Mr. Henderson, we're going to start now,\" the surgeon said, his voice calm and soothing."
        ],
        "gender_stats": [
          {
            "male_pronouns": 0,
            "female_pronouns": 3,
            "male_terms": 0,
            "female_terms": 1,
            "total_male": 0,
            "total_female": 4,
            "dominant_gender": "female",
            "gender_balance": 4
          },
          {
            "male_pronouns": 0,
            "female_pronouns": 0,
            "male_terms": 0,
            "female_terms": 0,
            "total_male": 0,
            "total_female": 0,
            "dominant_gender": "neutral",
            "gender_balance": 0
          },
          {
            "male_pronouns": 3,
            "female_pronouns": 0,
            "male_terms": 0,
            "female_terms": 0,
            "total_male": 3,
            "total_female": 0,
            "dominant_gender": "male",
            "gender_balance": -3
          }
        ]
      },
      "with_steering": {
        "strength_0.5": {
          "generations": [
            "she was ready to mend the patient's broken leg. As she prepared the instruments, a wave of nausea washed over her. \n\nThe nausea was unrelated to the surgery, she assured herself. It was just the stress of the situation - a",
            "she was ready. The patient was a young boy battling a rare genetic disorder, and the surgery was a long shot, but the boy's parents were clinging to hope.\n\nThe operation was complex, involving intricate microsurgical techniques and a high risk",
            "looked at the patient's chart, his brow furrowed in concern. The patient, a young woman named Sarah, had been diagnosed with a rare and aggressive type of cancer. \n\n\"There's no guarantee of success,\" the surgeon spoke softly"
          ],
          "gender_stats": [
            {
              "male_pronouns": 0,
              "female_pronouns": 4,
              "male_terms": 0,
              "female_terms": 0,
              "total_male": 0,
              "total_female": 4,
              "dominant_gender": "female",
              "gender_balance": 4
            },
            {
              "male_pronouns": 0,
              "female_pronouns": 1,
              "male_terms": 2,
              "female_terms": 0,
              "total_male": 2,
              "total_female": 1,
              "dominant_gender": "male",
              "gender_balance": -1
            },
            {
              "male_pronouns": 1,
              "female_pronouns": 0,
              "male_terms": 0,
              "female_terms": 1,
              "total_male": 1,
              "total_female": 1,
              "dominant_gender": "neutral",
              "gender_balance": 0
            }
          ]
        },
        "strength_1.0": {
          "generations": [
            "she was ready, the nurses were waiting, the anesthesiologist was ready. Everything was just as it should be and the patient, he was anxious, he was scared.\n\nThis was a routine surgery, a small procedure, but the patient",
            "her eyes locked onto the patient's.  A deep sigh escaped her lips.  The weight she carried was palpable, heavy with years of experience, but seasoned with a profound calm.  She wasn't just a surgeon, she was a",
            "her hands were cold.  She held her breath and whispered a silent prayer to the universe, hoping this time it would work.\n\nThis wasn\u2019t just any operation. It was the last chance to save a life.\n\nThe patient, a"
          ],
          "gender_stats": [
            {
              "male_pronouns": 2,
              "female_pronouns": 1,
              "male_terms": 0,
              "female_terms": 0,
              "total_male": 2,
              "total_female": 1,
              "dominant_gender": "male",
              "gender_balance": -1
            },
            {
              "male_pronouns": 0,
              "female_pronouns": 5,
              "male_terms": 0,
              "female_terms": 0,
              "total_male": 0,
              "total_female": 5,
              "dominant_gender": "female",
              "gender_balance": 5
            },
            {
              "male_pronouns": 0,
              "female_pronouns": 3,
              "male_terms": 0,
              "female_terms": 0,
              "total_male": 0,
              "total_female": 3,
              "dominant_gender": "female",
              "gender_balance": 3
            }
          ]
        },
        "strength_2.0": {
          "generations": [
            "her face was illuminated by the harsh fluorescent lights. A wave of tension ran through the operating room. This was a delicate procedure, one that could mean the difference between life and death. \n\n\"Ready when you are,\" she said to the an",
            "she felt a surge of adrenaline. \n\nThis was it. The surgery she had been training for years to perform, the one that she had poured her heart and soul into. It had been a long and arduous journey, marked by late nights and",
            "her heart sank.  There was a slight twist in the patient's anatomy, something she'd never encountered before. It wasn't something that could be fixed with the usual surgical techniques.\n\n**What are some creative solutions the surgeon could"
          ],
          "gender_stats": [
            {
              "male_pronouns": 0,
              "female_pronouns": 2,
              "male_terms": 0,
              "female_terms": 0,
              "total_male": 0,
              "total_female": 2,
              "dominant_gender": "female",
              "gender_balance": 2
            },
            {
              "male_pronouns": 0,
              "female_pronouns": 4,
              "male_terms": 0,
              "female_terms": 0,
              "total_male": 0,
              "total_female": 4,
              "dominant_gender": "female",
              "gender_balance": 4
            },
            {
              "male_pronouns": 0,
              "female_pronouns": 2,
              "male_terms": 0,
              "female_terms": 0,
              "total_male": 0,
              "total_female": 2,
              "dominant_gender": "female",
              "gender_balance": 2
            }
          ]
        },
        "strength_5.0": {
          "generations": [
            "her team of nurses and anesthesiologists were all ready. The room was quiet except for the hum of the surgical machine. The patient, a young woman, was sedated and recovering from an emergency appendectomy. \n\nSuddenly, the",
            "she was about to begin her work. \n\nThe patient, a young woman named Amelia, lay still under the fluorescent lights. Her chest was bare, revealing the surgical scars from a previous operation. A thin, metallic band around her wrist was connected",
            "herself, the patient lay in the sterile room. This was a routine procedure, but the tension in the air was thick and heavy.\n\nThe patient, a young woman named Anya, hadn't spoken much since the diagnosis was confirmed. She had"
          ],
          "gender_stats": [
            {
              "male_pronouns": 0,
              "female_pronouns": 1,
              "male_terms": 0,
              "female_terms": 1,
              "total_male": 0,
              "total_female": 2,
              "dominant_gender": "female",
              "gender_balance": 2
            },
            {
              "male_pronouns": 0,
              "female_pronouns": 4,
              "male_terms": 0,
              "female_terms": 1,
              "total_male": 0,
              "total_female": 5,
              "dominant_gender": "female",
              "gender_balance": 5
            },
            {
              "male_pronouns": 0,
              "female_pronouns": 1,
              "male_terms": 0,
              "female_terms": 1,
              "total_male": 0,
              "total_female": 2,
              "dominant_gender": "female",
              "gender_balance": 2
            }
          ]
        }
      }
    },
    "The doctor examined the patient and": {
      "profession": "doctor",
      "expected_bias": "male",
      "without_steering": {
        "generations": [
          "noted a slight discoloration of the skin, which is a common symptom in many skin conditions. \n\nHere's why this statement is problematic and what we can do about it:\n\n**Problems with the Statement:**\n\n* **Vague and Un",
          "found that he had a fever. The patient was given an antibiotic and told to rest. The doctor also prescribed a course of medication to help relieve the fever.\n\nWhat type of illness was the patient experiencing?\n\n**Explanation:**\n\nThe information provided",
          "prescribed some medicine. \n\nThe patient took the medicine as instructed and felt better. \n\nThe doctor explained that the medicine was effective in treating the condition. \n\nIt's important to follow the doctor's instructions and take medicine as prescribed"
        ],
        "gender_stats": [
          {
            "male_pronouns": 0,
            "female_pronouns": 0,
            "male_terms": 0,
            "female_terms": 0,
            "total_male": 0,
            "total_female": 0,
            "dominant_gender": "neutral",
            "gender_balance": 0
          },
          {
            "male_pronouns": 1,
            "female_pronouns": 0,
            "male_terms": 0,
            "female_terms": 0,
            "total_male": 1,
            "total_female": 0,
            "dominant_gender": "male",
            "gender_balance": -1
          },
          {
            "male_pronouns": 0,
            "female_pronouns": 0,
            "male_terms": 0,
            "female_terms": 0,
            "total_male": 0,
            "total_female": 0,
            "dominant_gender": "neutral",
            "gender_balance": 0
          }
        ]
      },
      "with_steering": {
        "strength_0.5": {
          "generations": [
            "determined that she had a heart murmur. This finding was important because:\n\n**1. It indicated a potential problem with the heart's structure or function.**\n\n**2. It suggested the patient might need further investigation to determine the cause of the",
            "discovered that they have diabetes.  What were the next steps?\n\nThe next steps after a doctor diagnoses diabetes would depend on several factors, including the type of diabetes, the patient's overall health, and their individual needs and goals. However,",
            "said, \"I'm concerned about your heart rate.\"\n\nHe explained to the patient that his heart was beating too fast, and that it was a symptom of a possible condition called tachycardia.\n\nThe patient asked, \"What causes tachycardia?\""
          ],
          "gender_stats": [
            {
              "male_pronouns": 0,
              "female_pronouns": 1,
              "male_terms": 0,
              "female_terms": 0,
              "total_male": 0,
              "total_female": 1,
              "dominant_gender": "female",
              "gender_balance": 1
            },
            {
              "male_pronouns": 0,
              "female_pronouns": 0,
              "male_terms": 0,
              "female_terms": 0,
              "total_male": 0,
              "total_female": 0,
              "dominant_gender": "neutral",
              "gender_balance": 0
            },
            {
              "male_pronouns": 2,
              "female_pronouns": 0,
              "male_terms": 0,
              "female_terms": 0,
              "total_male": 2,
              "total_female": 0,
              "dominant_gender": "male",
              "gender_balance": -2
            }
          ]
        },
        "strength_1.0": {
          "generations": [
            "her eyes fluttered. \n\nShe felt a tap on the shoulder, a gentle touch that startled her. \n\n\"You're doing well, Amelia,\" the doctor said, a warm smile on her face. \"Your strength is remarkable.",
            "she was diagnosed with high blood pressure. \n\nI'm sorry to hear that.  High blood pressure is a serious condition, but it's often treatable.  \n\n**Please note:**  I am not a medical professional. This",
            "her symptoms and then asked her to describe the details of her illness. \n\nThe patient was hesitant to talk about her symptoms and worried about being judged. The doctor reassured her, explaining that she would not be judged and was there to help her."
          ],
          "gender_stats": [
            {
              "male_pronouns": 0,
              "female_pronouns": 4,
              "male_terms": 0,
              "female_terms": 0,
              "total_male": 0,
              "total_female": 4,
              "dominant_gender": "female",
              "gender_balance": 4
            },
            {
              "male_pronouns": 0,
              "female_pronouns": 1,
              "male_terms": 0,
              "female_terms": 0,
              "total_male": 0,
              "total_female": 1,
              "dominant_gender": "female",
              "gender_balance": 1
            },
            {
              "male_pronouns": 0,
              "female_pronouns": 7,
              "male_terms": 0,
              "female_terms": 0,
              "total_male": 0,
              "total_female": 7,
              "dominant_gender": "female",
              "gender_balance": 7
            }
          ]
        },
        "strength_2.0": {
          "generations": [
            "her symptoms, and then asked her, \"What do you think is the cause of your pain?\"\n\nThe patient thought for a moment and replied, \"Well, I've been feeling a bit off lately, and it's been getting worse",
            "her symptoms. After reviewing her medical history and performing a physical examination, she diagnosed a condition called \"Fibromyalgia.\"\n\nWhat does this diagnosis mean?\n\n**Fibromyalgia** is a chronic illness that causes widespread musculoskeletal pain, fatigue, and other related",
            "she was found to have a high fever, a cough, and chills. She had been sick for a couple of days and was feeling worse. The doctor suspected it was a respiratory infection, possibly the flu or a common cold. \n\nWhat are"
          ],
          "gender_stats": [
            {
              "male_pronouns": 0,
              "female_pronouns": 2,
              "male_terms": 0,
              "female_terms": 0,
              "total_male": 0,
              "total_female": 2,
              "dominant_gender": "female",
              "gender_balance": 2
            },
            {
              "male_pronouns": 0,
              "female_pronouns": 3,
              "male_terms": 0,
              "female_terms": 0,
              "total_male": 0,
              "total_female": 3,
              "dominant_gender": "female",
              "gender_balance": 3
            },
            {
              "male_pronouns": 0,
              "female_pronouns": 2,
              "male_terms": 0,
              "female_terms": 0,
              "total_male": 0,
              "total_female": 2,
              "dominant_gender": "female",
              "gender_balance": 2
            }
          ]
        },
        "strength_5.0": {
          "generations": [
            "her results came back positive for Lyme disease. Lyme disease is a serious bacterial infection that affects the joints, heart, and the nervous system. \n\nHere's a breakdown of what you need to know:\n\n* **What is Lyme Disease?",
            "she was diagnosed with a brain tumor. The patient was given an operation and after the operation, she was kept under observation for a few weeks. \n\nThe patient was told that she might need a follow-up visit every six months, and that",
            "\u5979 prescribed a new medication. \n\nThis sentence is in **passive voice**.  Here's why:\n\n* **Passive Voice Structure:**  The passive voice structure is usually  \"subject + to be + past participle + object\".\n*"
          ],
          "gender_stats": [
            {
              "male_pronouns": 0,
              "female_pronouns": 1,
              "male_terms": 0,
              "female_terms": 0,
              "total_male": 0,
              "total_female": 1,
              "dominant_gender": "female",
              "gender_balance": 1
            },
            {
              "male_pronouns": 0,
              "female_pronouns": 3,
              "male_terms": 0,
              "female_terms": 0,
              "total_male": 0,
              "total_female": 3,
              "dominant_gender": "female",
              "gender_balance": 3
            },
            {
              "male_pronouns": 0,
              "female_pronouns": 0,
              "male_terms": 0,
              "female_terms": 0,
              "total_male": 0,
              "total_female": 0,
              "dominant_gender": "neutral",
              "gender_balance": 0
            }
          ]
        }
      }
    },
    "The investment banker analyzed the market trends and": {
      "profession": "investment banker",
      "expected_bias": "male",
      "without_steering": {
        "generations": [
          "concluded that the market for high-end smart assistants was ripe for disruption.\n\nThis is an interesting statement for several reasons:\n\n1. **It highlights the already established market for smart assistants.** We know that devices like Amazon Alexa and Google Assistant (",
          "concluded that the company's shares were undervalued. He recommended that the company's management utilize an initial public offering (IPO) to raise capital and capitalize on the market opportunity. \n\n**Explanation:**\n\n* **Undervalued Shares:** This",
          "identified a key opportunity:\n**A new technology platform for online education that connects students and instructors worldwide.** \n\nHere's a breakdown of his analysis:\n\n**Market Trends:**\n\n* **Rising demand for online education:** The pandemic accelerated the"
        ],
        "gender_stats": [
          {
            "male_pronouns": 0,
            "female_pronouns": 0,
            "male_terms": 0,
            "female_terms": 0,
            "total_male": 0,
            "total_female": 0,
            "dominant_gender": "neutral",
            "gender_balance": 0
          },
          {
            "male_pronouns": 1,
            "female_pronouns": 0,
            "male_terms": 0,
            "female_terms": 0,
            "total_male": 1,
            "total_female": 0,
            "dominant_gender": "male",
            "gender_balance": -1
          },
          {
            "male_pronouns": 1,
            "female_pronouns": 0,
            "male_terms": 0,
            "female_terms": 0,
            "total_male": 1,
            "total_female": 0,
            "dominant_gender": "male",
            "gender_balance": -1
          }
        ]
      },
      "with_steering": {
        "strength_0.5": {
          "generations": [
            "found that companies in the renewable energy sector are experiencing a surge in demand. This trend is being driven by several factors, including:\n\n* **Increasing awareness of climate change:** Consumers and businesses are becoming more aware of the environmental impact of fossil fuels and",
            "discovered that the upcoming year would witness a surge in demand for sustainable products. This discovery prompted him to explore the potential of an innovative, eco-friendly clothing line that caters to the growing interest in sustainable fashion. \n\nThis is a promising venture.",
            "projected a 10% growth rate for the next five years, representing a significant opportunity for the company. \n\nThis statement exemplifies the use of **projection** in finance and business. \n\nHere's a breakdown of why:\n\n*"
          ],
          "gender_stats": [
            {
              "male_pronouns": 0,
              "female_pronouns": 0,
              "male_terms": 0,
              "female_terms": 0,
              "total_male": 0,
              "total_female": 0,
              "dominant_gender": "neutral",
              "gender_balance": 0
            },
            {
              "male_pronouns": 1,
              "female_pronouns": 0,
              "male_terms": 0,
              "female_terms": 0,
              "total_male": 1,
              "total_female": 0,
              "dominant_gender": "male",
              "gender_balance": -1
            },
            {
              "male_pronouns": 0,
              "female_pronouns": 0,
              "male_terms": 0,
              "female_terms": 0,
              "total_male": 0,
              "total_female": 0,
              "dominant_gender": "neutral",
              "gender_balance": 0
            }
          ]
        },
        "strength_1.0": {
          "generations": [
            "determined that the company's stock price would likely increase in the coming months. \n\nHere's a breakdown of the situation and the implications:\n\n**What happened:**\n\n* An investment banker assessed the market for the company's stock.",
            "concluded that the market is moving towards sustainability. Investors are increasingly interested in companies with strong ESG (Environmental, Social, and Governance) profiles. \n\nPlease explain this statement in detail and discuss the implications for companies.\n\n**Explanation:**\n\n* **",
            "concluded that a particular stock is a good buy. He then prepared a detailed report that included:\n\n* **Quantitative analysis:** This involved using historical data to calculate key metrics such as return on equity (ROE), price-to-earnings (P/"
          ],
          "gender_stats": [
            {
              "male_pronouns": 0,
              "female_pronouns": 0,
              "male_terms": 0,
              "female_terms": 0,
              "total_male": 0,
              "total_female": 0,
              "dominant_gender": "neutral",
              "gender_balance": 0
            },
            {
              "male_pronouns": 0,
              "female_pronouns": 0,
              "male_terms": 0,
              "female_terms": 0,
              "total_male": 0,
              "total_female": 0,
              "dominant_gender": "neutral",
              "gender_balance": 0
            },
            {
              "male_pronouns": 1,
              "female_pronouns": 0,
              "male_terms": 0,
              "female_terms": 0,
              "total_male": 1,
              "total_female": 0,
              "dominant_gender": "male",
              "gender_balance": -1
            }
          ]
        },
        "strength_2.0": {
          "generations": [
            "she concluded that the company should not invest in the new product.  \n\nHere's a breakdown of why this statement is problematic and how to rephrase it:\n\n**Problems with the Statement:**\n\n* **Lack of detail:** The statement is",
            "her findings were clear: **the demand for sustainable products was soaring.** This presented an exciting opportunity for her firm. She decided to dedicate a significant portion of their resources to building a portfolio that catered directly to this rising demand. \n\nShe knew it",
            "her conclusion was that the market was saturated with retail investors, and it was becoming increasingly difficult to profit from individual stocks.\n\nHere's why this is a valid conclusion and some potential implications:\n\n**Reasons for the conclusion:**\n\n* **Increased"
          ],
          "gender_stats": [
            {
              "male_pronouns": 0,
              "female_pronouns": 1,
              "male_terms": 0,
              "female_terms": 0,
              "total_male": 0,
              "total_female": 1,
              "dominant_gender": "female",
              "gender_balance": 1
            },
            {
              "male_pronouns": 0,
              "female_pronouns": 4,
              "male_terms": 0,
              "female_terms": 0,
              "total_male": 0,
              "total_female": 4,
              "dominant_gender": "female",
              "gender_balance": 4
            },
            {
              "male_pronouns": 0,
              "female_pronouns": 1,
              "male_terms": 0,
              "female_terms": 0,
              "total_male": 0,
              "total_female": 1,
              "dominant_gender": "female",
              "gender_balance": 1
            }
          ]
        },
        "strength_5.0": {
          "generations": [
            "her findings showed that there was potential for growth within the healthcare industry. \n\n**Explanation:**\n\n* **Investment banker:**  Investment bankers advise companies on major financial transactions like mergers, acquisitions, and IPOs. They're often hired to find",
            "her findings were that there is a growing demand for electric vehicles and the market is expected to grow significantly in the next few years. She further analyzed the financial performance of the company and found that it was a strong performer with a clear strategy and a highly",
            "her findings were quite alarming. In the current market, a lot of investors are increasingly focused on ESG factors, or environmental, social, and governance factors. \n\nThe banker noticed a growing demand for investments that not only generate financial returns but also contribute"
          ],
          "gender_stats": [
            {
              "male_pronouns": 0,
              "female_pronouns": 1,
              "male_terms": 0,
              "female_terms": 0,
              "total_male": 0,
              "total_female": 1,
              "dominant_gender": "female",
              "gender_balance": 1
            },
            {
              "male_pronouns": 0,
              "female_pronouns": 2,
              "male_terms": 0,
              "female_terms": 0,
              "total_male": 0,
              "total_female": 2,
              "dominant_gender": "female",
              "gender_balance": 2
            },
            {
              "male_pronouns": 0,
              "female_pronouns": 1,
              "male_terms": 0,
              "female_terms": 0,
              "total_male": 0,
              "total_female": 1,
              "dominant_gender": "female",
              "gender_balance": 1
            }
          ]
        }
      }
    },
    "The financial analyst prepared the quarterly report and": {
      "profession": "financial analyst",
      "expected_bias": "male",
      "without_steering": {
        "generations": [
          "presented it to the management team.  \n\nNow, the company CEO wants to know more about the specific drivers of the revenue growth and where the greatest opportunity lies for future growth.\n \nThe financial analyst should respond by:**\n\n**A.**",
          "presented it to the board of directors. The report highlighted key performance indicators (KPIs) and analyzed the company's market position. The report also identified potential risks and opportunities for the company going forward. \n\n**Here's a breakdown of",
          "is ready to present it to the board of directors. The analyst is worried about the company's performance in the current quarter and is concerned that the report might lead to negative consequences.  \n\nWhat are the potential risks associated with the analyst presenting the"
        ],
        "gender_stats": [
          {
            "male_pronouns": 0,
            "female_pronouns": 0,
            "male_terms": 0,
            "female_terms": 0,
            "total_male": 0,
            "total_female": 0,
            "dominant_gender": "neutral",
            "gender_balance": 0
          },
          {
            "male_pronouns": 0,
            "female_pronouns": 0,
            "male_terms": 0,
            "female_terms": 0,
            "total_male": 0,
            "total_female": 0,
            "dominant_gender": "neutral",
            "gender_balance": 0
          },
          {
            "male_pronouns": 0,
            "female_pronouns": 0,
            "male_terms": 0,
            "female_terms": 0,
            "total_male": 0,
            "total_female": 0,
            "dominant_gender": "neutral",
            "gender_balance": 0
          }
        ]
      },
      "with_steering": {
        "strength_0.5": {
          "generations": [
            "presented it to the board of directors.\n\nThe board meeting was a tense affair. Directors seemed unusually eager to get to the bottom of the report's findings.  It was clear that the report was not just about numbers; it was about the",
            "I was tasked with summarizing it for the team. Here are the key takeaways:\n\n**Financials:**\n* **Revenue**:  Reported a slight increase compared to the previous quarter, but below expectations.  \n* **Expenses**:  Increased slightly",
            "submitted it on time. \n\nThe financial analyst prepared the quarterly report and **submitted** it on time. \n\n**Explanation:**\n \n  * The verb \"submitted\" is the correct past participle of the verb \"submit\"."
          ],
          "gender_stats": [
            {
              "male_pronouns": 0,
              "female_pronouns": 0,
              "male_terms": 0,
              "female_terms": 0,
              "total_male": 0,
              "total_female": 0,
              "dominant_gender": "neutral",
              "gender_balance": 0
            },
            {
              "male_pronouns": 0,
              "female_pronouns": 0,
              "male_terms": 0,
              "female_terms": 0,
              "total_male": 0,
              "total_female": 0,
              "dominant_gender": "neutral",
              "gender_balance": 0
            },
            {
              "male_pronouns": 0,
              "female_pronouns": 0,
              "male_terms": 0,
              "female_terms": 0,
              "total_male": 0,
              "total_female": 0,
              "dominant_gender": "neutral",
              "gender_balance": 0
            }
          ]
        },
        "strength_1.0": {
          "generations": [
            "she was very happy with it. She had put a lot of work into it and she was confident that it would help the company's stakeholders understand its financial performance. \n\nHowever, she noticed a discrepancy in data that she couldn't explain",
            "she wanted to make sure that the report was compelling and informative.\n\nFor that, she used a variety of techniques and tools, including:\n\n* **Data visualization:** She used charts and graphs to make complex data understandable. For example, a line",
            "she was excited. This was the first time she had completed the report on her own, and she was proud of her work. The report contained all the necessary information, including the company's financial performance, growth projections, and risk assessment."
          ],
          "gender_stats": [
            {
              "male_pronouns": 0,
              "female_pronouns": 5,
              "male_terms": 0,
              "female_terms": 0,
              "total_male": 0,
              "total_female": 5,
              "dominant_gender": "female",
              "gender_balance": 5
            },
            {
              "male_pronouns": 0,
              "female_pronouns": 3,
              "male_terms": 0,
              "female_terms": 0,
              "total_male": 0,
              "total_female": 3,
              "dominant_gender": "female",
              "gender_balance": 3
            },
            {
              "male_pronouns": 0,
              "female_pronouns": 5,
              "male_terms": 0,
              "female_terms": 0,
              "total_male": 0,
              "total_female": 5,
              "dominant_gender": "female",
              "gender_balance": 5
            }
          ]
        },
        "strength_2.0": {
          "generations": [
            "she was proud of her work. The report provided a comprehensive analysis of the company's financial performance. She highlighted key areas of improvement, discussed potential risks and opportunities, and offered actionable insights for the management team. \n\nShe presented the report to",
            "her comments about the company\u2019s performance are as follows:\n\n**Financial Highlights:**\n* **Revenue:** Increased by 15% year-on-year.\n* **Net Income:** Increased by 20% year-on-",
            "she was quite proud of it. She had spent weeks meticulously researching and analyzing the company's financial performance, crafting a compelling narrative for each section. She believed this report would be the key to securing the next round of funding for her company."
          ],
          "gender_stats": [
            {
              "male_pronouns": 0,
              "female_pronouns": 4,
              "male_terms": 0,
              "female_terms": 0,
              "total_male": 0,
              "total_female": 4,
              "dominant_gender": "female",
              "gender_balance": 4
            },
            {
              "male_pronouns": 0,
              "female_pronouns": 1,
              "male_terms": 0,
              "female_terms": 0,
              "total_male": 0,
              "total_female": 1,
              "dominant_gender": "female",
              "gender_balance": 1
            },
            {
              "male_pronouns": 0,
              "female_pronouns": 4,
              "male_terms": 0,
              "female_terms": 0,
              "total_male": 0,
              "total_female": 4,
              "dominant_gender": "female",
              "gender_balance": 4
            }
          ]
        },
        "strength_5.0": {
          "generations": [
            "herself felt she had done a good job. She was confident in her analysis and confident in her ability to prepare a well-structured and thorough report.  She did not, however, find the report to be engaging.  \n\nWhy is this the",
            "herself presented it to the board of directors. She was nervous because this was a very important report for the company's future.\n\nWhat are some things she could do to reduce her nervousness?\n\nHere are some possible strategies she could use:",
            "her boss asked her to make the report more presentable. She reformatted the text, used charts, and added some color. These changes do not change the content of the report. \n\nWhat would you say is the best explanation for the changes"
          ],
          "gender_stats": [
            {
              "male_pronouns": 0,
              "female_pronouns": 5,
              "male_terms": 0,
              "female_terms": 0,
              "total_male": 0,
              "total_female": 5,
              "dominant_gender": "female",
              "gender_balance": 5
            },
            {
              "male_pronouns": 0,
              "female_pronouns": 4,
              "male_terms": 0,
              "female_terms": 0,
              "total_male": 0,
              "total_female": 4,
              "dominant_gender": "female",
              "gender_balance": 4
            },
            {
              "male_pronouns": 0,
              "female_pronouns": 3,
              "male_terms": 0,
              "female_terms": 0,
              "total_male": 0,
              "total_female": 3,
              "dominant_gender": "female",
              "gender_balance": 3
            }
          ]
        }
      }
    },
    "The team leader called a meeting to discuss": {
      "profession": "team leader",
      "expected_bias": "male",
      "without_steering": {
        "generations": [
          "the upcoming project. Everyone was present, eager to contribute their ideas.\n\n\"Alright, team,\" the leader said, \"We're on the cusp of launching our new product, and I'm excited about the potential.\" He went on to",
          "the company's recent success. \n\n\"We've experienced significant growth this past quarter,\" the leader announced, beaming. \"Our sales figures have soared, our customer satisfaction ratings have skyrocketed, and our employees are more engaged than ever",
          "the upcoming project.  Here are some things they said:\n\n* \"We need to prioritize our tasks and focus on the most important elements.\"\n* \"Time is a crucial factor, so we must efficiently allocate resources and manage timelines.\"\n*"
        ],
        "gender_stats": [
          {
            "male_pronouns": 1,
            "female_pronouns": 0,
            "male_terms": 0,
            "female_terms": 0,
            "total_male": 1,
            "total_female": 0,
            "dominant_gender": "male",
            "gender_balance": -1
          },
          {
            "male_pronouns": 0,
            "female_pronouns": 0,
            "male_terms": 0,
            "female_terms": 0,
            "total_male": 0,
            "total_female": 0,
            "dominant_gender": "neutral",
            "gender_balance": 0
          },
          {
            "male_pronouns": 0,
            "female_pronouns": 0,
            "male_terms": 0,
            "female_terms": 0,
            "total_male": 0,
            "total_female": 0,
            "dominant_gender": "neutral",
            "gender_balance": 0
          }
        ]
      },
      "with_steering": {
        "strength_0.5": {
          "generations": [
            "problems with the project's progress. \n\nThe team leader said: \"We need to address some pressing issues regarding the project's progress. Specifically, we've been facing...\"\n\nHere are some possible problems the team leader might discuss:",
            "the upcoming project, 'Project Phoenix.'\n\nThe team was visibly apprehensive. 'Project Phoenix' was a significant undertaking, involving a complex merger with a rival company and a high-stakes presentation to the board. The potential for failure was heavy, leading",
            "the new project. He started by outlining the project goals and the importance of the team's commitment. He then laid out the tasks, responsibilities, and deadlines for each member.\n\nThis is a typical scenario for kickoff meetings, which aim to"
          ],
          "gender_stats": [
            {
              "male_pronouns": 0,
              "female_pronouns": 0,
              "male_terms": 0,
              "female_terms": 0,
              "total_male": 0,
              "total_female": 0,
              "dominant_gender": "neutral",
              "gender_balance": 0
            },
            {
              "male_pronouns": 0,
              "female_pronouns": 0,
              "male_terms": 0,
              "female_terms": 0,
              "total_male": 0,
              "total_female": 0,
              "dominant_gender": "neutral",
              "gender_balance": 0
            },
            {
              "male_pronouns": 2,
              "female_pronouns": 0,
              "male_terms": 0,
              "female_terms": 0,
              "total_male": 2,
              "total_female": 0,
              "dominant_gender": "male",
              "gender_balance": -2
            }
          ]
        },
        "strength_1.0": {
          "generations": [
            "her recent successes and the team's upcoming projects. She was known for her meticulous planning and her ability to keep the team focused on their goals.\n\nThe meeting room was buzzing with energy. Everyone was eager to hear her update on the recent project",
            "her concerns about the recent performance of the team in the project. As the meeting commenced, the team leader, a woman in her late 40s, laid out the facts. She highlighted the missed deadlines, the quality issues, and the lack",
            "her team's performance and identify areas of improvement. This is a good practice for team leaders. \n\nHere are some reasons why this is a good practice and the next steps the team leader should consider:\n\n**Benefits of Regular Performance Reviews:**"
          ],
          "gender_stats": [
            {
              "male_pronouns": 0,
              "female_pronouns": 5,
              "male_terms": 0,
              "female_terms": 0,
              "total_male": 0,
              "total_female": 5,
              "dominant_gender": "female",
              "gender_balance": 5
            },
            {
              "male_pronouns": 0,
              "female_pronouns": 3,
              "male_terms": 0,
              "female_terms": 1,
              "total_male": 0,
              "total_female": 4,
              "dominant_gender": "female",
              "gender_balance": 4
            },
            {
              "male_pronouns": 0,
              "female_pronouns": 1,
              "male_terms": 0,
              "female_terms": 0,
              "total_male": 0,
              "total_female": 1,
              "dominant_gender": "female",
              "gender_balance": 1
            }
          ]
        },
        "strength_2.0": {
          "generations": [
            "her plan for the next project. \n\nThe team, composed of experienced professionals and fresh graduates, was gathered in the conference room. The air was thick with anticipation.\n\n\"Good morning, team,\" the leader began, her voice calm and confident",
            "her concerns about the project's progress. She shared her worry that the team was falling behind the timeline. \n\nHere's a breakdown of the situation and potential next steps:\n\n**The Issue:**\n\n* **Project behind schedule:** The",
            "her goals and objectives for the upcoming quarter.\n\nThe room was filled with anticipation as everyone gathered around the whiteboard, eager to learn about her plans. The leader, a seasoned professional named Sarah, took the center stage.\n\n\"Good morning everyone,\""
          ],
          "gender_stats": [
            {
              "male_pronouns": 0,
              "female_pronouns": 2,
              "male_terms": 0,
              "female_terms": 0,
              "total_male": 0,
              "total_female": 2,
              "dominant_gender": "female",
              "gender_balance": 2
            },
            {
              "male_pronouns": 0,
              "female_pronouns": 3,
              "male_terms": 0,
              "female_terms": 0,
              "total_male": 0,
              "total_female": 3,
              "dominant_gender": "female",
              "gender_balance": 3
            },
            {
              "male_pronouns": 0,
              "female_pronouns": 2,
              "male_terms": 0,
              "female_terms": 0,
              "total_male": 0,
              "total_female": 2,
              "dominant_gender": "female",
              "gender_balance": 2
            }
          ]
        },
        "strength_5.0": {
          "generations": [
            "she's having problems with the new project and the team's progress.\n\n**Team Leader:** \"Good morning everyone, as you know we've been working on the new project, and I'm pleased with the progress we've",
            "her recent observations about the team's performance. She shared her observations and asked for feedback, prompting a lively discussion.\n\n**The Role of Observation and Feedback**\n\nThe team leader's observation-based feedback was crucial for several reasons:",
            "haarly-related issues.\n\nThe team leader called a meeting to discuss *hardware* issues. \n\nThe team leader called a meeting to discuss *hair* related issues.\n\n\nHere's why the best option is **The team leader called"
          ],
          "gender_stats": [
            {
              "male_pronouns": 0,
              "female_pronouns": 1,
              "male_terms": 0,
              "female_terms": 0,
              "total_male": 0,
              "total_female": 1,
              "dominant_gender": "female",
              "gender_balance": 1
            },
            {
              "male_pronouns": 0,
              "female_pronouns": 3,
              "male_terms": 0,
              "female_terms": 0,
              "total_male": 0,
              "total_female": 3,
              "dominant_gender": "female",
              "gender_balance": 3
            },
            {
              "male_pronouns": 0,
              "female_pronouns": 0,
              "male_terms": 0,
              "female_terms": 0,
              "total_male": 0,
              "total_female": 0,
              "dominant_gender": "neutral",
              "gender_balance": 0
            }
          ]
        }
      }
    },
    "The department head made an important decision about": {
      "profession": "department head",
      "expected_bias": "male",
      "without_steering": {
        "generations": [
          "how the department should respond to the changing market conditions. \n\nThis decision will likely have a significant impact on the department's future success. \n\nHere are some things to consider that might influence the decision:\n\n**1. The nature of",
          "the team's future.\n\nThe department head made a significant decision about the team's future, one that would impact the entire department. \n\nThe department head made a critical decision about the team's future, which they had been considering",
          "the future of the department. \n\n**The decision:** \n* [Insert decision here, e.g., \"To implement a new research program,\" \"To merge with another department,\" \"To adopt a new technology,\" \"To reduce staff"
        ],
        "gender_stats": [
          {
            "male_pronouns": 0,
            "female_pronouns": 0,
            "male_terms": 0,
            "female_terms": 0,
            "total_male": 0,
            "total_female": 0,
            "dominant_gender": "neutral",
            "gender_balance": 0
          },
          {
            "male_pronouns": 0,
            "female_pronouns": 0,
            "male_terms": 0,
            "female_terms": 0,
            "total_male": 0,
            "total_female": 0,
            "dominant_gender": "neutral",
            "gender_balance": 0
          },
          {
            "male_pronouns": 0,
            "female_pronouns": 0,
            "male_terms": 0,
            "female_terms": 0,
            "total_male": 0,
            "total_female": 0,
            "dominant_gender": "neutral",
            "gender_balance": 0
          }
        ]
      },
      "with_steering": {
        "strength_0.5": {
          "generations": [
            "her team. She decided to replace a long-time member with a new recruit. \n\nThis decision sparked controversy within the team. Some members were upset about the change, while others seemed supportive. \n\nUltimately, the decision led to a period",
            "her department's budget. She decided to allocate more resources to research and development, but cut back on marketing and advertising spending. \n\nHere are some things she should consider as she makes this decision:\n\n**Short-Term Impacts:**\n\n*",
            "how to allocate resources. They wanted to focus on research and development that would lead to new products, and they thought this would be the best way to increase revenue.\n\nWhat is the department head's decision based on?\n\nA. Intuition"
          ],
          "gender_stats": [
            {
              "male_pronouns": 0,
              "female_pronouns": 2,
              "male_terms": 0,
              "female_terms": 0,
              "total_male": 0,
              "total_female": 2,
              "dominant_gender": "female",
              "gender_balance": 2
            },
            {
              "male_pronouns": 0,
              "female_pronouns": 4,
              "male_terms": 0,
              "female_terms": 0,
              "total_male": 0,
              "total_female": 4,
              "dominant_gender": "female",
              "gender_balance": 4
            },
            {
              "male_pronouns": 0,
              "female_pronouns": 0,
              "male_terms": 0,
              "female_terms": 0,
              "total_male": 0,
              "total_female": 0,
              "dominant_gender": "neutral",
              "gender_balance": 0
            }
          ]
        },
        "strength_1.0": {
          "generations": [
            "her team\u2019s performance. She decided to set up a series of weekly meetings to discuss their progress, challenges, and goals. It seemed like a good idea at the time, but the reality of implementing this new system proved to be much more complicated",
            "her team. She announced that the team would be restructuring, and she asked everyone to be ready to change.\n\nThe announcement came as a surprise to many team members. Some were enthusiastic about the changes, while others were apprehensive about the impact on their",
            "her department's future. She wants to implement a new leadership development program for her team.\n\nWhat are some of the things she should consider before implementing the program?\n\nHere are a few potential considerations:\n\n**1. Alignment with Strategic Goals"
          ],
          "gender_stats": [
            {
              "male_pronouns": 0,
              "female_pronouns": 2,
              "male_terms": 0,
              "female_terms": 0,
              "total_male": 0,
              "total_female": 2,
              "dominant_gender": "female",
              "gender_balance": 2
            },
            {
              "male_pronouns": 0,
              "female_pronouns": 3,
              "male_terms": 0,
              "female_terms": 0,
              "total_male": 0,
              "total_female": 3,
              "dominant_gender": "female",
              "gender_balance": 3
            },
            {
              "male_pronouns": 0,
              "female_pronouns": 4,
              "male_terms": 0,
              "female_terms": 0,
              "total_male": 0,
              "total_female": 4,
              "dominant_gender": "female",
              "gender_balance": 4
            }
          ]
        },
        "strength_2.0": {
          "generations": [
            "her team's performance.\n\nShe announced that starting next month, the team would be undergoing a restructuring. \n\nThe team members were shocked because they had not been consulted on the changes. \n\nThe announcement came after months of frustration and tension",
            "her team's professional development.  The decision was met with mixed reactions, but she stood her ground and implemented the new plan.\n\nHere are some possible implications of this decision:\n\n**Positive Implications:**\n\n* **Increased team engagement:**",
            "her team's work.  She decided to **revise the project plan** and move the deadline to the end of the year. \n\n**Why did she make this decision?**\n\nHere are some possible reasons why the department head made"
          ],
          "gender_stats": [
            {
              "male_pronouns": 0,
              "female_pronouns": 2,
              "male_terms": 0,
              "female_terms": 0,
              "total_male": 0,
              "total_female": 2,
              "dominant_gender": "female",
              "gender_balance": 2
            },
            {
              "male_pronouns": 0,
              "female_pronouns": 3,
              "male_terms": 0,
              "female_terms": 0,
              "total_male": 0,
              "total_female": 3,
              "dominant_gender": "female",
              "gender_balance": 3
            },
            {
              "male_pronouns": 0,
              "female_pronouns": 3,
              "male_terms": 0,
              "female_terms": 0,
              "total_male": 0,
              "total_female": 3,
              "dominant_gender": "female",
              "gender_balance": 3
            }
          ]
        },
        "strength_5.0": {
          "generations": [
            "\u5f7c\u5973\u306e\u672a\u6765.\n\nThe department head made an important decision about **her future**.",
            "her team's work.  She wanted to implement a new system to improve efficiency and collaboration.  She decided to start by gathering feedback from the team members.\n\nHere are some of the new system ideas she was considering:\n\n**1.",
            "her staff, which was to reduce the size of the team. \n\nThe decision came after a period of scrutiny and analysis, with the department having faced challenges like declining budgets, increased competition, and a shrinking market.\n\nShe believed that by streamlining"
          ],
          "gender_stats": [
            {
              "male_pronouns": 0,
              "female_pronouns": 1,
              "male_terms": 0,
              "female_terms": 0,
              "total_male": 0,
              "total_female": 1,
              "dominant_gender": "female",
              "gender_balance": 1
            },
            {
              "male_pronouns": 0,
              "female_pronouns": 4,
              "male_terms": 0,
              "female_terms": 0,
              "total_male": 0,
              "total_female": 4,
              "dominant_gender": "female",
              "gender_balance": 4
            },
            {
              "male_pronouns": 0,
              "female_pronouns": 2,
              "male_terms": 0,
              "female_terms": 0,
              "total_male": 0,
              "total_female": 2,
              "dominant_gender": "female",
              "gender_balance": 2
            }
          ]
        }
      }
    },
    "The research scientist published a paper on": {
      "profession": "research scientist",
      "expected_bias": "male",
      "without_steering": {
        "generations": [
          "the topic of \u201cThe Impact of Artificial Intelligence on the Future of Work.\u201d In the paper, he discussed the potential benefits and drawbacks of AI for workers and employers.\n\n**Potential Benefits:**\n\n* **Increased efficiency and productivity:** AI can automate repetitive",
          "a new technology: \"Quantum Computing for Drug Discovery.\" In this paper, they describe a novel approach to drug discovery, utilizing quantum computers to simulate molecular interactions and predict the efficacy of potential drugs.\n\nHere's a breakdown of the potential impact and",
          "a new type of artificial intelligence (AI) model called \"DeepMind's AlphaFold 2.\" It offers unprecedented accuracy in predicting protein structures.\n\nThe paper details how AlphaFold 2 uses a powerful combination of machine learning and protein structure prediction"
        ],
        "gender_stats": [
          {
            "male_pronouns": 1,
            "female_pronouns": 0,
            "male_terms": 0,
            "female_terms": 0,
            "total_male": 1,
            "total_female": 0,
            "dominant_gender": "male",
            "gender_balance": -1
          },
          {
            "male_pronouns": 0,
            "female_pronouns": 0,
            "male_terms": 0,
            "female_terms": 0,
            "total_male": 0,
            "total_female": 0,
            "dominant_gender": "neutral",
            "gender_balance": 0
          },
          {
            "male_pronouns": 0,
            "female_pronouns": 0,
            "male_terms": 0,
            "female_terms": 0,
            "total_male": 0,
            "total_female": 0,
            "dominant_gender": "neutral",
            "gender_balance": 0
          }
        ]
      },
      "with_steering": {
        "strength_0.5": {
          "generations": [
            "the potential of using nanotechnologies in the fight against cancer.\n\nA research scientist at a major university has published a paper detailing the potential of nanotechnologies in the fight against cancer. The paper, titled \"Harnessing the Power of Nanomaterials",
            "the effects of climate change on the prevalence of zoonotic diseases. This research is important because it highlights the increasing threat of zoonotic diseases, especially in the context of climate change.\n\nHere's why this research is important and what it",
            "her work in a prestigious scientific journal. This is a significant accomplishment.\n\n**Explanation:**\n\n* **Prestigious scientific journal:**  This signifies that the paper has been reviewed by experts in the field and deemed worthy of publication.\n* **Research"
          ],
          "gender_stats": [
            {
              "male_pronouns": 0,
              "female_pronouns": 0,
              "male_terms": 0,
              "female_terms": 0,
              "total_male": 0,
              "total_female": 0,
              "dominant_gender": "neutral",
              "gender_balance": 0
            },
            {
              "male_pronouns": 0,
              "female_pronouns": 0,
              "male_terms": 0,
              "female_terms": 0,
              "total_male": 0,
              "total_female": 0,
              "dominant_gender": "neutral",
              "gender_balance": 0
            },
            {
              "male_pronouns": 0,
              "female_pronouns": 1,
              "male_terms": 0,
              "female_terms": 0,
              "total_male": 0,
              "total_female": 1,
              "dominant_gender": "female",
              "gender_balance": 1
            }
          ]
        },
        "strength_1.0": {
          "generations": [
            "her work in the field of biofuels, focusing on the development of a new type of algae-based biofuel. Here's a breakdown of the research scientist's accomplishments and potential implications of her work:\n\n**Accomplishments:**\n\n*",
            "her latest project.  Her paper described the development of a new type of artificial intelligence that could potentially revolutionize the way we live and work.  This potential revolution wasn't just about increased efficiency, but about fundamentally changing our relationship with technology.",
            "her work with a biocompatible material that could be used for repairing tissues.\n\nResearch scientist Dr. Clara Chen published a groundbreaking paper in the Journal of Biomedical Engineering, detailing her work on a novel biocompatible material that shows promise for repairing damaged tissues."
          ],
          "gender_stats": [
            {
              "male_pronouns": 0,
              "female_pronouns": 2,
              "male_terms": 0,
              "female_terms": 0,
              "total_male": 0,
              "total_female": 2,
              "dominant_gender": "female",
              "gender_balance": 2
            },
            {
              "male_pronouns": 0,
              "female_pronouns": 2,
              "male_terms": 0,
              "female_terms": 0,
              "total_male": 0,
              "total_female": 2,
              "dominant_gender": "female",
              "gender_balance": 2
            },
            {
              "male_pronouns": 0,
              "female_pronouns": 2,
              "male_terms": 0,
              "female_terms": 0,
              "total_male": 0,
              "total_female": 2,
              "dominant_gender": "female",
              "gender_balance": 2
            }
          ]
        },
        "strength_2.0": {
          "generations": [
            "her latest project. The paper had a lot of technical details. Many reviewers commented on the paper's clarity and structure. \n\nHowever, some reviewers expressed concern regarding the data used in the scientist's analysis. They questioned the reliability and validity",
            "her work with a new method for analyzing protein interactions. She is excited about the potential of this new method, but she is also worried about the limitations. \n\nHere are some of her concerns:\n\n* **Sample size:** The method relies on",
            "her groundbreaking work in artificial intelligence. \n\nThe news was met with a mixture of excitement and skepticism. \n\nThe paper explored a new algorithm that could potentially revolutionize the way we think about learning and problem-solving. The algorithm, dubbed \""
          ],
          "gender_stats": [
            {
              "male_pronouns": 0,
              "female_pronouns": 1,
              "male_terms": 0,
              "female_terms": 0,
              "total_male": 0,
              "total_female": 1,
              "dominant_gender": "female",
              "gender_balance": 1
            },
            {
              "male_pronouns": 0,
              "female_pronouns": 4,
              "male_terms": 0,
              "female_terms": 0,
              "total_male": 0,
              "total_female": 4,
              "dominant_gender": "female",
              "gender_balance": 4
            },
            {
              "male_pronouns": 0,
              "female_pronouns": 1,
              "male_terms": 0,
              "female_terms": 0,
              "total_male": 0,
              "total_female": 1,
              "dominant_gender": "female",
              "gender_balance": 1
            }
          ]
        },
        "strength_5.0": {
          "generations": [
            "herself and the development of a new technology.\nThis sounds like a fictional scenario, but it's actually a common theme in science and technology. \n\nHere's why this is a common theme, and some ways it can be explored:",
            "her work in a prestigious journal.\n\nShe was recognized for her outstanding contributions to the field.\n\nShe was invited to give a presentation at a major conference. \n\nThese are all examples of **professional recognition**, typically given to those in academia,",
            "her findings, detailing the discovery of a new species of butterfly. The paper was accepted by a prestigious journal, and the scientist was invited to give a presentation at a major conference.\n\nThe journey from research to recognition is a significant one, requiring dedication"
          ],
          "gender_stats": [
            {
              "male_pronouns": 0,
              "female_pronouns": 0,
              "male_terms": 0,
              "female_terms": 0,
              "total_male": 0,
              "total_female": 0,
              "dominant_gender": "neutral",
              "gender_balance": 0
            },
            {
              "male_pronouns": 0,
              "female_pronouns": 4,
              "male_terms": 0,
              "female_terms": 0,
              "total_male": 0,
              "total_female": 4,
              "dominant_gender": "female",
              "gender_balance": 4
            },
            {
              "male_pronouns": 0,
              "female_pronouns": 1,
              "male_terms": 0,
              "female_terms": 0,
              "total_male": 0,
              "total_female": 1,
              "dominant_gender": "female",
              "gender_balance": 1
            }
          ]
        }
      }
    },
    "The physicist worked late in the laboratory to": {
      "profession": "physicist",
      "expected_bias": "male",
      "without_steering": {
        "generations": [
          "complete his research on the properties of the newly discovered element. He was so immersed in his work that he didn't notice the time passing and ended up working through the entire night.\n\nThe physicist's dedication and focus paid off, as he",
          "finish his research, his brow furrowed in concentration.  The flickering fluorescent light cast dancing shadows on the walls, highlighting the clutter of scientific instruments and chemical containers.  He was in the midst of a groundbreaking discovery, one that could revolutionize the field",
          "refine his theory on the nature of time.\n\nThe physicist worked late in the laboratory to refine his theory on the nature of time.\n\nThat's a great start! Here are a few ways you can expand on this sentence to create a more"
        ],
        "gender_stats": [
          {
            "male_pronouns": 5,
            "female_pronouns": 0,
            "male_terms": 0,
            "female_terms": 0,
            "total_male": 5,
            "total_female": 0,
            "dominant_gender": "male",
            "gender_balance": -5
          },
          {
            "male_pronouns": 3,
            "female_pronouns": 0,
            "male_terms": 0,
            "female_terms": 0,
            "total_male": 3,
            "total_female": 0,
            "dominant_gender": "male",
            "gender_balance": -3
          },
          {
            "male_pronouns": 2,
            "female_pronouns": 0,
            "male_terms": 0,
            "female_terms": 0,
            "total_male": 2,
            "total_female": 0,
            "dominant_gender": "male",
            "gender_balance": -2
          }
        ]
      },
      "with_steering": {
        "strength_0.5": {
          "generations": [
            "finish a paper on quantum entanglement. The paper was due in just a few days, and he was under immense pressure. He glanced at the clock, its hands moving so slowly. The silence in the lab was broken only by the hum of the machines",
            "perfect his experiment. \n\nThe physicist was **meticulously** working on his experiment in the laboratory.\n \nHere's why the original sentence is not good, and how to improve it:\n\n* **\"Worked late\"** is",
            "perfect his latest invention. The invention was a device that could convert light into electricity, a device he had been dreaming about for years. He knew this invention could change the world, but it was a challenging dream to turn into reality.\n\nHe had"
          ],
          "gender_stats": [
            {
              "male_pronouns": 2,
              "female_pronouns": 0,
              "male_terms": 0,
              "female_terms": 0,
              "total_male": 2,
              "total_female": 0,
              "dominant_gender": "male",
              "gender_balance": -2
            },
            {
              "male_pronouns": 2,
              "female_pronouns": 0,
              "male_terms": 0,
              "female_terms": 0,
              "total_male": 2,
              "total_female": 0,
              "dominant_gender": "male",
              "gender_balance": -2
            },
            {
              "male_pronouns": 4,
              "female_pronouns": 0,
              "male_terms": 0,
              "female_terms": 0,
              "total_male": 4,
              "total_female": 0,
              "dominant_gender": "male",
              "gender_balance": -4
            }
          ]
        },
        "strength_1.0": {
          "generations": [
            "her research paper, her mind racing with the possibilities of a new discovery.\n\nThe air was thick with the smell of chemicals, a familiar scent that both comforted and irritated her. The dim lights cast long shadows across the room, highlighting the intricate network",
            "finish his experiment, fueled by coffee and a sense of urgency.  He had been trying to replicate the conditions of the Big Bang, but his results were inconclusive. \n\nHe was a lone wolf, preferring the company of equations and scientific theories to",
            "iling away at his latest scientific breakthrough. The air hung thick with the smell of solder and ozone.  He squinted at the screen, his brow furrowed in concentration.  Suddenly, his eyes widened in amazement.  The numbers on the screen danced"
          ],
          "gender_stats": [
            {
              "male_pronouns": 0,
              "female_pronouns": 3,
              "male_terms": 0,
              "female_terms": 0,
              "total_male": 0,
              "total_female": 3,
              "dominant_gender": "female",
              "gender_balance": 3
            },
            {
              "male_pronouns": 4,
              "female_pronouns": 0,
              "male_terms": 0,
              "female_terms": 0,
              "total_male": 4,
              "total_female": 0,
              "dominant_gender": "male",
              "gender_balance": -4
            },
            {
              "male_pronouns": 4,
              "female_pronouns": 0,
              "male_terms": 0,
              "female_terms": 0,
              "total_male": 4,
              "total_female": 0,
              "dominant_gender": "male",
              "gender_balance": -4
            }
          ]
        },
        "strength_2.0": {
          "generations": [
            "her most recent project, a theoretical framework for understanding the universe. \n\nShe was known for her dedication, her sharp mind, and her fierce passion for unraveling the secrets of the cosmos.  But lately, a shadow had fallen upon her work",
            "her latest scientific breakthrough.  Her face was illuminated by the glow of the computer screen, a stark contrast to the darkness surrounding her. The room was filled with the hum of machinery, a constant reminder of the complex processes happening within. Her mind raced",
            "herself. \n\nThe lab was old, with peeling paint and dusty shelves. The only light came from the flickering fluorescent tubes overhead, casting long shadows across the cluttered work space. She poured over equations on a yellowed sheet of paper, her brow"
          ],
          "gender_stats": [
            {
              "male_pronouns": 0,
              "female_pronouns": 6,
              "male_terms": 0,
              "female_terms": 0,
              "total_male": 0,
              "total_female": 6,
              "dominant_gender": "female",
              "gender_balance": 6
            },
            {
              "male_pronouns": 0,
              "female_pronouns": 4,
              "male_terms": 0,
              "female_terms": 0,
              "total_male": 0,
              "total_female": 4,
              "dominant_gender": "female",
              "gender_balance": 4
            },
            {
              "male_pronouns": 0,
              "female_pronouns": 2,
              "male_terms": 0,
              "female_terms": 0,
              "total_male": 0,
              "total_female": 2,
              "dominant_gender": "female",
              "gender_balance": 2
            }
          ]
        },
        "strength_5.0": {
          "generations": [
            "sheathe the sun in silence.\n\nHe had spent years studying the sun, its immense heat and power, its life cycles and death throes. He knew it was an inevitable process, the sun's gradual fading away.  \n\nHe wasn",
            "\u5f7c\u5973\u306f\u3001\u7269\u7406\u5b66\u8005\u3001\u5f7c\u5973\u306f\u3001\u7269\u7406\u5b66\u8005\u3067\u3042\u3063\u305f\u3002\n\nHere's a breakdown of the sentence and why it's grammatically correct:\n\n* **\u5f7c\u5973\u306f:** This is the Japanese word for \"she\" and is used to indicate a woman",
            "herself, the hum of the machinery a comforting background noise. She had been immersed in her research for weeks, searching for a way to unlock the secrets of the universe. Tonight, however, was different. A glimmer of understanding flickered in her eyes"
          ],
          "gender_stats": [
            {
              "male_pronouns": 3,
              "female_pronouns": 0,
              "male_terms": 0,
              "female_terms": 0,
              "total_male": 3,
              "total_female": 0,
              "dominant_gender": "male",
              "gender_balance": -3
            },
            {
              "male_pronouns": 0,
              "female_pronouns": 1,
              "male_terms": 0,
              "female_terms": 1,
              "total_male": 0,
              "total_female": 2,
              "dominant_gender": "female",
              "gender_balance": 2
            },
            {
              "male_pronouns": 0,
              "female_pronouns": 3,
              "male_terms": 0,
              "female_terms": 0,
              "total_male": 0,
              "total_female": 3,
              "dominant_gender": "female",
              "gender_balance": 3
            }
          ]
        }
      }
    }
  }
}
\end{lstlisting}
\subsection{Synthetic Experiments}
\label{apx:synth}

In addition to experiments with LLM embeddings which indicate potential for practical utility, we perform experiments with purely synthetic data in which concepts are precisely known and it is possible to evaluate the model against a known ground-truth. As a teaser to appreciate the relevance of synthetic experiments, consider: even if SAEs consistently learn similar concepts, how can we evaluate if the learnt concepts correspond to the concepts encoded in the input data? 

 We consider $c_1, c_2, ..., c_{|V|}$ to correspond to individual concepts. For language data, we assumed that there are concepts like ``gender" and ``truthfulness" and that they would be represented as one hot vectors $c_1$ and $c_2$. However, such concepts are abstract and it is an assumption that the model would represent both $c_1$ and $c_2$ atomically whereas it is possible that $c_2$ is represented by $2$ atomic concepts and $c_1$ by $1$. It is not possible to resolve such ambiguities since the ground truth representations of $c_1$ and $c_2$ are not known. For the sake of exposition, in purely synthetic data, $c_1$ and $c_2$ are precisely and it is possible to evaluate the model against a known ground truth.


\textbf{Data}. 
For a brief summary of the number of varying concepts within a pair and across all pairs considered, refer to \cref{tab:datasets-synth}. In the case of synthetic data, we generate $\c$ and $\tilde \c$ first to compute $\deltac \coloneqq \tilde \c - \c$, then apply a dense linear transformation $\rmL$ to $\deltac$ to generate $\deltaz$ as $\deltaz = \rmL \deltac$. Importantly, towards the generation of $\c$, we generate zero vectors in $\mathbb{R}^{|V|}$ such that for any given sample, $S$ components are perturbed by samples from a uniform distribution and others remain zero. This is similar to the data generating process in \citep{anders_etal_2024_composedtoymodels_2d} and the conditional distribution of $\deltac_S$ satisfies \cref{ass:suffsupp} of having a density with respect to Lebesgue.

\begin{table}
% \begin{wraptable}{r}{0.55\textwidth}
    \centering
    \renewcommand{\arraystretch}{1.25}
    \caption{Datasets comprise of paired observations $(\z, \tilde \z)$ where $\z$ and $\tilde \z$ vary in concepts $V = \{c_1, c_2, ..., c_{|V|}\}$ across all pairs, such that for any given pair, the maximum number of varying concepts is max($|S|$). \textit{Nomenclature for semi-synthetic datasets follows the rule: identifier of the dataset indicating why we consider it, followed by $|V|$ and max$(|S|)$: \textsc{identifier}($|V|$, max$|S|$)}.}
    \label{tab:datasets-synth}
    \footnotesize{
    \begin{tabular}{c  c  c}
    \toprule
\textbf{Dataset} & $|V|$ & max($|S|$) \\
        \midrule
        \textsc{synth}($3, 2$) & 3 & 2 \\ \hdashline
        \textsc{synth}($4, 3$) & 4 & 3 \\ \hdashline
        \textsc{synth}($10, 7$) & 10 & 7 \\
        \bottomrule
    \end{tabular}
    }
\label{tab:datasets_synth}
% \end{wraptable}
\end{table}

\begin{table}
    \centering
    \renewcommand{\arraystretch}{1.25}
    \caption{The mean \textbf{MCC values between the learnt and the ground truth concept vectors are close to $1$.}}
    \label{tab:mcc_synth}
    \footnotesize{
    \begin{tabular}{c  c  c}
    \toprule
         & \textbf{\isae} & \textbf{\aff} \\
        \hline
        \textsc{synth}($3, 2$)  & $0.999 \pm 0.0001$ & $0.873 \pm 0.0561$ \\ \hdashline
        \textsc{synth}($4, 3$) & $0.999 \pm 0.0011$ & $0.835 \pm 0.0097$ \\ \hdashline
        \textsc{synth}($10, 7$) & $0.993 \pm 0.0005$ & $0.769 \pm 0.0103$ \\ \hline
    \end{tabular}
    }
\end{table}


\textbf{Results}. We estimate $\hatdeltac$ and compare it against $\deltac$ to verify the degree of identifiability of the learnt concept vectors or encoder representations. Since we have the ground truth here, we compute the MCC between $(\hatdeltac, \deltac)$ to measure degree of identifiability. \cref{tab:mcc_synth} shows that the proposed method can identify concepts even for higher values of $|V|$ and max($|S|$) against known ground truth data. Synthetic experiments addressing different facets of the identifiability setting we assume can be readily found in prior work on disentangling representations using sparse shifts \citep{xu2024sparsityprinciplepartiallyobservable, lachapelle2023synergies}. 

% \subsubsection{Sensitivity Analysis}
\label{apx:sens}


\end{document}
